% Options for packages loaded elsewhere
% Options for packages loaded elsewhere
\PassOptionsToPackage{unicode}{hyperref}
\PassOptionsToPackage{hyphens}{url}
\PassOptionsToPackage{dvipsnames,svgnames,x11names}{xcolor}
%
\documentclass[
]{article}
\usepackage{xcolor}
\usepackage[margin=1in]{geometry}
\usepackage{amsmath,amssymb}
\setcounter{secnumdepth}{5}
\usepackage{iftex}
\ifPDFTeX
  \usepackage[T1]{fontenc}
  \usepackage[utf8]{inputenc}
  \usepackage{textcomp} % provide euro and other symbols
\else % if luatex or xetex
  \usepackage{unicode-math} % this also loads fontspec
  \defaultfontfeatures{Scale=MatchLowercase}
  \defaultfontfeatures[\rmfamily]{Ligatures=TeX,Scale=1}
\fi
\usepackage{lmodern}
\ifPDFTeX\else
  % xetex/luatex font selection
\fi
% Use upquote if available, for straight quotes in verbatim environments
\IfFileExists{upquote.sty}{\usepackage{upquote}}{}
\IfFileExists{microtype.sty}{% use microtype if available
  \usepackage[]{microtype}
  \UseMicrotypeSet[protrusion]{basicmath} % disable protrusion for tt fonts
}{}
\makeatletter
\@ifundefined{KOMAClassName}{% if non-KOMA class
  \IfFileExists{parskip.sty}{%
    \usepackage{parskip}
  }{% else
    \setlength{\parindent}{0pt}
    \setlength{\parskip}{6pt plus 2pt minus 1pt}}
}{% if KOMA class
  \KOMAoptions{parskip=half}}
\makeatother
% Make \paragraph and \subparagraph free-standing
\makeatletter
\ifx\paragraph\undefined\else
  \let\oldparagraph\paragraph
  \renewcommand{\paragraph}{
    \@ifstar
      \xxxParagraphStar
      \xxxParagraphNoStar
  }
  \newcommand{\xxxParagraphStar}[1]{\oldparagraph*{#1}\mbox{}}
  \newcommand{\xxxParagraphNoStar}[1]{\oldparagraph{#1}\mbox{}}
\fi
\ifx\subparagraph\undefined\else
  \let\oldsubparagraph\subparagraph
  \renewcommand{\subparagraph}{
    \@ifstar
      \xxxSubParagraphStar
      \xxxSubParagraphNoStar
  }
  \newcommand{\xxxSubParagraphStar}[1]{\oldsubparagraph*{#1}\mbox{}}
  \newcommand{\xxxSubParagraphNoStar}[1]{\oldsubparagraph{#1}\mbox{}}
\fi
\makeatother

\usepackage{color}
\usepackage{fancyvrb}
\newcommand{\VerbBar}{|}
\newcommand{\VERB}{\Verb[commandchars=\\\{\}]}
\DefineVerbatimEnvironment{Highlighting}{Verbatim}{commandchars=\\\{\}}
% Add ',fontsize=\small' for more characters per line
\usepackage{framed}
\definecolor{shadecolor}{RGB}{241,243,245}
\newenvironment{Shaded}{\begin{snugshade}}{\end{snugshade}}
\newcommand{\AlertTok}[1]{\textcolor[rgb]{0.68,0.00,0.00}{#1}}
\newcommand{\AnnotationTok}[1]{\textcolor[rgb]{0.37,0.37,0.37}{#1}}
\newcommand{\AttributeTok}[1]{\textcolor[rgb]{0.40,0.45,0.13}{#1}}
\newcommand{\BaseNTok}[1]{\textcolor[rgb]{0.68,0.00,0.00}{#1}}
\newcommand{\BuiltInTok}[1]{\textcolor[rgb]{0.00,0.23,0.31}{#1}}
\newcommand{\CharTok}[1]{\textcolor[rgb]{0.13,0.47,0.30}{#1}}
\newcommand{\CommentTok}[1]{\textcolor[rgb]{0.37,0.37,0.37}{#1}}
\newcommand{\CommentVarTok}[1]{\textcolor[rgb]{0.37,0.37,0.37}{\textit{#1}}}
\newcommand{\ConstantTok}[1]{\textcolor[rgb]{0.56,0.35,0.01}{#1}}
\newcommand{\ControlFlowTok}[1]{\textcolor[rgb]{0.00,0.23,0.31}{\textbf{#1}}}
\newcommand{\DataTypeTok}[1]{\textcolor[rgb]{0.68,0.00,0.00}{#1}}
\newcommand{\DecValTok}[1]{\textcolor[rgb]{0.68,0.00,0.00}{#1}}
\newcommand{\DocumentationTok}[1]{\textcolor[rgb]{0.37,0.37,0.37}{\textit{#1}}}
\newcommand{\ErrorTok}[1]{\textcolor[rgb]{0.68,0.00,0.00}{#1}}
\newcommand{\ExtensionTok}[1]{\textcolor[rgb]{0.00,0.23,0.31}{#1}}
\newcommand{\FloatTok}[1]{\textcolor[rgb]{0.68,0.00,0.00}{#1}}
\newcommand{\FunctionTok}[1]{\textcolor[rgb]{0.28,0.35,0.67}{#1}}
\newcommand{\ImportTok}[1]{\textcolor[rgb]{0.00,0.46,0.62}{#1}}
\newcommand{\InformationTok}[1]{\textcolor[rgb]{0.37,0.37,0.37}{#1}}
\newcommand{\KeywordTok}[1]{\textcolor[rgb]{0.00,0.23,0.31}{\textbf{#1}}}
\newcommand{\NormalTok}[1]{\textcolor[rgb]{0.00,0.23,0.31}{#1}}
\newcommand{\OperatorTok}[1]{\textcolor[rgb]{0.37,0.37,0.37}{#1}}
\newcommand{\OtherTok}[1]{\textcolor[rgb]{0.00,0.23,0.31}{#1}}
\newcommand{\PreprocessorTok}[1]{\textcolor[rgb]{0.68,0.00,0.00}{#1}}
\newcommand{\RegionMarkerTok}[1]{\textcolor[rgb]{0.00,0.23,0.31}{#1}}
\newcommand{\SpecialCharTok}[1]{\textcolor[rgb]{0.37,0.37,0.37}{#1}}
\newcommand{\SpecialStringTok}[1]{\textcolor[rgb]{0.13,0.47,0.30}{#1}}
\newcommand{\StringTok}[1]{\textcolor[rgb]{0.13,0.47,0.30}{#1}}
\newcommand{\VariableTok}[1]{\textcolor[rgb]{0.07,0.07,0.07}{#1}}
\newcommand{\VerbatimStringTok}[1]{\textcolor[rgb]{0.13,0.47,0.30}{#1}}
\newcommand{\WarningTok}[1]{\textcolor[rgb]{0.37,0.37,0.37}{\textit{#1}}}

\usepackage{longtable,booktabs,array}
\newcounter{none} % for unnumbered tables
\usepackage{calc} % for calculating minipage widths
% Correct order of tables after \paragraph or \subparagraph
\usepackage{etoolbox}
\makeatletter
\patchcmd\longtable{\par}{\if@noskipsec\mbox{}\fi\par}{}{}
\makeatother
% Allow footnotes in longtable head/foot
\IfFileExists{footnotehyper.sty}{\usepackage{footnotehyper}}{\usepackage{footnote}}
\makesavenoteenv{longtable}
\usepackage{graphicx}
\makeatletter
\newsavebox\pandoc@box
\newcommand*\pandocbounded[1]{% scales image to fit in text height/width
  \sbox\pandoc@box{#1}%
  \Gscale@div\@tempa{\textheight}{\dimexpr\ht\pandoc@box+\dp\pandoc@box\relax}%
  \Gscale@div\@tempb{\linewidth}{\wd\pandoc@box}%
  \ifdim\@tempb\p@<\@tempa\p@\let\@tempa\@tempb\fi% select the smaller of both
  \ifdim\@tempa\p@<\p@\scalebox{\@tempa}{\usebox\pandoc@box}%
  \else\usebox{\pandoc@box}%
  \fi%
}
% Set default figure placement to htbp
\def\fps@figure{htbp}
\makeatother





\setlength{\emergencystretch}{3em} % prevent overfull lines

\providecommand{\tightlist}{%
  \setlength{\itemsep}{0pt}\setlength{\parskip}{0pt}}



 


\usepackage{fvextra}
\fvset{breaklines,breaksymbol=\tiny\ensuremath{\hookrightarrow},fontsize=\small}
\DefineVerbatimEnvironment{Highlighting}{Verbatim}{breaklines,breaksymbol=\tiny\ensuremath{\hookrightarrow},breaksymbolleft=,fontsize=\small,commandchars=\\\{\}}
\makeatletter
\@ifpackageloaded{caption}{}{\usepackage{caption}}
\AtBeginDocument{%
\ifdefined\contentsname
  \renewcommand*\contentsname{Table of contents}
\else
  \newcommand\contentsname{Table of contents}
\fi
\ifdefined\listfigurename
  \renewcommand*\listfigurename{List of Figures}
\else
  \newcommand\listfigurename{List of Figures}
\fi
\ifdefined\listtablename
  \renewcommand*\listtablename{List of Tables}
\else
  \newcommand\listtablename{List of Tables}
\fi
\ifdefined\figurename
  \renewcommand*\figurename{Figure}
\else
  \newcommand\figurename{Figure}
\fi
\ifdefined\tablename
  \renewcommand*\tablename{Table}
\else
  \newcommand\tablename{Table}
\fi
}
\@ifpackageloaded{float}{}{\usepackage{float}}
\floatstyle{ruled}
\@ifundefined{c@chapter}{\newfloat{codelisting}{h}{lop}}{\newfloat{codelisting}{h}{lop}[chapter]}
\floatname{codelisting}{Listing}
\newcommand*\listoflistings{\listof{codelisting}{List of Listings}}
\makeatother
\makeatletter
\makeatother
\makeatletter
\@ifpackageloaded{caption}{}{\usepackage{caption}}
\@ifpackageloaded{subcaption}{}{\usepackage{subcaption}}
\makeatother
\makeatletter
\definecolor{QuartoInternalColor1}{rgb}{0,0,0}
\definecolor{QuartoInternalColor2}{rgb}{0.63,0.19,0.59}
\definecolor{QuartoInternalColor3}{rgb}{0.70,0.49,0.07}
\makeatother
\usepackage{bookmark}
\IfFileExists{xurl.sty}{\usepackage{xurl}}{} % add URL line breaks if available
\urlstyle{same}
\hypersetup{
  pdftitle={Unit 5: More language features for software architecture},
  colorlinks=true,
  linkcolor={blue},
  filecolor={Maroon},
  citecolor={Blue},
  urlcolor={Blue},
  pdfcreator={LaTeX via pandoc}}


\title{Unit 5: More language features for software architecture}
\author{}
\date{}
\begin{document}
\maketitle

\renewcommand*\contentsname{Table of contents}
{
\setcounter{tocdepth}{3}
\tableofcontents
}

\section{Unit 5: More language features for software
architecture}\label{unit-5-more-language-features-for-software-architecture}

In the previous three units we explored basics of programming and
computation (Unit 2), algorithms and data structures (Unit 3), and data
files and numerics (Unit 4). In this unit we take a deeper and more
thorough approach at basic Julia language features.

Programming languages, including Julia, are designed mostly with these
aims in mind:

\begin{enumerate}
\def\labelenumi{\arabic{enumi}.}
\tightlist
\item
  \textbf{Execution speed} --- programs should run fast enough to solve
  the problem at hand.
\item
  \textbf{Coding speed} --- writing programs to solve problems should be
  quick, easy and enjoyable for the coder.
\item
  \textbf{Scalable engineering} --- creating and maintaining large
  programs with many contributors should be no harder than writing small
  programs yourself.
\end{enumerate}

Different programming languages have different goals and target
audiences, and make different trade-offs and affordances for the above.
The trickiest one to get right is the third one --- as codebases grow it
is \emph{typical} to get stuck in a tarpit of complexity, where it gets
harder and harder to change and improve the code, and progress grinds to
a halt. A corollary of that is that code in large pieces of software is
\emph{read} far more often than it is \emph{written}, so it becomes
crucial to make code that is written easy to understand.

Julia is a high-performance language for solving large mathematical
problems, and compiled code needs to run ``close to the metal''. It
makes many affordances to solve mathematical problems quickly and
efficiently. But most crucially, it tries to solve these in such a way
that is ``scalable'' --- where generic code can be reused and
repurposed, and pieces of a program can be connected together elegantly
like lego bricks.

The goal of this unit is to show you the most important language
features of Julia and how to use them to best effect. The most important
language features that we explore are the syntax, type system,
user-defined types, and multiple-dispatch. At the end we'll consider how
to use these to construct your code with an emphasis on types and
functions (data and transformations / nouns and verbs). Choosing good
\emph{nouns} and \emph{verbs} for your problem domain is the core of
good program design. The notes below often refer to the Julia
documentation.

\subsection{Values and their types}\label{values-and-their-types}

One way Julia allows for easy and fast mathematical code is to emphasise
\emph{values} (or \emph{objects}), and transformations of values (via
functions).

A few ways that is apparent is:

\begin{enumerate}
\def\labelenumi{\arabic{enumi}.}
\tightlist
\item
  All valid Julia code is an \emph{expression} that evaluates to a value
  (no statements, no \texttt{void}).
\item
  Every value has a specific (or ``concrete'') \emph{type} (like
  \texttt{String} or \texttt{Int64}).
\item
  Functions are values, which can be stored in variables or passed to
  other functions.
\end{enumerate}

Many modern languages feature the above, but Julia takes it even further
than that:

\begin{enumerate}
\def\labelenumi{\arabic{enumi}.}
\setcounter{enumi}{3}
\tightlist
\item
  Types are also values that can be passed around and manipulated.
\item
  Julia code is a value (you can use and create macros, access generated
  native code, etc).
\end{enumerate}

With these last two, you can do some pretty powerful things with Julia.
In fact, most of Julia is itself written in Julia. However, we won't
focus on those here.

\subsubsection{Why is it so?}\label{why-is-it-so}

Remember, all this is designed to make it easy for mathematicians to
write code that fast to run and easy to scale.

In mathematics, we are interested in values or objects, as well as
operations and transformations on such objects. Basing the language
around values (objects), and allowing functions to behave like values
(objects), makes our code resemble mathematics and have a similar
interpretation. It does this without ``getting in your way'' or stopping
you writing imperative code.

From a historical perspective, Julia is strongly influenced by lisp
(scheme) which is a functional language where everything is an
expression.

The fact that values have \emph{types} is for two reasons:

\begin{enumerate}
\def\labelenumi{\arabic{enumi}.}
\tightlist
\item
  Types let us specify our data and define the \emph{behaviour} of
  functions that operator on them (using multiple dispatch).
\item
  Types let the compiler know the \emph{representation} of the value in
  memory (e.g.~how many bytes it is) so it can generate fast code.
\end{enumerate}

Types have two purposes: help \emph{humans} organise their data and
define behaviour, and the help \emph{the computer} to execute the
program quickly.

\subsubsection{\texorpdfstring{All valid code is an
\emph{expression}}{All valid code is an expression}}\label{all-valid-code-is-an-expression}

An expression is a piece of code that evaluates to a value. (Some
languages, like Python or C, also have other pieces of syntax which are
not expressions, typically known as ``statements'' - a statement tells
to computer or compiler to ``do something'' but does not produce a value
you can use).

Here are some expressions in Julia:

\begin{Shaded}
\begin{Highlighting}[]
\FloatTok{1} \OperatorTok{+} \FloatTok{1}
\NormalTok{[}\StringTok{"a"}\NormalTok{, }\StringTok{"b"}\NormalTok{, }\StringTok{"c"}\NormalTok{]}
\FunctionTok{println}\NormalTok{(}\ConstantTok{pi}\NormalTok{)}
\end{Highlighting}
\end{Shaded}

\begin{verbatim}
π
\end{verbatim}

One nice aspects of expressions is you can copy them and paste them
anywhere that expects an expression, and it is guaranteed to by
syntactically valid. Each of these expressions evaluate to a value,
which you can for example bind to a variable:

\begin{Shaded}
\begin{Highlighting}[]
\NormalTok{x }\OperatorTok{=} \FloatTok{1} \OperatorTok{+} \FloatTok{1}
\NormalTok{y }\OperatorTok{=}\NormalTok{ [}\StringTok{"a"}\NormalTok{, }\StringTok{"b"}\NormalTok{, }\StringTok{"c"}\NormalTok{]}
\NormalTok{z }\OperatorTok{=} \FunctionTok{println}\NormalTok{(}\ConstantTok{pi}\NormalTok{)}
\end{Highlighting}
\end{Shaded}

\begin{verbatim}
π
\end{verbatim}

\subsubsection{Every value has a concrete
type}\label{every-value-has-a-concrete-type}

The types of the first value is straightforward:

\begin{Shaded}
\begin{Highlighting}[]
\FunctionTok{typeof}\NormalTok{(}\FloatTok{1}\OperatorTok{+}\FloatTok{1}\NormalTok{)}
\end{Highlighting}
\end{Shaded}

\begin{verbatim}
Int64
\end{verbatim}

This is just a 64-bit signed integer, which requires 64 bits (or 8
bytes) of memory to represent at runtime.

The second is a 1D array, or vector:

\begin{verbatim}
typeof(["a", "b", "c"])
\end{verbatim}

Vectors can be arbitrarily long and use an unknown amount of memory at
runtime. We'll talk about type parameters (the \texttt{\{String\}} part)
and aliases in more detail later, but first let's look at the third:

\begin{Shaded}
\begin{Highlighting}[]
\FunctionTok{typeof}\NormalTok{(}\FunctionTok{println}\NormalTok{(}\ConstantTok{pi}\NormalTok{))}
\end{Highlighting}
\end{Shaded}

\begin{verbatim}
π
\end{verbatim}

\begin{verbatim}
Nothing
\end{verbatim}

The \texttt{println} function takes input and prints it to the console
(followed by a new line). It doesn't really have any ``value'' to
return. But I said earlier that all expressions in Julia return
\emph{some} value. Julia has an value defined called \texttt{nothing}
for this purpose, which is of type \texttt{Nothing}. There is nothing
special about this type, except that it has no associated data (meaning
it needs 0 bytes of data to represent at runtime).

\subsubsection{What is a concrete type?}\label{what-is-a-concrete-type}

A type can be thought of as a set of values.

Every value is a concrete type, and Julia has just these concrete types:

\begin{enumerate}
\def\labelenumi{\arabic{enumi}.}
\tightlist
\item
  \texttt{primitive\ type}
\item
  \texttt{struct} and \texttt{mutable\ struct}
\item
  \texttt{Memory\{T\}} (a 1D non-resizable container of \texttt{T}, used
  to store the elements of \texttt{Array\{N,\ T\}})
\item
  \texttt{String}
\item
  \texttt{Tuple}
\end{enumerate}

\subsubsection{Primitive types}\label{primitive-types}

Primitive types is where we begin describing data at the bits-and-bytes
level. Here is the definition of \texttt{UInt8}:

\begin{Shaded}
\begin{Highlighting}[]
\KeywordTok{primitive type} \DataTypeTok{UInt8}
   \FloatTok{8}
\KeywordTok{end}
\end{Highlighting}
\end{Shaded}

There are 256 possible \texttt{UInt8} values:
\texttt{0x00,\ 0x01,\ ...,\ 0xfe,\ 0xff}. You could say \texttt{UInt8}
is a set with 256 elements.

Here are the definitions of \texttt{Int64} and \texttt{Float64}:

\begin{Shaded}
\begin{Highlighting}[]
\KeywordTok{primitive type} \DataTypeTok{Int64}
   \FloatTok{64}
\KeywordTok{end}

\KeywordTok{primitive type} \DataTypeTok{Float64}
   \FloatTok{64}
\KeywordTok{end}
\end{Highlighting}
\end{Shaded}

Each of these are sets containing \texttt{2\^{}64} (about
\texttt{10\^{}18}) possible values. However, for every value in your
program Julia keeps track of the type of the value, as well as the data.
You can consider \texttt{Int64} and \texttt{Float64} to be
\emph{disjoint} sets. The values \texttt{0} (an integer) and
\texttt{0.0} (a float) are represented by exactly the same bits (all
zero), yet they are not ``the same'':

\begin{Shaded}
\begin{Highlighting}[]
\FunctionTok{reinterpret}\NormalTok{(}\DataTypeTok{Float64}\NormalTok{, }\FloatTok{0}\NormalTok{)}
\end{Highlighting}
\end{Shaded}

\begin{verbatim}
0.0
\end{verbatim}

\begin{Shaded}
\begin{Highlighting}[]
\FunctionTok{reinterpret}\NormalTok{(}\DataTypeTok{Int64}\NormalTok{, }\FloatTok{0.0}\NormalTok{)}
\end{Highlighting}
\end{Shaded}

\begin{verbatim}
0
\end{verbatim}

\begin{Shaded}
\begin{Highlighting}[]
\FloatTok{0} \OperatorTok{===} \FloatTok{0.0}
\end{Highlighting}
\end{Shaded}

\begin{verbatim}
false
\end{verbatim}

(Type tracking typically occurs at compile time, before the program
runs, so Julia doesn't need to store extra data describing whether every
value is an integer or a float or something else.)

The advantage of distinct types is we can describe different behaviour
for these types. For example, there are separate methods for addition
representing the different ways the CPU handles integer and
floating-point arithmetic, \texttt{+(x::Int64,\ y::Int64)} and
\texttt{+(x::Float64,\ y::Float64)}.

(Note that all primitive types are immutable values. You cannot
``mutate'' \texttt{0} to become \texttt{1} - you construct a new value
instead.)

\subsubsection{Structs}\label{structs}

Structs allow us to organise different pieces of data into a
``structure''. In general, structs represent composition of data: e.g.~a
complex number consists of a real component and an imaginary component:

\begin{Shaded}
\begin{Highlighting}[]
\KeywordTok{struct} \DataTypeTok{Complex64}
\NormalTok{   re}\OperatorTok{::}\DataTypeTok{Float64}
   \ConstantTok{im}\OperatorTok{::}\DataTypeTok{Float64}
\KeywordTok{end}
\end{Highlighting}
\end{Shaded}

The compiler knows that a \texttt{Complex64} needs 128 bits of data to
represent - 64 bits each for \texttt{re} and \texttt{im}. This type is a
distinct set with \texttt{2\^{}128} possible values.

You can construct a struct with the ``constructor'' - where the type
doubles as a kind of function that lets you specify the fields:

\begin{Shaded}
\begin{Highlighting}[]
\NormalTok{z }\OperatorTok{=} \FunctionTok{Complex64}\NormalTok{(}\FloatTok{1.0}\NormalTok{, }\FloatTok{2.0}\NormalTok{)}
\end{Highlighting}
\end{Shaded}

\begin{verbatim}
Complex64(1.0, 2.0)
\end{verbatim}

\begin{Shaded}
\begin{Highlighting}[]
\FunctionTok{sizeof}\NormalTok{(z)}
\end{Highlighting}
\end{Shaded}

\begin{verbatim}
16
\end{verbatim}

You can get a field from the struct with the \texttt{.} operator, for
example

\begin{Shaded}
\begin{Highlighting}[]
\NormalTok{z.re}
\end{Highlighting}
\end{Shaded}

\begin{verbatim}
1.0
\end{verbatim}

\begin{Shaded}
\begin{Highlighting}[]
\NormalTok{z.}\ConstantTok{im}
\end{Highlighting}
\end{Shaded}

\begin{verbatim}
2.0
\end{verbatim}

\paragraph{Empty structs}\label{empty-structs}

A struct can have any number of fields, including zero:

\begin{Shaded}
\begin{Highlighting}[]
\KeywordTok{struct} \DataTypeTok{Nothing}
\KeywordTok{end}

\KeywordTok{const} \ConstantTok{nothing} \OperatorTok{=} \FunctionTok{Nothing}\NormalTok{()}
\end{Highlighting}
\end{Shaded}

This is exactly how \texttt{nothing} is defined in Julia! It takes zero
bytes to represent nothing, and you could say \texttt{Nothing} is like a
set containing \texttt{2\^{}0\ =\ 1} possible value - \texttt{nothing}.

\subparagraph{Mutable structs}\label{mutable-structs}

By default structs are immutable - once constructed you cannot edit the
fields. It is possible to create structs that you can modify with the
\texttt{mutable} keyword:

\begin{Shaded}
\begin{Highlighting}[]
\KeywordTok{mutable struct}\NormalTok{ MutableComplex64}
\NormalTok{   re}\OperatorTok{::}\DataTypeTok{Float64}
   \ConstantTok{im}\OperatorTok{::}\DataTypeTok{Float64}
\KeywordTok{end}

\NormalTok{z }\OperatorTok{=} \FunctionTok{MutableComplex64}\NormalTok{(}\FloatTok{1.0}\NormalTok{, }\FloatTok{2.0}\NormalTok{)}

\NormalTok{z.}\ConstantTok{im} \OperatorTok{=} \FloatTok{0.0}  \CommentTok{\# this would throw an error if z were immutable}
\end{Highlighting}
\end{Shaded}

\paragraph{Generic structs}\label{generic-structs}

Finally, one more thing you can do is create \emph{generic} structs. In
Julia the type for complex numbers is more like:

\begin{Shaded}
\begin{Highlighting}[]
\KeywordTok{struct} \DataTypeTok{Complex}\NormalTok{\{T\}}
\NormalTok{   re}\OperatorTok{::}\DataTypeTok{T}
   \ConstantTok{im}\OperatorTok{::}\DataTypeTok{T}
\KeywordTok{end}
\end{Highlighting}
\end{Shaded}

Given some type \texttt{T}, \texttt{Complex\{T\}} is a struct with
fields \texttt{re} and \texttt{im} of type \texttt{T}.

This useful because now we don't need a separate \texttt{struct}
definition for complex numbers with \texttt{Float32} components, or
\texttt{Int64} components, or whatever. Just one definition suffices for
them all. This lets us write quite generic code that is oblivious to the
exact type \texttt{T} but uses properties common to all complex numbers.
For example, the absolute value squared can be written:

\begin{Shaded}
\begin{Highlighting}[]
\KeywordTok{function} \FunctionTok{abs2}\NormalTok{(z}\OperatorTok{::}\DataTypeTok{Complex}\NormalTok{)}
   \ControlFlowTok{return}\NormalTok{ z.re}\OperatorTok{*}\NormalTok{z.re }\OperatorTok{+}\NormalTok{ z.}\ConstantTok{im}\OperatorTok{*}\NormalTok{z.}\ConstantTok{im}
\KeywordTok{end}
\end{Highlighting}
\end{Shaded}

\begin{verbatim}
abs2 (generic function with 1 method)
\end{verbatim}

and this single definition is sufficient for all possible \texttt{T}.

Note that the type \texttt{Complex} is an \emph{abstract} type (not a
\emph{concrete} type). We'll talk more about abstract types and multiple
dispatch later. For now, note that Julia lets you write both generic
data types and generic functions.

\paragraph{Type aliases}\label{type-aliases}

Earlier we saw an example of a type ``alias''. You can create them like
so:

\begin{Shaded}
\begin{Highlighting}[]
\KeywordTok{const} \DataTypeTok{Complex64} \OperatorTok{=} \DataTypeTok{Complex}\NormalTok{\{}\DataTypeTok{Float64}\NormalTok{\}}
\end{Highlighting}
\end{Shaded}

Type aliases are useful to make your code shorter or more readable.

\paragraph{Tuples}\label{tuples}

Tuples in Julia are a kind of ``lightweight'' struct with no type or
field names that you need to define in advance. You construct tuples
with brackets like \texttt{(1,\ 2,\ 3)}. Note the types don't have to be
uniform, \texttt{t\ =\ (1,\ true,\ "abc")} is equally valid. You access
the fields with the index, such as \texttt{t{[}3{]}\ ===\ "abc"}.

Note the syntax of tuples is much like the arguments to a function -
\texttt{(x,\ y)} is a tuple and \texttt{f(x,\ y)} calls function
\texttt{f} with inputs \texttt{(x,\ y)}. The correspondence is
intentional.

Julia also has ``named'' tuples, where you can name the fields, like
\texttt{(a\ =\ 1,\ b\ =\ true,\ c\ =\ "abc")}. (These look a bit like
keyword arguments to functions).

\subsubsection{Arrays}\label{arrays}

Structs are useful for collecting data, but only if you know how many
elements you want at compile time.

Arrays are dynamically-sized collections, that can contain zero-or-more
elements. We use arrays a lot in Julia.

Arrays in Julia are multi-dimensional and have the type
\texttt{Array\{T,\ N\}}. Each array value knows the type of its contents
\texttt{T} as well as how many dimensions \texttt{N} it has. However,
since we don't know how many elements might be in an array, the computer
needs to allocate memory to store the array at runtime. In terms of
sets, arrays are a type contain an unbounded number of possible values.

Arrays are mutable collections - elements can be edited, inserted or
deleted.

Remember that vectors and matrices are just aliases for 1D and 2D
arrays.

\begin{Shaded}
\begin{Highlighting}[]
\KeywordTok{const} \DataTypeTok{Vector}\NormalTok{\{T\} }\OperatorTok{=} \DataTypeTok{Array}\NormalTok{\{}\FloatTok{1}\NormalTok{, T\}}
\KeywordTok{const} \DataTypeTok{Matrix}\NormalTok{\{T\} }\OperatorTok{=} \DataTypeTok{Array}\NormalTok{\{}\FloatTok{2}\NormalTok{, T\}}
\end{Highlighting}
\end{Shaded}

Underneath, \texttt{Array\{N,T\}} is now (as of Julia 1.11) implemented
in terms of the low level \texttt{Memory\{T\}} - a one dimensional array
of elements of type \texttt{T} with a length which is fixed at the time
of construction. However, \texttt{Vector\{T\}} is more convenient to
work with and you'll generally see that used in most Julia code.

\subsubsection{Strings}\label{strings}

Strings are built into the Julia language, for convenience. They are
much like a \texttt{Vector\{UInt8\}} containing UTF-8 encoded string
data, but are immutable.

(If a future version of Julia had immutable arrays, \texttt{String}
could theoretically become a ``normal'' Julia type, but for now this is
a pragmatic tradeoff).

\textbf{\ldots and, that's it!}

Julia provides other important types, like \texttt{Set}s and
\texttt{Dict}s, but these are implemented in terms of the above.
\emph{Every} value is one of:

\begin{enumerate}
\def\labelenumi{\arabic{enumi}.}
\tightlist
\item
  A primitive type
\item
  A struct (immutable, mutable, tuple or named tuple)
\item
  A \texttt{Memory}
\item
  A \texttt{String}
\end{enumerate}

\subsection{Logic and syntax}\label{logic-and-syntax}

Now that we've covered all the possible types of values in Julia, we'll
take a look at how data is manipulated via logic and functions.

Eventually we'll come back to abstract types and multiple dispatch, but
first we'll cover some fundamentals of Julia code syntax and functions
so we understand \emph{what} we are looking at later. A lot of this
might be intuitive to some of you, but since we are mathematicians, we'd
like to spell it out more formally :)

\subsubsection{Everything is an expression, everything has a value, and
everything has a
type}\label{everything-is-an-expression-everything-has-a-value-and-everything-has-a-type}

We saw earlier that an expression is a piece of code that, when
evaluated, produces a value. For example:

\begin{Shaded}
\begin{Highlighting}[]
\FloatTok{1} \OperatorTok{+} \FloatTok{1}
\end{Highlighting}
\end{Shaded}

\begin{verbatim}
2
\end{verbatim}

One thing you should understand about Julia is that it uses type
inference to predict the output type of an expression.

\begin{Shaded}
\begin{Highlighting}[]
\ImportTok{using} \BuiltInTok{InteractiveUtils}

\PreprocessorTok{@code\_typed} \FloatTok{1} \OperatorTok{+} \FloatTok{1}
\end{Highlighting}
\end{Shaded}

\begin{verbatim}
CodeInfo(
1 ─ %1 = intrinsic Base.add_int(x, y)::Int64
└──      return %1
) => Int64
\end{verbatim}

The compiler essentially contains a ``proof engine'' that can formally
prove constraints on the outputs of expressions. In this case it knows
for sure that the output will be an \texttt{Int64}. Sometimes no proof
can be found - in Julia this just makes the code slower, while in some
other programming languages that could be a compilation error. This
choice makes Julia feel more ``dynamic'' and makes prototype code easier
to write.

\subsubsection{\texorpdfstring{\texttt{if}-\texttt{else}}{if-else}}\label{if-else}

In Julia, \emph{every} piece of syntactically valid code is an
expression --- which means at run time it generates a value which you
can assign to a variable. Even things like \texttt{if} statements!

\begin{Shaded}
\begin{Highlighting}[]
\NormalTok{x }\OperatorTok{=} \FloatTok{42}

\NormalTok{is\_x\_even }\OperatorTok{=} \ControlFlowTok{if}\NormalTok{ x }\OperatorTok{\%} \FloatTok{2} \OperatorTok{==} \FloatTok{0}
   \StringTok{"}\SpecialCharTok{$}\NormalTok{x}\StringTok{ is even"}
\ControlFlowTok{else}
   \StringTok{"}\SpecialCharTok{$}\NormalTok{x}\StringTok{ is odd"}
\ControlFlowTok{end}
\end{Highlighting}
\end{Shaded}

\begin{verbatim}
"42 is even"
\end{verbatim}

Like many languages, Julia has a shorthand ``ternary operator''
(i.e.~operator taking 3 inputs) for \emph{if-else} using \texttt{?} and
\texttt{:}

\begin{Shaded}
\begin{Highlighting}[]
\NormalTok{is\_x\_even }\OperatorTok{=}\NormalTok{ x }\OperatorTok{\%} \FloatTok{2} \OperatorTok{==} \FloatTok{0}\NormalTok{ ? }\StringTok{"}\SpecialCharTok{$}\NormalTok{x}\StringTok{ is even"} \OperatorTok{:} \StringTok{"}\SpecialCharTok{$}\NormalTok{x}\StringTok{ is odd"}
\end{Highlighting}
\end{Shaded}

\begin{verbatim}
"42 is even"
\end{verbatim}

It's worth noting these are 100\% identical. We say that
\texttt{x\ ?\ y\ :\ z} is ``syntax sugar'' for the \texttt{if} statement
--- it's a sweetener to make the programmer's life easier, but not an
entirely new feature since you can always achieve the same thing with
more characters of code. Another piece of syntax sugar is that
\texttt{x\ +\ y} has exactly the same meaning as \texttt{+(x,\ y)}.

\paragraph{Our first abstract type}\label{our-first-abstract-type}

What if the types of the \texttt{true} branch and the \texttt{false}
branch of \texttt{if}-\texttt{else} don't match?

\begin{Shaded}
\begin{Highlighting}[]
\FunctionTok{is\_even}\NormalTok{(x) }\OperatorTok{=}\NormalTok{ x }\OperatorTok{\%} \FloatTok{2} \OperatorTok{==} \FloatTok{0}\NormalTok{ ? }\ConstantTok{true} \OperatorTok{:} \StringTok{"false"}
\end{Highlighting}
\end{Shaded}

\begin{verbatim}
is_even (generic function with 1 method)
\end{verbatim}

\emph{Q: what is the output type of \texttt{is\_even}?}

\begin{Shaded}
\begin{Highlighting}[]
\FunctionTok{is\_even}\NormalTok{(}\FloatTok{2}\NormalTok{)}
\end{Highlighting}
\end{Shaded}

\begin{verbatim}
true
\end{verbatim}

\begin{Shaded}
\begin{Highlighting}[]
\FunctionTok{is\_even}\NormalTok{(}\FloatTok{3}\NormalTok{)}
\end{Highlighting}
\end{Shaded}

\begin{verbatim}
"false"
\end{verbatim}

Sometimes it's a \texttt{Bool} and sometimes it's a \texttt{String}.

\begin{Shaded}
\begin{Highlighting}[]
\PreprocessorTok{@code\_typed} \FunctionTok{is\_even}\NormalTok{(}\FloatTok{42}\NormalTok{)}
\end{Highlighting}
\end{Shaded}

\begin{verbatim}
CodeInfo(
1 ─ %1 = intrinsic Base.checked_srem_int(x, 2)::Int64
│   %2 =   builtin (%1 === 0)::Bool
└──      goto #3 if not %2
2 ─      return true
3 ─      return "false"
) => Union{Bool, String}
\end{verbatim}

If \texttt{Bool} represents the set containing \texttt{true} and
\texttt{false}, and \texttt{String} represents the set of all possible
strings, then the type \texttt{Union\{Bool,\ String\}} represents the
union of these two sets.

\texttt{Union} is the first \textbf{abstract type} we'll cover, and is
relatively straightforward. \texttt{Union} itself is not a concrete
type; there are no values of type \texttt{Union}. It mostly exists in
Julia to aid with type inference for cases like
\texttt{if}-\texttt{else} above -- to abstractly analyse the outputs and
intermediate values of programs during compilation (where more precise
inference leads to more efficient code generation). In cases like this
Julia can keep track which of the two possibilities the data is in
relatively efficiently (e.g.~with a single byte for a tag).

There are two more ways abstract types can show up Julia that we'll see
later.

\paragraph{Short cutting logic}\label{short-cutting-logic}

Julia has operators for Boolean arithmetic, including \texttt{\&} for
\emph{AND} (logical conjunction) and \texttt{\textbar{}} for \emph{OR}
(logical disjunction).

\begin{Shaded}
\begin{Highlighting}[]
\PreprocessorTok{@show} \ConstantTok{false} \OperatorTok{\&} \ConstantTok{false}
\PreprocessorTok{@show} \ConstantTok{false} \OperatorTok{\&} \ConstantTok{true}
\PreprocessorTok{@show} \ConstantTok{true} \OperatorTok{\&} \ConstantTok{false}
\PreprocessorTok{@show} \ConstantTok{true} \OperatorTok{\&} \ConstantTok{true}
\end{Highlighting}
\end{Shaded}

\begin{verbatim}
false & false = false
false & true = false
true & false = false
true & true = true
\end{verbatim}

\begin{verbatim}
true
\end{verbatim}

\begin{Shaded}
\begin{Highlighting}[]
\PreprocessorTok{@show} \ConstantTok{false} \OperatorTok{|} \ConstantTok{false}
\PreprocessorTok{@show} \ConstantTok{false} \OperatorTok{|} \ConstantTok{true}
\PreprocessorTok{@show} \ConstantTok{true} \OperatorTok{|} \ConstantTok{false}
\PreprocessorTok{@show} \ConstantTok{true} \OperatorTok{|} \ConstantTok{true}
\end{Highlighting}
\end{Shaded}

\begin{verbatim}
false | false = false
false | true = true
true | false = true
true | true = true
\end{verbatim}

\begin{verbatim}
true
\end{verbatim}

The operation \texttt{x\ \&\ y} is equivalent to \texttt{\&(x,\ y)} and
will evaluate \texttt{x}, then evaluate \texttt{y} and then determine
the result.

The operation \texttt{x\ \&\&\ y} is similar but it sometimes takes a
shortcut. If \texttt{x} is false, it doesn't bother compute a value for
\texttt{y} (as the answer is \texttt{false} no matter the value for
\texttt{y}).

This matters when \texttt{y} is particularly difficult/slow to
calculate, or when it may be something you want to avoid doing when
\texttt{x} is false.

Note that \texttt{x\ \&\&\ y} is 100\% equivalent to:

\begin{Shaded}
\begin{Highlighting}[]
\ControlFlowTok{if}\NormalTok{ x}
\NormalTok{   y}
\ControlFlowTok{else}
   \ConstantTok{false}
\ControlFlowTok{end}
\end{Highlighting}
\end{Shaded}

Note that the type of \texttt{y} doesn't even need to be \texttt{Bool},
since it is never compared to anything. Sometimes people use this in a
cute way:

\begin{Shaded}
\begin{Highlighting}[]
\NormalTok{x }\OperatorTok{\textless{}} \FloatTok{0} \OperatorTok{\&\&} \FunctionTok{println}\NormalTok{(}\StringTok{"x is negative"}\NormalTok{)}
\end{Highlighting}
\end{Shaded}

People tend to use that because it takes less space to write than:

\begin{Shaded}
\begin{Highlighting}[]
\ControlFlowTok{if}\NormalTok{ x }\OperatorTok{\textless{}} \FloatTok{0}
   \FunctionTok{println}\NormalTok{(}\StringTok{"x is negative"}\NormalTok{)}
\ControlFlowTok{end}
\end{Highlighting}
\end{Shaded}

Don't feel you need to use ``cute'' code like this, but if you see it
somewhere I don't want you to feel confused.

\emph{Q: What does
\texttt{(x\ \textless{}\ 0)\ \&\ println("x\ is\ negative")} do?}

There is an equivalent shortcut operator for \texttt{\textbar{}}:

\begin{Shaded}
\begin{Highlighting}[]
\NormalTok{x }\OperatorTok{||}\NormalTok{ y}
\end{Highlighting}
\end{Shaded}

which is equivalent to

\begin{Shaded}
\begin{Highlighting}[]
\ControlFlowTok{if}\NormalTok{ x}
   \ConstantTok{true}
\ControlFlowTok{else}
\NormalTok{   y}
\ControlFlowTok{end}
\end{Highlighting}
\end{Shaded}

One ``cute'' way to people use this is to perform assertions:

\begin{Shaded}
\begin{Highlighting}[]
\NormalTok{x }\OperatorTok{\textgreater{}=} \FloatTok{0} \OperatorTok{||} \FunctionTok{error}\NormalTok{(}\StringTok{"x is negative"}\NormalTok{)}
\end{Highlighting}
\end{Shaded}

\subsubsection{Blocks of code}\label{blocks-of-code}

Above, we saw blocks of code containing values. They didn't really
\emph{do} anything, like assign values to variables.

The rule in Julia is that the last expression in a block of code is the
value returned from that block of code. If for some reason you want to
\emph{explicitly} create a code block (sometimes useful) you can use
\texttt{begin} like so:

\begin{Shaded}
\begin{Highlighting}[]
\ControlFlowTok{begin}
   \FloatTok{1}
   \FloatTok{2}
   \FloatTok{3}
\ControlFlowTok{end}
\end{Highlighting}
\end{Shaded}

\begin{verbatim}
3
\end{verbatim}

We see blocks of code everywhere in Julia --- inside \texttt{if} and
\texttt{else} statements, within \texttt{for} and \texttt{while} loops,
within \texttt{function} bodies, etc. Blocks of code can be nested ---
we can have \texttt{if} statements inside \texttt{while} loops inside
\texttt{function} bodies, and we tend to represent this visually with
indentation:

\begin{Shaded}
\begin{Highlighting}[]
\KeywordTok{function} \FunctionTok{f}\NormalTok{()}
\NormalTok{   x }\OperatorTok{=} \FloatTok{0}
   \ControlFlowTok{while}\NormalTok{ x }\OperatorTok{\textless{}} \FloatTok{10}
   \ControlFlowTok{if}\NormalTok{ x }\OperatorTok{\%} \FloatTok{2} \OperatorTok{==} \FloatTok{0}
   \FunctionTok{println}\NormalTok{(}\StringTok{"}\SpecialCharTok{$}\NormalTok{x}\StringTok{ is even"}\NormalTok{)}
   \ControlFlowTok{else}
   \FunctionTok{println}\NormalTok{(}\StringTok{"}\SpecialCharTok{$}\NormalTok{x}\StringTok{ is odd"}\NormalTok{)}
   \ControlFlowTok{end}
\NormalTok{   x }\OperatorTok{=}\NormalTok{ x }\OperatorTok{+} \FloatTok{1}
   \ControlFlowTok{end}
   \ControlFlowTok{return}\NormalTok{ x}
\KeywordTok{end}

\FunctionTok{f}\NormalTok{()}
\end{Highlighting}
\end{Shaded}

\begin{verbatim}
0 is even
1 is odd
2 is even
3 is odd
4 is even
5 is odd
6 is even
7 is odd
8 is even
9 is odd
\end{verbatim}

\begin{verbatim}
10
\end{verbatim}

\emph{Note: \texttt{for} and \texttt{while} loops always produce
\texttt{nothing}!}

If you prefer, you can separate multiple expressions in a code block
with \texttt{;} instead of on separate lines. Sometimes you'd even use
parentheses to clarify where the code block begins and ends.

\begin{Shaded}
\begin{Highlighting}[]
\NormalTok{x }\OperatorTok{=}\NormalTok{ (}\FloatTok{1}\NormalTok{; }\FloatTok{2}\NormalTok{; }\FloatTok{3}\NormalTok{)}
\end{Highlighting}
\end{Shaded}

\begin{verbatim}
3
\end{verbatim}

\subsubsection{\texorpdfstring{Empty blocks of code return
\texttt{nothing}}{Empty blocks of code return nothing}}\label{empty-blocks-of-code-return-nothing}

So what happens when the block of code is empty?

\begin{Shaded}
\begin{Highlighting}[]
\ControlFlowTok{begin}
\ControlFlowTok{end}
\end{Highlighting}
\end{Shaded}

Since every expression has a value in Julia, it must return
\emph{something}! That something is called \texttt{nothing}. It is a
special builtin value that contains 0 bytes of data and has the builtin
type \texttt{Nothing}.

\begin{Shaded}
\begin{Highlighting}[]
\FunctionTok{typeof}\NormalTok{(}\ConstantTok{nothing}\NormalTok{)}
\end{Highlighting}
\end{Shaded}

\begin{verbatim}
Nothing
\end{verbatim}

While the above is contrived, there are some more likely places to see
this.

\begin{Shaded}
\begin{Highlighting}[]
\NormalTok{half\_of\_x }\OperatorTok{=} \ControlFlowTok{if}\NormalTok{ x }\OperatorTok{\%} \FloatTok{2} \OperatorTok{==} \FloatTok{0}
\NormalTok{   x }\OperatorTok{÷} \FloatTok{2}
\ControlFlowTok{end}
\end{Highlighting}
\end{Shaded}

This is \emph{implicitly} the same as

\begin{Shaded}
\begin{Highlighting}[]
\NormalTok{half\_of\_x }\OperatorTok{=} \ControlFlowTok{if}\NormalTok{ x }\OperatorTok{\%} \FloatTok{2} \OperatorTok{==} \FloatTok{0}
\NormalTok{   x }\OperatorTok{÷} \FloatTok{2}
\ControlFlowTok{else}
\ControlFlowTok{end}
\end{Highlighting}
\end{Shaded}

which is \emph{implicitly} the same as

\begin{Shaded}
\begin{Highlighting}[]
\NormalTok{half\_of\_x }\OperatorTok{=} \ControlFlowTok{if}\NormalTok{ x }\OperatorTok{\%} \FloatTok{2} \OperatorTok{==} \FloatTok{0}
\NormalTok{   x }\OperatorTok{÷} \FloatTok{2}
\ControlFlowTok{else}
   \ConstantTok{nothing}
\ControlFlowTok{end}
\end{Highlighting}
\end{Shaded}

So if \texttt{x} is odd, then \texttt{half\_of\_x} is \texttt{nothing}.

We also use \texttt{nothing} as the return value for functions that
don't have a value to return --- e.g.~when the purpose of the function
is to perform some action rather that compute some value.

\begin{Shaded}
\begin{Highlighting}[]
\KeywordTok{function} \FunctionTok{hello}\NormalTok{()}
   \FunctionTok{println}\NormalTok{(}\StringTok{"Hello class"}\NormalTok{)}
   \ControlFlowTok{return} \ConstantTok{nothing}
\KeywordTok{end}

\NormalTok{x }\OperatorTok{=} \FunctionTok{hello}\NormalTok{()}
\PreprocessorTok{@show}\NormalTok{ x;}
\end{Highlighting}
\end{Shaded}

\begin{verbatim}
Hello class
x = nothing
\end{verbatim}

\emph{Q: Does anyone know what the equivalent thing in \texttt{C} is?}

\subsubsection{Why expressions?
Metaprogramming!}\label{why-expressions-metaprogramming}

Having the syntax rule that ``all valid syntax is an expression'' makes
it easier to analyse and manipulate code than it would be otherwise. It
lets you rearrange code much like you can manipulate a mathematical
expression with algebra (via substitution, etc), resulting in code that
is correct and still compiles.

For human coders, this means it is easier to copy-paste code from one
place to another, or refactor it into a function.

A ``metaprogram'' is a program that writes a program. Julia supports
metaprogramming through macros and other advanced techniques. We won't
be teaching how to create a macro, but using a macro is pretty easy. A
good example is the \texttt{@show} macro:

\begin{Shaded}
\begin{Highlighting}[]
\NormalTok{x }\OperatorTok{=} \FloatTok{10}
\PreprocessorTok{@show}\NormalTok{ x }\OperatorTok{+} \FloatTok{1}\NormalTok{;}
\end{Highlighting}
\end{Shaded}

\begin{verbatim}
x + 1 = 11
\end{verbatim}

Effectively the macro takes the ``expression'' \texttt{x\ +\ 1}, prints
the expression, prints equals, and prints the evaluated value. While
Julia is busy replacing the builtin syntax sugar (like operators,
\texttt{x\ +\ 1}) with more fundamental expressions (like function
calls, \texttt{+(x,\ 1)}), it also ``expands'' the macro, resulting in
code that is roughly:

\begin{Shaded}
\begin{Highlighting}[]
\NormalTok{x }\OperatorTok{=} \FloatTok{10}
\FunctionTok{println}\NormalTok{(}\StringTok{"x + 1 = "}\NormalTok{, x }\OperatorTok{+} \FloatTok{1}\NormalTok{)}
\end{Highlighting}
\end{Shaded}

\begin{verbatim}
x + 1 = 11
\end{verbatim}

Each macro can be thought of as like user-defined syntax sugar. We can
extend the language with our own ideas of how we can write programs.

\subsubsection{An even better way to
metaprogram}\label{an-even-better-way-to-metaprogram}

Another form of metaprogramming in Julia is multiple dispatch. Multiple
dispatch allows us to construct entire classes of programs that depend
on the \emph{types} of values --- and Julia will construct on-demand a
particular program \emph{specialized} on the inputs you actually
provide.

Next we'll learn more about Julia functions and how multiple dispatch
works. Afterwards we'll learn about abstract types and how to use this
metaprogramming approach.

\subsection{Functions in Julia}\label{functions-in-julia}

Functions are at the heart of Julia, and there's lots of ways of
expressing different functions. The
\href{https://docs.julialang.org/en/v1/manual/functions/}{Julia
documentation of functions} provides a rich description of all of the
details. We now overview a few special features that were perhaps not
evident from Units 1-3. There are also links to the documentation for
it.

\subsubsection{\texorpdfstring{\href{https://docs.julialang.org/en/v1/manual/functions/\#man-functions}{Function
syntaxes}}{Function syntaxes}}\label{function-syntaxes}

In Julia there are multiple ways of creating a function. There is the
``long form'':

\begin{Shaded}
\begin{Highlighting}[]
\KeywordTok{function} \FunctionTok{f1}\NormalTok{(x)}
   \ControlFlowTok{return}\NormalTok{ x}\OperatorTok{\^{}}\FloatTok{2}
\KeywordTok{end}
\end{Highlighting}
\end{Shaded}

\begin{verbatim}
f1 (generic function with 1 method)
\end{verbatim}

The \texttt{return} is optional in this case --- why?

There is a ``short form'':

\begin{Shaded}
\begin{Highlighting}[]
\FunctionTok{f2}\NormalTok{(x) }\OperatorTok{=}\NormalTok{ x}\OperatorTok{\^{}}\FloatTok{2}
\end{Highlighting}
\end{Shaded}

\begin{verbatim}
f2 (generic function with 1 method)
\end{verbatim}

And there are ``anonymous functions'' (or ``arrow functions''):

\begin{Shaded}
\begin{Highlighting}[]
\KeywordTok{const}\NormalTok{ f3 }\OperatorTok{=}\NormalTok{ x }\OperatorTok{{-}\textgreater{}}\NormalTok{ x}\OperatorTok{\^{}}\FloatTok{2}
\end{Highlighting}
\end{Shaded}

\begin{verbatim}
#2 (generic function with 1 method)
\end{verbatim}

In this form \texttt{x\ -\textgreater{}\ x\^{}2} is an expression that
constructs a function which is ``anonymous'' --- it has no name. In this
case we have bound it to the name \texttt{f3} so it can be referred to
later, but anonymous functions are most useful to write simple functions
inline within some other code. For example to pass to \texttt{map}:

\begin{Shaded}
\begin{Highlighting}[]
\FunctionTok{map}\NormalTok{(x }\OperatorTok{{-}\textgreater{}}\NormalTok{ x}\OperatorTok{\^{}}\FloatTok{2}\NormalTok{, [}\FloatTok{1}\NormalTok{,}\FloatTok{2}\NormalTok{,}\FloatTok{4}\NormalTok{])}
\end{Highlighting}
\end{Shaded}

\begin{verbatim}
3-element Vector{Int64}:
  1
  4
 16
\end{verbatim}

\paragraph{What type is a function?}\label{what-type-is-a-function}

You can think of the above functions as being a struct with zero fields.
Each function has a unique type:

\begin{Shaded}
\begin{Highlighting}[]
\FunctionTok{typeof}\NormalTok{(f1)}
\end{Highlighting}
\end{Shaded}

\begin{verbatim}
typeof(f1) (singleton type of function f1, subtype of Function)
\end{verbatim}

\begin{Shaded}
\begin{Highlighting}[]
\FunctionTok{typeof}\NormalTok{(f2)}
\end{Highlighting}
\end{Shaded}

\begin{verbatim}
typeof(f2) (singleton type of function f2, subtype of Function)
\end{verbatim}

\begin{Shaded}
\begin{Highlighting}[]
\FunctionTok{typeof}\NormalTok{(f1) }\OperatorTok{==} \FunctionTok{typeof}\NormalTok{(f2)}
\end{Highlighting}
\end{Shaded}

\begin{verbatim}
false
\end{verbatim}

And like \texttt{nothing}, there is no data associated with these:

\begin{Shaded}
\begin{Highlighting}[]
\FunctionTok{sizeof}\NormalTok{(f1)}
\end{Highlighting}
\end{Shaded}

\begin{verbatim}
0
\end{verbatim}

These types have only one possible value, which is why Julia describes
them as ``singleton types''. Singleton types have a one-to-one
correspondence between type and value.

\paragraph{Function methods}\label{function-methods}

Note that in Julia you can attach \emph{function methods} to any type
(primitive, struct, etc), not just those defined like above.

\begin{Shaded}
\begin{Highlighting}[]
\FunctionTok{methods}\NormalTok{(f1)}
\end{Highlighting}
\end{Shaded}

\begin{verbatim}
# 1 method for generic function "f1" from Main.Notebook:
 [1] f1(x)
     @ ~/nextcloud/MATH2504_2026/lecture/quarto/lecture-unit-5.qmd:600
\end{verbatim}

\begin{Shaded}
\begin{Highlighting}[]
\FunctionTok{methods}\NormalTok{(}\ConstantTok{nothing}\NormalTok{)}
\end{Highlighting}
\end{Shaded}

\begin{verbatim}
# 0 methods for callable object
\end{verbatim}

A ``method'' encodes the logic performed when the function executes.
Types without methods are not callable. Types can have more than one
method. We'll talk more about multiple methods when we discuss multiple
dispatch.

For now, just know you can attach methods to \emph{any} type with
special syntax like this:

\begin{Shaded}
\begin{Highlighting}[]
\NormalTok{(x}\OperatorTok{::}\DataTypeTok{Bool}\NormalTok{)() }\OperatorTok{=} \StringTok{"this is }\SpecialCharTok{$}\NormalTok{x}\StringTok{"}

\KeywordTok{function}\NormalTok{ (x}\OperatorTok{::}\DataTypeTok{Bool}\NormalTok{)()}
   \ControlFlowTok{return} \StringTok{"this is }\SpecialCharTok{$}\NormalTok{x}\StringTok{"}
\KeywordTok{end}
\end{Highlighting}
\end{Shaded}

For ``singleton'' types you can use the singleton value instead of the
type, so these are equivalent:

\begin{Shaded}
\begin{Highlighting}[]
\FunctionTok{nothing}\NormalTok{() }\OperatorTok{=} \StringTok{"Nothing else matters"}
\NormalTok{(}\OperatorTok{::}\DataTypeTok{Nothing}\NormalTok{)() }\OperatorTok{=} \StringTok{"Nothing else matters"}
\end{Highlighting}
\end{Shaded}

If this seems a bit crazy or abstract - don't worry! Just remember that
\emph{anything} can be made callable - functions are not special.

\subsubsection{Nested functions:
``closures''}\label{nested-functions-closures}

In Julia you can define functions within functions.

\begin{Shaded}
\begin{Highlighting}[]
\KeywordTok{function} \FunctionTok{adder}\NormalTok{(x)}
   \ControlFlowTok{return}\NormalTok{ y }\OperatorTok{{-}\textgreater{}}\NormalTok{ y }\OperatorTok{+}\NormalTok{ x}
\KeywordTok{end}
\end{Highlighting}
\end{Shaded}

\begin{verbatim}
adder (generic function with 1 method)
\end{verbatim}

This function returns another function! We can use the output to add
things.

\begin{Shaded}
\begin{Highlighting}[]
\NormalTok{add\_2 }\OperatorTok{=} \FunctionTok{adder}\NormalTok{(}\FloatTok{2}\NormalTok{)}
\FunctionTok{add\_2}\NormalTok{(}\FloatTok{3}\NormalTok{)}
\end{Highlighting}
\end{Shaded}

\begin{verbatim}
5
\end{verbatim}

Such functions are called ``closures'' because they ``close over'' (or
``capture'') the values of variables outside their scope --- in this
case it captured the value of \texttt{x}, \texttt{2}, and implicitly
stores it inside \texttt{add\_2}.

\begin{Shaded}
\begin{Highlighting}[]
\NormalTok{add\_2.x}
\end{Highlighting}
\end{Shaded}

\begin{verbatim}
2
\end{verbatim}

\begin{Shaded}
\begin{Highlighting}[]
\FunctionTok{sizeof}\NormalTok{(add\_2)}
\end{Highlighting}
\end{Shaded}

\begin{verbatim}
8
\end{verbatim}

We can see the full internal structure using the \texttt{dump()}
function:

\begin{Shaded}
\begin{Highlighting}[]
\FunctionTok{dump}\NormalTok{(add\_2)}
\end{Highlighting}
\end{Shaded}

\begin{verbatim}
#adder##0 (function of type var"#adder##0#adder##1"{Int64})
  x: Int64 2
\end{verbatim}

Note again it doesn't matter which form of function syntax we use. It's
just different syntax for the same thing. We could have equally written:

\begin{Shaded}
\begin{Highlighting}[]
\KeywordTok{function} \FunctionTok{adder}\NormalTok{(x)}
   \ControlFlowTok{return} \KeywordTok{function}\NormalTok{ (y)}
   \ControlFlowTok{return}\NormalTok{ y }\OperatorTok{+}\NormalTok{ x}
   \KeywordTok{end}
\KeywordTok{end}
\end{Highlighting}
\end{Shaded}

\begin{verbatim}
adder (generic function with 1 method)
\end{verbatim}

In Julia, closures are given their own types and methods. The closure is
just an automatically generated \texttt{struct} datatype with a method
attached.

\begin{Shaded}
\begin{Highlighting}[]
\FunctionTok{typeof}\NormalTok{(add\_2)}
\end{Highlighting}
\end{Shaded}

\begin{verbatim}
var"#adder##0#adder##1"{Int64}
\end{verbatim}

\begin{Shaded}
\begin{Highlighting}[]
\FunctionTok{methods}\NormalTok{(add\_2)}
\end{Highlighting}
\end{Shaded}

\begin{verbatim}
# 1 method for anonymous function "#adder##0":
 [1] (::Main.Notebook.var"#adder##0#adder##1")(y)
     @ ~/nextcloud/MATH2504_2026/lecture/quarto/lecture-unit-5.qmd:689
\end{verbatim}

Because closures are ``just'' a convenient way to do things you can
already do in Julia with structs and methods, you can think of them as
an advanced type of syntax sugar. The above is equivalent to this:

\begin{Shaded}
\begin{Highlighting}[]
\KeywordTok{struct}\NormalTok{ Adder\{T\}}
\NormalTok{   x}\OperatorTok{::}\DataTypeTok{T}
\KeywordTok{end}

\KeywordTok{function}\NormalTok{ (a}\OperatorTok{::}\DataTypeTok{Adder}\NormalTok{)(y)}
   \ControlFlowTok{return}\NormalTok{ a.x }\OperatorTok{+}\NormalTok{ y}
\KeywordTok{end}

\KeywordTok{function} \FunctionTok{adder}\NormalTok{(x)}
   \ControlFlowTok{return} \FunctionTok{Adder}\NormalTok{(x)}
\KeywordTok{end}

\NormalTok{add\_2 }\OperatorTok{=} \FunctionTok{adder}\NormalTok{(}\FloatTok{2}\NormalTok{)}
\end{Highlighting}
\end{Shaded}

\begin{verbatim}
Adder{Int64}(2)
\end{verbatim}

\begin{Shaded}
\begin{Highlighting}[]
\FunctionTok{add\_2}\NormalTok{(}\FloatTok{3}\NormalTok{)}
\end{Highlighting}
\end{Shaded}

\begin{verbatim}
5
\end{verbatim}

Note that \emph{every} function in Julia has its own type ---
e.g.~\texttt{sin} and \texttt{cos} have different types. We'll come back
to types and methods later.

\subsubsection{\texorpdfstring{\href{https://docs.julialang.org/en/v1/manual/functions/\#Argument-type-declarations}{Argument
types}}{Argument types}}\label{argument-types}

To constrain the type of a function argument we use the \texttt{::}
operator: When written in the argument list, \texttt{x::T} means the
argument named \texttt{x} must be of type \texttt{T} for this method to
match.

We have seen Julia methods from the start, e.g.

\begin{Shaded}
\begin{Highlighting}[]
\FunctionTok{f}\NormalTok{(x}\OperatorTok{::}\DataTypeTok{Int}\NormalTok{) }\OperatorTok{=}\NormalTok{ x}\OperatorTok{\^{}}\FloatTok{2}
\end{Highlighting}
\end{Shaded}

\begin{verbatim}
f (generic function with 2 methods)
\end{verbatim}

or,

\begin{Shaded}
\begin{Highlighting}[]
\FunctionTok{f}\NormalTok{(x}\OperatorTok{::}\DataTypeTok{Number}\NormalTok{) }\OperatorTok{=}\NormalTok{ x}\OperatorTok{\^{}}\FloatTok{2}
\end{Highlighting}
\end{Shaded}

\begin{verbatim}
f (generic function with 3 methods)
\end{verbatim}

or simply,

\begin{Shaded}
\begin{Highlighting}[]
\FunctionTok{f}\NormalTok{(x) }\OperatorTok{=}\NormalTok{ x}\OperatorTok{\^{}}\FloatTok{2}
\end{Highlighting}
\end{Shaded}

\begin{verbatim}
f (generic function with 4 methods)
\end{verbatim}

Remember that \texttt{x} and \texttt{x::Any} are the same thing. The
final method will accept \emph{any} inputs. Whether or not
\texttt{x\^{}2} works or throws an error depends on \texttt{x}.

In a sense, the third method isn't \emph{safe} since someone could
provide an input like \texttt{(1,\ 2,\ 3)} that can't be squared. What
happens if someone provides a string? Is this expected? On the other
hand, the third method provides maximum flexibility --- someone can
introduce a new type that works with \texttt{f} at any point in time,
without having to edit the definition of \texttt{f} to make it work.

Note that there are two other uses of the \texttt{::} operator outside
of function arguments - type assertion and typed variable declarations -
but these aren't used as commonly and we won't cover them here.

\subsubsection{Dispatch}\label{dispatch}

``Dispatch'' is the process of determining which method of a function is
called.

\begin{Shaded}
\begin{Highlighting}[]
\FunctionTok{example\_f}\NormalTok{(x}\OperatorTok{::}\DataTypeTok{Int}\NormalTok{) }\OperatorTok{=} \StringTok{"input is an integer"}
\FunctionTok{example\_f}\NormalTok{(x}\OperatorTok{::}\DataTypeTok{String}\NormalTok{) }\OperatorTok{=} \StringTok{"input is a string"}

\PreprocessorTok{@show} \FunctionTok{example\_f}\NormalTok{(}\FloatTok{1}\NormalTok{)}
\PreprocessorTok{@show} \FunctionTok{example\_f}\NormalTok{(}\StringTok{"abc"}\NormalTok{)}

\FunctionTok{methods}\NormalTok{(example\_f)}
\end{Highlighting}
\end{Shaded}

\begin{verbatim}
example_f(1) = "input is an integer"
example_f("abc") = "input is a string"
\end{verbatim}

\begin{verbatim}
# 2 methods for generic function "example_f" from Main.Notebook:
 [1] example_f(x::String)
     @ ~/nextcloud/MATH2504_2026/lecture/quarto/lecture-unit-5.qmd:798
 [2] example_f(x::Int64)
     @ ~/nextcloud/MATH2504_2026/lecture/quarto/lecture-unit-5.qmd:797
\end{verbatim}

You can use it to tailor the logic of a function to its inputs. Try type
\texttt{methods(+)} at the REPL!

Argument types can be abstract types, like \texttt{Any} or
\texttt{Union}.

\begin{Shaded}
\begin{Highlighting}[]
\FunctionTok{example\_f2}\NormalTok{(x}\OperatorTok{::}\DataTypeTok{Int}\NormalTok{) }\OperatorTok{=} \StringTok{"input is an integer"}
\FunctionTok{example\_f2}\NormalTok{(x}\OperatorTok{::}\DataTypeTok{Any}\NormalTok{) }\OperatorTok{=} \StringTok{"input is not an integer"}

\PreprocessorTok{@show} \FunctionTok{example\_f2}\NormalTok{(}\FloatTok{1}\NormalTok{)}
\PreprocessorTok{@show} \FunctionTok{example\_f2}\NormalTok{(}\StringTok{"abc"}\NormalTok{)}

\FunctionTok{methods}\NormalTok{(example\_f2)}
\end{Highlighting}
\end{Shaded}

\begin{verbatim}
example_f2(1) = "input is an integer"
example_f2("abc") = "input is not an integer"
\end{verbatim}

\begin{verbatim}
# 2 methods for generic function "example_f2" from Main.Notebook:
 [1] example_f2(x::Int64)
     @ ~/nextcloud/MATH2504_2026/lecture/quarto/lecture-unit-5.qmd:811
 [2] example_f2(x)
     @ ~/nextcloud/MATH2504_2026/lecture/quarto/lecture-unit-5.qmd:812
\end{verbatim}

Note that \texttt{Int\ \textless{}:\ Any} and the two methods overlap
somewhat. In this case Julia will pick the most ``specific'' method. In
terms of ``types as sets of values'' it matches to the ``smallest'' set.

The term ``multiple dispatch'' means that dispatch can depend on the
types of multiple arguments:

\begin{Shaded}
\begin{Highlighting}[]
\FunctionTok{example\_f3}\NormalTok{(x}\OperatorTok{::}\DataTypeTok{Int}\NormalTok{, y}\OperatorTok{::}\DataTypeTok{Int}\NormalTok{) }\OperatorTok{=} \StringTok{"both input are integers"}
\FunctionTok{example\_f3}\NormalTok{(x}\OperatorTok{::}\DataTypeTok{String}\NormalTok{, y}\OperatorTok{::}\DataTypeTok{Int}\NormalTok{) }\OperatorTok{=} \StringTok{"the inputs have different types"}
\FunctionTok{example\_f3}\NormalTok{(x}\OperatorTok{::}\DataTypeTok{Int}\NormalTok{, y}\OperatorTok{::}\DataTypeTok{String}\NormalTok{) }\OperatorTok{=} \StringTok{"the inputs have different types"}
\FunctionTok{example\_f3}\NormalTok{(x}\OperatorTok{::}\DataTypeTok{String}\NormalTok{, y}\OperatorTok{::}\DataTypeTok{String}\NormalTok{) }\OperatorTok{=} \StringTok{"both input are strings"}

\PreprocessorTok{@show} \FunctionTok{example\_f3}\NormalTok{(}\FloatTok{1}\NormalTok{, }\FloatTok{2}\NormalTok{);}
\PreprocessorTok{@show} \FunctionTok{example\_f3}\NormalTok{(}\FloatTok{1}\NormalTok{, }\StringTok{"abc"}\NormalTok{);}

\FunctionTok{methods}\NormalTok{(example\_f3)}
\end{Highlighting}
\end{Shaded}

\begin{verbatim}
example_f3(1, 2) = "both input are integers"
example_f3(1, "abc") = "the inputs have different types"
\end{verbatim}

\begin{verbatim}
# 4 methods for generic function "example_f3" from Main.Notebook:
 [1] example_f3(x::String, y::String)
     @ ~/nextcloud/MATH2504_2026/lecture/quarto/lecture-unit-5.qmd:828
 [2] example_f3(x::Int64, y::String)
     @ ~/nextcloud/MATH2504_2026/lecture/quarto/lecture-unit-5.qmd:827
 [3] example_f3(x::String, y::Int64)
     @ ~/nextcloud/MATH2504_2026/lecture/quarto/lecture-unit-5.qmd:826
 [4] example_f3(x::Int64, y::Int64)
     @ ~/nextcloud/MATH2504_2026/lecture/quarto/lecture-unit-5.qmd:825
\end{verbatim}

In object-oriented languages (like Python, Javascript, C\#, Java or C++)
dynamic dispatch occurs only on the type of first object (the class
method being called). Multiple dispatch is a lot more flexible and
changes the way we design software.

\subsubsection{\texorpdfstring{\href{https://docs.julialang.org/en/v1/manual/functions/\#Varargs-Functions}{Varargs
functions}}{Varargs functions}}\label{varargs-functions}

Sometimes functions have a variable number of arguments --- sometimes
called ``varargs'' (or ``variadic functions'').

A simple thing to do might be to add up all the inputs:

\begin{Shaded}
\begin{Highlighting}[]
\KeywordTok{function} \FunctionTok{add\_all}\NormalTok{(inputs}\OperatorTok{...}\NormalTok{)}
\NormalTok{   out }\OperatorTok{=} \FloatTok{0}
   \ControlFlowTok{for}\NormalTok{ x }\KeywordTok{in}\NormalTok{ inputs}
\NormalTok{   out }\OperatorTok{+=}\NormalTok{ x}
   \ControlFlowTok{end}
   \ControlFlowTok{return}\NormalTok{ out}
\KeywordTok{end}

\FunctionTok{add\_all}\NormalTok{(}\FloatTok{1}\NormalTok{, }\FloatTok{2}\NormalTok{, }\FloatTok{3}\NormalTok{, }\FloatTok{4}\NormalTok{, }\FloatTok{5}\NormalTok{)}
\end{Highlighting}
\end{Shaded}

\begin{verbatim}
15
\end{verbatim}

\begin{Shaded}
\begin{Highlighting}[]
\KeywordTok{function} \FunctionTok{polynomialGenerator}\NormalTok{(a}\OperatorTok{...}\NormalTok{)}
\NormalTok{   n }\OperatorTok{=} \FunctionTok{length}\NormalTok{(a) }\OperatorTok{{-}} \FloatTok{1}
\NormalTok{   poly }\OperatorTok{=} \KeywordTok{function}\NormalTok{(x)}
   \ControlFlowTok{return} \FunctionTok{sum}\NormalTok{(a[i}\OperatorTok{+}\FloatTok{1}\NormalTok{] }\OperatorTok{*}\NormalTok{ x}\OperatorTok{\^{}}\NormalTok{i }\ControlFlowTok{for}\NormalTok{ i }\KeywordTok{in} \FloatTok{0}\OperatorTok{:}\NormalTok{n)}
   \ControlFlowTok{end}
   \ControlFlowTok{return}\NormalTok{ poly}
\KeywordTok{end}

\NormalTok{polynomial }\OperatorTok{=} \FunctionTok{polynomialGenerator}\NormalTok{(}\FloatTok{1}\NormalTok{, }\FloatTok{3}\NormalTok{, }\OperatorTok{{-}}\FloatTok{10}\NormalTok{)}

\NormalTok{[}\FunctionTok{polynomial}\NormalTok{(}\OperatorTok{{-}}\FloatTok{1}\NormalTok{), }\FunctionTok{polynomial}\NormalTok{(}\FloatTok{0}\NormalTok{), }\FunctionTok{polynomial}\NormalTok{(}\FloatTok{1}\NormalTok{)]}
\end{Highlighting}
\end{Shaded}

\begin{verbatim}
3-element Vector{Int64}:
 -12
   1
  -6
\end{verbatim}

You could use then use \texttt{Roots} package to automatically find out
where all the inputs produce zero output:

\begin{Shaded}
\begin{Highlighting}[]
\ImportTok{using} \BuiltInTok{Roots}

\NormalTok{zero\_vals }\OperatorTok{=} \FunctionTok{find\_zeros}\NormalTok{(polynomial, }\OperatorTok{{-}}\FloatTok{10}\NormalTok{, }\FloatTok{10}\NormalTok{)}
\FunctionTok{println}\NormalTok{(}\StringTok{"Zeros of the function f(x): "}\NormalTok{, zero\_vals)}
\end{Highlighting}
\end{Shaded}

\begin{verbatim}
Zeros of the function f(x): [-0.19999999999999998, 0.5]
\end{verbatim}

Here you start to see how Julia features like varargs functions,
closures, and packages let you solve a wide range of problems without
writing very much code.

\subsubsection{\texorpdfstring{\href{https://docs.julialang.org/en/v1/manual/functions/\#Optional-Arguments}{Optional
arguments}}{Optional arguments}}\label{optional-arguments}

You can have optional arguments (or default values):

\begin{Shaded}
\begin{Highlighting}[]
\ImportTok{using} \BuiltInTok{Distributions}

\KeywordTok{function} \FunctionTok{my\_density}\NormalTok{(x}\OperatorTok{::}\DataTypeTok{Float64}\NormalTok{, μ}\OperatorTok{::}\DataTypeTok{Float64 }\OperatorTok{=} \FloatTok{0.0}\NormalTok{, σ}\OperatorTok{::}\DataTypeTok{Float64 }\OperatorTok{=} \FloatTok{1.0}\NormalTok{)}
   \ControlFlowTok{return} \FunctionTok{exp}\NormalTok{(}\FunctionTok{{-}}\NormalTok{(x}\OperatorTok{{-}}\NormalTok{μ)}\OperatorTok{\^{}}\FloatTok{2} \OperatorTok{/}\NormalTok{ (}\FloatTok{2}\NormalTok{σ}\OperatorTok{\^{}}\FloatTok{2}\NormalTok{)) }\OperatorTok{/}\NormalTok{ (}\FunctionTok{σ*√}\NormalTok{(}\FloatTok{2}\NormalTok{π))}
\KeywordTok{end}

\NormalTok{x }\OperatorTok{=} \FloatTok{1.5}
\PreprocessorTok{@show} \FunctionTok{pdf}\NormalTok{(}\FunctionTok{Normal}\NormalTok{(), x) }\FunctionTok{my\_density}\NormalTok{(x)}
\PreprocessorTok{@show} \FunctionTok{pdf}\NormalTok{(}\FunctionTok{Normal}\NormalTok{(}\FloatTok{0.5}\NormalTok{), x) }\FunctionTok{my\_density}\NormalTok{(x, }\FloatTok{0.5}\NormalTok{);}
\end{Highlighting}
\end{Shaded}

\begin{verbatim}
pdf(Normal(), x) = 0.12951759566589174
my_density(x) = 0.12951759566589174
pdf(Normal(0.5), x) = 0.24197072451914337
my_density(x, 0.5) = 0.24197072451914337
\end{verbatim}

\subsubsection{\texorpdfstring{\href{https://docs.julialang.org/en/v1/manual/functions/\#Keyword-Arguments}{Keyword
arguments}}{Keyword arguments}}\label{keyword-arguments}

Arguments following the \texttt{;} character in the function definition
or the function call are called \textbf{keyword arguments}. These are
named as they are used.

\begin{Shaded}
\begin{Highlighting}[]
\KeywordTok{function} \FunctionTok{my\_density}\NormalTok{(x}\OperatorTok{::}\DataTypeTok{Float64}\NormalTok{; μ}\OperatorTok{::}\DataTypeTok{Float64 }\OperatorTok{=} \FloatTok{0.0}\NormalTok{, σ}\OperatorTok{::}\DataTypeTok{Float64 }\OperatorTok{=} \FloatTok{1.0}\NormalTok{)}
   \ControlFlowTok{return} \FunctionTok{exp}\NormalTok{(}\FunctionTok{{-}}\NormalTok{(x}\OperatorTok{{-}}\NormalTok{μ)}\OperatorTok{\^{}}\FloatTok{2}\OperatorTok{/}\NormalTok{(}\FloatTok{2}\NormalTok{σ}\OperatorTok{\^{}}\FloatTok{2}\NormalTok{) ) }\OperatorTok{/}\NormalTok{ (}\FunctionTok{σ*√}\NormalTok{(}\FloatTok{2}\NormalTok{π)) }
\KeywordTok{end}

\PreprocessorTok{@show} \FunctionTok{pdf}\NormalTok{(}\FunctionTok{Normal}\NormalTok{(}\FloatTok{0.0}\NormalTok{,}\FloatTok{2.5}\NormalTok{),x) }\FunctionTok{my\_density}\NormalTok{(x, σ}\OperatorTok{=}\FloatTok{2.5}\NormalTok{);}
\end{Highlighting}
\end{Shaded}

\begin{verbatim}
pdf(Normal(0.0, 2.5), x) = 0.13328984115671988
my_density(x, σ = 2.5) = 0.13328984115671988
\end{verbatim}

We can even make a function that takes arbitrary numbers of positional
and keyword arguments:

\begin{Shaded}
\begin{Highlighting}[]
\KeywordTok{function} \FunctionTok{very\_flexible\_function}\NormalTok{(args}\OperatorTok{...}\NormalTok{; kwargs}\OperatorTok{...}\NormalTok{)}
   \PreprocessorTok{@info} \StringTok{"All the args"}\NormalTok{ args kwargs}
\KeywordTok{end}

\FunctionTok{very\_flexible\_function}\NormalTok{(}\FloatTok{2.5}\NormalTok{, }\ConstantTok{false}\NormalTok{, a}\OperatorTok{=}\FloatTok{1}\NormalTok{, b}\OperatorTok{=}\StringTok{"two"}\NormalTok{, c}\OperatorTok{=:}\NormalTok{three)}
\end{Highlighting}
\end{Shaded}

\begin{verbatim}
┌ Info: All the args
│   args = (2.5, false)
│   kwargs =
│    pairs(::NamedTuple) with 3 entries:
│      :a => 1
│      :b => "two"
└      :c => :three
\end{verbatim}

The standard non-keyword arguments are called ``positional'' arguments
since they are identified by their position not their name.

Note that when \emph{calling} a function you may either put a \texttt{;}
or a \texttt{,} before the keyword arguments. In a function definition
there's an ambiguity between optional arguments and keyword

\subsubsection{\texorpdfstring{\href{https://docs.julialang.org/en/v1/manual/functions/\#Argument-destructuring}{Destructing}}{Destructing}}\label{destructing}

In Julia you can \emph{destructure} tuples (and other iterables) in the
reverse way you \emph{construct} them:

\begin{Shaded}
\begin{Highlighting}[]
\NormalTok{my\_tuple }\OperatorTok{=}\NormalTok{ (}\FloatTok{42}\NormalTok{, }\StringTok{"abc"}\NormalTok{)  }\CommentTok{\# construct}
\NormalTok{(a, b) }\OperatorTok{=}\NormalTok{ my\_tuple       }\CommentTok{\# destruct}

\PreprocessorTok{@show}\NormalTok{ a}
\PreprocessorTok{@show}\NormalTok{ b;}
\end{Highlighting}
\end{Shaded}

\begin{verbatim}
a = 42
b = "abc"
\end{verbatim}

A really common pattern here is swapping two variables:

\begin{Shaded}
\begin{Highlighting}[]
\NormalTok{x }\OperatorTok{=} \FloatTok{1}
\NormalTok{y }\OperatorTok{=} \FloatTok{2}
\NormalTok{(x, y) }\OperatorTok{=}\NormalTok{ (y, x)}
\PreprocessorTok{@show}\NormalTok{ x}
\PreprocessorTok{@show}\NormalTok{ y;}
\end{Highlighting}
\end{Shaded}

\begin{verbatim}
x = 2
y = 1
\end{verbatim}

You can also destructure fields of structs and named tuples in the
reverse way you create a named tuple:

\begin{Shaded}
\begin{Highlighting}[]
\KeywordTok{function} \FunctionTok{f}\NormalTok{(x}\OperatorTok{::}\DataTypeTok{Complex}\NormalTok{)}
\NormalTok{   (; re, }\ConstantTok{im}\NormalTok{) }\OperatorTok{=}\NormalTok{ x}
   \FunctionTok{println}\NormalTok{(}\StringTok{"Real part: }\SpecialCharTok{$}\NormalTok{re}\StringTok{"}\NormalTok{)}
   \FunctionTok{println}\NormalTok{(}\StringTok{"Imaginary part: }\SpecialCharTok{$}\NormalTok{im}\StringTok{"}\NormalTok{)}
\KeywordTok{end}

\FunctionTok{f}\NormalTok{(}\FloatTok{10} \OperatorTok{+} \FloatTok{42im}\NormalTok{)}
\end{Highlighting}
\end{Shaded}

\begin{verbatim}
-1664 + 840im
\end{verbatim}

(But note that the \texttt{x.re} and \texttt{x.im} fields are internal -
you should use the \texttt{real()} and \texttt{imag()} functions to work
with \texttt{Complex} values!)

\subsubsection{\texorpdfstring{\href{https://docs.julialang.org/en/v1/manual/functions/\#Function-composition-and-piping}{Function
composition and
piping}}{Function composition and piping}}\label{function-composition-and-piping}

Just a little bit of syntax sugar for ``piping'' values into function
calls:

\begin{Shaded}
\begin{Highlighting}[]
\ConstantTok{π}\OperatorTok{/}\FloatTok{4} \OperatorTok{|\textgreater{}}\NormalTok{ cos }\OperatorTok{|\textgreater{}}\NormalTok{ acos }\OperatorTok{|\textgreater{}}\NormalTok{ x }\OperatorTok{{-}\textgreater{}} \FloatTok{4}\NormalTok{x}
\end{Highlighting}
\end{Shaded}

\begin{verbatim}
3.141592653589793
\end{verbatim}

Here is a function just like \texttt{identity}:

\begin{Shaded}
\begin{Highlighting}[]
\KeywordTok{const}\NormalTok{ ii }\OperatorTok{=}\NormalTok{ cos }\OperatorTok{∘}\NormalTok{ acos }\CommentTok{\#\textbackslash{}circ + [TAB]}
\NormalTok{(}\FunctionTok{ii}\NormalTok{(}\ConstantTok{π}\OperatorTok{/}\FloatTok{4}\NormalTok{), }\ConstantTok{π}\OperatorTok{/}\FloatTok{4}\NormalTok{)}
\end{Highlighting}
\end{Shaded}

\begin{verbatim}
(0.7853981633974483, 0.7853981633974483)
\end{verbatim}

\subsubsection{\texorpdfstring{\href{https://docs.julialang.org/en/v1/manual/functions/\#man-vectorized}{Dot
syntax for broadcasting
functions}}{Dot syntax for broadcasting functions}}\label{dot-syntax-for-broadcasting-functions}

You already know the broadcast operator. It is syntax sugar for the
\texttt{broadcast} function.

\begin{Shaded}
\begin{Highlighting}[]
\NormalTok{x\_range }\OperatorTok{=} \FloatTok{0}\OperatorTok{:}\FloatTok{0.5}\OperatorTok{:}\ConstantTok{π}
\FunctionTok{cos}\NormalTok{.(x\_range)}
\end{Highlighting}
\end{Shaded}

\begin{verbatim}
7-element Vector{Float64}:
  1.0
  0.8775825618903728
  0.5403023058681398
  0.0707372016677029
 -0.4161468365471424
 -0.8011436155469337
 -0.9899924966004454
\end{verbatim}

See the docs for a discussion of performance of broadcasting.

You can also use the macro \texttt{@.}

\begin{Shaded}
\begin{Highlighting}[]
\NormalTok{x\_range }\OperatorTok{=} \FloatTok{0}\OperatorTok{:}\FloatTok{0.5}\OperatorTok{:}\ConstantTok{π}
\NormalTok{@. }\FunctionTok{cos}\NormalTok{(x\_range }\OperatorTok{+} \FloatTok{2}\NormalTok{)}\OperatorTok{\^{}}\FloatTok{2}
\end{Highlighting}
\end{Shaded}

\begin{verbatim}
7-element Vector{Float64}:
 0.17317818956819406
 0.6418310927316131
 0.9800851433251829
 0.8769511271716524
 0.4272499830956933
 0.0444348690576615
 0.08046423546177377
\end{verbatim}

\subsection{More items from control
flow}\label{more-items-from-control-flow}

You are already very familiar with conditional statements (\texttt{if},
\texttt{elseif}, \texttt{else}), with loops (\texttt{for} and
\texttt{while}), with short circuit evaluation, and with many other
variants. For example you can use \texttt{continue} and \texttt{break}
in loops, and Julia even supports \texttt{@goto} and \texttt{@label} for
when that isn't enough.

You can find more control flow details here:
\href{https://docs.julialang.org/en/v1/manual/control-flow/}{control
flow in Julia docs}. We'll just dive into \texttt{for} loops and error
handling below.

\subsubsection{\texorpdfstring{Iteration and \texttt{for}
loops}{Iteration and for loops}}\label{iteration-and-for-loops}

One impressive piece of syntax sugar is the \texttt{for} loop. It is in
fact a special case of a \texttt{while} loop, like so:

\begin{Shaded}
\begin{Highlighting}[]
\CommentTok{\# this loop...}
\ControlFlowTok{for}\NormalTok{ x }\KeywordTok{in}\NormalTok{ iterable}
   \CommentTok{\# \textless{}CODE\textgreater{}}
\ControlFlowTok{end}

\CommentTok{\# ...is equivalent to}
\NormalTok{tmp }\OperatorTok{=} \FunctionTok{iterate}\NormalTok{(iterable)}
\ControlFlowTok{while}\NormalTok{ tmp }\OperatorTok{!==} \ConstantTok{nothing}
\NormalTok{   (x, state) }\OperatorTok{=}\NormalTok{ tmp}

   \CommentTok{\# \textless{}CODE\textgreater{}}
   
\NormalTok{   tmp }\OperatorTok{=} \FunctionTok{iterate}\NormalTok{(iterable, state)}
\ControlFlowTok{end}
\end{Highlighting}
\end{Shaded}

In the above, \texttt{tmp} is either \texttt{nothing} or it's a tuple of
length two --- containing the next value and whatever metadata is needed
to iterate to the next element.

You can make any of your data types iterable in a \texttt{for} loop by
implementing a method on the \texttt{iterate} function. We won't be
doing this in this course, but it is an excellent example of an
\emph{interface} in Julia. We'll talk more about interfaces later.

\subsubsection{\texorpdfstring{\href{https://docs.julialang.org/en/v1/manual/control-flow/\#Exception-Handling}{Errors
and exceptions}}{Errors and exceptions}}\label{errors-and-exceptions}

One additional thing to know about is exception handling.

\begin{Shaded}
\begin{Highlighting}[]
\KeywordTok{function} \FunctionTok{my\_2\_by\_2\_inv}\NormalTok{(A}\OperatorTok{::}\DataTypeTok{Matrix\{Float64\}}\NormalTok{)}
   \ControlFlowTok{if} \FunctionTok{size}\NormalTok{(A) }\OperatorTok{!=}\NormalTok{ (}\FloatTok{2}\NormalTok{,}\FloatTok{2}\NormalTok{)}
   \FunctionTok{error}\NormalTok{(}\StringTok{"This function only works for 2x2 matrices"}\NormalTok{)}
   \ControlFlowTok{end}

\NormalTok{   d }\OperatorTok{=}\NormalTok{ A[}\FloatTok{1}\NormalTok{,}\FloatTok{1}\NormalTok{]}\OperatorTok{*}\NormalTok{A[}\FloatTok{2}\NormalTok{,}\FloatTok{2}\NormalTok{] }\OperatorTok{{-}}\NormalTok{ A[}\FloatTok{2}\NormalTok{,}\FloatTok{1}\NormalTok{]}\OperatorTok{*}\NormalTok{A[}\FloatTok{1}\NormalTok{,}\FloatTok{2}\NormalTok{]}

   \ControlFlowTok{if}\NormalTok{ d }\OperatorTok{≈} \FloatTok{0}
   \FunctionTok{throw}\NormalTok{(}\FunctionTok{ArgumentError}\NormalTok{(}\StringTok{"matrix is singular or near singular"}\NormalTok{)) }\CommentTok{\#\textbackslash{}approx + [TAB]}
   \ControlFlowTok{end}

   \ControlFlowTok{return}\NormalTok{ [A[}\FloatTok{2}\NormalTok{,}\FloatTok{2}\NormalTok{] }\OperatorTok{{-}}\NormalTok{A[}\FloatTok{1}\NormalTok{,}\FloatTok{2}\NormalTok{]; }\OperatorTok{{-}}\NormalTok{A[}\FloatTok{2}\NormalTok{,}\FloatTok{1}\NormalTok{] A[}\FloatTok{1}\NormalTok{,}\FloatTok{1}\NormalTok{]] }\OperatorTok{./}\NormalTok{ d}
\KeywordTok{end}

\FunctionTok{my\_2\_by\_2\_inv}\NormalTok{(}\FunctionTok{rand}\NormalTok{(}\FloatTok{3}\NormalTok{,}\FloatTok{3}\NormalTok{))}
\end{Highlighting}
\end{Shaded}

\begin{Highlighting}
\textcolor{black}{ErrorException: ErrorException("This function only works for 2x2 matrices")}
\textcolor{black}{This function only works for 2x2 matrices}
\textcolor{black}{Stacktrace:}
\textcolor{black}{ [1] }\textcolor{QuartoInternalColor1}{}\textcolor{QuartoInternalColor1}{error}\textcolor{QuartoInternalColor1}{}\textcolor{QuartoInternalColor1}{}\textcolor{QuartoInternalColor1}{(}\textcolor{QuartoInternalColor1}{}\textcolor{QuartoInternalColor1}{s}\textcolor{QuartoInternalColor1}{::}\textcolor{QuartoInternalColor1}{String}\textcolor{QuartoInternalColor1}{}\textcolor{QuartoInternalColor1}{)}\textcolor{QuartoInternalColor1}{}
\textcolor{QuartoInternalColor1}{}\textcolor{QuartoInternalColor1}{   @}\textcolor{QuartoInternalColor1}{ }\textcolor{QuartoInternalColor1}{Base}\textcolor{QuartoInternalColor1}{ }\textcolor{QuartoInternalColor1}{./}\textcolor{QuartoInternalColor1}{}\textcolor{QuartoInternalColor1}{}\textcolor{QuartoInternalColor1}{error.jl:44}\textcolor{QuartoInternalColor1}{}\textcolor{QuartoInternalColor1}{}
\textcolor{QuartoInternalColor1}{ [2] }\textcolor{QuartoInternalColor1}{}\textcolor{QuartoInternalColor1}{my\_2\_by\_2\_inv}\textcolor{QuartoInternalColor1}{}\textcolor{QuartoInternalColor1}{}\textcolor{QuartoInternalColor1}{(}\textcolor{QuartoInternalColor1}{}\textcolor{QuartoInternalColor1}{A}\textcolor{QuartoInternalColor1}{::}\textcolor{QuartoInternalColor1}{Matrix}\textcolor{QuartoInternalColor1}{\{Float64\}}\textcolor{QuartoInternalColor1}{}\textcolor{QuartoInternalColor1}{}\textcolor{QuartoInternalColor1}{)}\textcolor{QuartoInternalColor1}{}
\textcolor{QuartoInternalColor1}{}\textcolor{QuartoInternalColor1}{   @}\textcolor{QuartoInternalColor1}{ }\textcolor{QuartoInternalColor2}{Main.Notebook}\textcolor{QuartoInternalColor1}{ }\textcolor{QuartoInternalColor1}{\textasciitilde{}/nextcloud/MATH2504\_2026/lecture/quarto/}\textcolor{QuartoInternalColor1}{}\textcolor{QuartoInternalColor1}{}\textcolor{QuartoInternalColor1}{lecture-unit-5.qmd:1032}\textcolor{QuartoInternalColor1}{}\textcolor{QuartoInternalColor1}{}
\textcolor{QuartoInternalColor1}{ [3] top-level scope}
\textcolor{QuartoInternalColor1}{}\textcolor{QuartoInternalColor1}{   @}\textcolor{QuartoInternalColor1}{ }\textcolor{QuartoInternalColor1}{\textasciitilde{}/nextcloud/MATH2504\_2026/lecture/quarto/}\textcolor{QuartoInternalColor1}{}\textcolor{QuartoInternalColor1}{}\textcolor{QuartoInternalColor1}{lecture-unit-5.qmd:1044}\textcolor{QuartoInternalColor1}{}\textcolor{QuartoInternalColor1}{}
\end{Highlighting}

\begin{Shaded}
\begin{Highlighting}[]
\FunctionTok{my\_2\_by\_2\_inv}\NormalTok{(}\FunctionTok{ones}\NormalTok{(}\FloatTok{2}\NormalTok{,}\FloatTok{2}\NormalTok{))}
\end{Highlighting}
\end{Shaded}

\begin{Highlighting}
\textcolor{black}{ArgumentError: ArgumentError("matrix is singular or near singular")}
\textcolor{black}{ArgumentError: matrix is singular or near singular}
\textcolor{black}{Stacktrace:}
\textcolor{black}{ [1] }\textcolor{QuartoInternalColor1}{}\textcolor{QuartoInternalColor1}{my\_2\_by\_2\_inv}\textcolor{QuartoInternalColor1}{}\textcolor{QuartoInternalColor1}{}\textcolor{QuartoInternalColor1}{(}\textcolor{QuartoInternalColor1}{}\textcolor{QuartoInternalColor1}{A}\textcolor{QuartoInternalColor1}{::}\textcolor{QuartoInternalColor1}{Matrix}\textcolor{QuartoInternalColor1}{\{Float64\}}\textcolor{QuartoInternalColor1}{}\textcolor{QuartoInternalColor1}{}\textcolor{QuartoInternalColor1}{)}\textcolor{QuartoInternalColor1}{}
\textcolor{QuartoInternalColor1}{}\textcolor{QuartoInternalColor1}{   @}\textcolor{QuartoInternalColor1}{ }\textcolor{QuartoInternalColor2}{Main.Notebook}\textcolor{QuartoInternalColor1}{ }\textcolor{QuartoInternalColor1}{\textasciitilde{}/nextcloud/MATH2504\_2026/lecture/quarto/}\textcolor{QuartoInternalColor1}{}\textcolor{QuartoInternalColor1}{}\textcolor{QuartoInternalColor1}{lecture-unit-5.qmd:1038}\textcolor{QuartoInternalColor1}{}\textcolor{QuartoInternalColor1}{}
\textcolor{QuartoInternalColor1}{ [2] top-level scope}
\textcolor{QuartoInternalColor1}{}\textcolor{QuartoInternalColor1}{   @}\textcolor{QuartoInternalColor1}{ }\textcolor{QuartoInternalColor1}{\textasciitilde{}/nextcloud/MATH2504\_2026/lecture/quarto/}\textcolor{QuartoInternalColor1}{}\textcolor{QuartoInternalColor1}{}\textcolor{QuartoInternalColor1}{lecture-unit-5.qmd:1048}\textcolor{QuartoInternalColor1}{}\textcolor{QuartoInternalColor1}{}
\end{Highlighting}

\begin{Shaded}
\begin{Highlighting}[]
\ImportTok{using} \BuiltInTok{LinearAlgebra}\NormalTok{, }\BuiltInTok{Random}

\BuiltInTok{Random}\NormalTok{.}\FunctionTok{seed!}\NormalTok{(}\FloatTok{0}\NormalTok{)}
\NormalTok{A }\OperatorTok{=} \FunctionTok{rand}\NormalTok{(}\FloatTok{2}\NormalTok{, }\FloatTok{2}\NormalTok{)}
\NormalTok{A\_inv }\OperatorTok{=} \FunctionTok{my\_2\_by\_2\_inv}\NormalTok{(A)}
\PreprocessorTok{@assert}\NormalTok{ A\_inv}\OperatorTok{*}\NormalTok{A }\OperatorTok{≈}\NormalTok{ I}
\end{Highlighting}
\end{Shaded}

\begin{Shaded}
\begin{Highlighting}[]
\BuiltInTok{Random}\NormalTok{.}\FunctionTok{seed!}\NormalTok{(}\FloatTok{0}\NormalTok{)}
\ControlFlowTok{for}\NormalTok{ \_ }\OperatorTok{∈} \FloatTok{1}\OperatorTok{:}\FloatTok{10} \CommentTok{\#\textbackslash{}in + [TAB]}
\NormalTok{   A }\OperatorTok{=} \FunctionTok{float}\NormalTok{.(}\FunctionTok{rand}\NormalTok{(}\FloatTok{1}\OperatorTok{:}\FloatTok{5}\NormalTok{,}\FloatTok{2}\NormalTok{,}\FloatTok{2}\NormalTok{))}
   \ControlFlowTok{try}
   \FunctionTok{my\_2\_by\_2\_inv}\NormalTok{(A)}
   \ControlFlowTok{catch} \ConstantTok{e}
   \FunctionTok{println}\NormalTok{(}\ConstantTok{e}\NormalTok{)}
   \ControlFlowTok{end}
\ControlFlowTok{end}
\end{Highlighting}
\end{Shaded}

An exception may be caught way down the call stack:

\begin{Shaded}
\begin{Highlighting}[]
\FunctionTok{f}\NormalTok{(mat) }\OperatorTok{=} \FloatTok{10}\FunctionTok{*my\_2\_by\_2\_inv}\NormalTok{(mat)}
\FunctionTok{g}\NormalTok{(mat) }\OperatorTok{=} \FunctionTok{f}\NormalTok{(mat }\OperatorTok{.+} \FloatTok{3}\NormalTok{)}
\FunctionTok{h}\NormalTok{(mat) }\OperatorTok{=} \FloatTok{2}\FunctionTok{g}\NormalTok{(mat)}
\NormalTok{A }\OperatorTok{=} \FunctionTok{ones}\NormalTok{(}\FloatTok{2}\NormalTok{,}\FloatTok{2}\NormalTok{)}
\FunctionTok{h}\NormalTok{(A)}
\end{Highlighting}
\end{Shaded}

\begin{Highlighting}
\textcolor{black}{ArgumentError: ArgumentError("matrix is singular or near singular")}
\textcolor{black}{ArgumentError: matrix is singular or near singular}
\textcolor{black}{Stacktrace:}
\textcolor{black}{ [1] }\textcolor{QuartoInternalColor1}{}\textcolor{QuartoInternalColor1}{my\_2\_by\_2\_inv}\textcolor{QuartoInternalColor1}{}\textcolor{QuartoInternalColor1}{}\textcolor{QuartoInternalColor1}{(}\textcolor{QuartoInternalColor1}{}\textcolor{QuartoInternalColor1}{A}\textcolor{QuartoInternalColor1}{::}\textcolor{QuartoInternalColor1}{Matrix}\textcolor{QuartoInternalColor1}{\{Float64\}}\textcolor{QuartoInternalColor1}{}\textcolor{QuartoInternalColor1}{}\textcolor{QuartoInternalColor1}{)}\textcolor{QuartoInternalColor1}{}
\textcolor{QuartoInternalColor1}{}\textcolor{QuartoInternalColor1}{   @}\textcolor{QuartoInternalColor1}{ }\textcolor{QuartoInternalColor2}{Main.Notebook}\textcolor{QuartoInternalColor1}{ }\textcolor{QuartoInternalColor1}{\textasciitilde{}/nextcloud/MATH2504\_2026/lecture/quarto/}\textcolor{QuartoInternalColor1}{}\textcolor{QuartoInternalColor1}{}\textcolor{QuartoInternalColor1}{lecture-unit-5.qmd:1038}\textcolor{QuartoInternalColor1}{}\textcolor{QuartoInternalColor1}{}
\textcolor{QuartoInternalColor1}{ [2] }\textcolor{QuartoInternalColor1}{}\textcolor{QuartoInternalColor1}{f}\textcolor{QuartoInternalColor1}{}
\textcolor{QuartoInternalColor1}{}\textcolor{QuartoInternalColor1}{   @}\textcolor{QuartoInternalColor1}{ }\textcolor{QuartoInternalColor1}{\textasciitilde{}/nextcloud/MATH2504\_2026/lecture/quarto/}\textcolor{QuartoInternalColor1}{}\textcolor{QuartoInternalColor1}{}\textcolor{QuartoInternalColor1}{lecture-unit-5.qmd:1075}\textcolor{QuartoInternalColor1}{}\textcolor{QuartoInternalColor1}{}\textcolor{QuartoInternalColor1}{ [inlined]}\textcolor{QuartoInternalColor1}{}
\textcolor{QuartoInternalColor1}{ [3] }\textcolor{QuartoInternalColor1}{}\textcolor{QuartoInternalColor1}{g}\textcolor{QuartoInternalColor1}{}\textcolor{QuartoInternalColor1}{}\textcolor{QuartoInternalColor1}{(}\textcolor{QuartoInternalColor1}{}\textcolor{QuartoInternalColor1}{mat}\textcolor{QuartoInternalColor1}{::}\textcolor{QuartoInternalColor1}{Matrix}\textcolor{QuartoInternalColor1}{\{Float64\}}\textcolor{QuartoInternalColor1}{}\textcolor{QuartoInternalColor1}{}\textcolor{QuartoInternalColor1}{)}\textcolor{QuartoInternalColor1}{}
\textcolor{QuartoInternalColor1}{}\textcolor{QuartoInternalColor1}{   @}\textcolor{QuartoInternalColor1}{ }\textcolor{QuartoInternalColor2}{Main.Notebook}\textcolor{QuartoInternalColor1}{ }\textcolor{QuartoInternalColor1}{\textasciitilde{}/nextcloud/MATH2504\_2026/lecture/quarto/}\textcolor{QuartoInternalColor1}{}\textcolor{QuartoInternalColor1}{}\textcolor{QuartoInternalColor1}{lecture-unit-5.qmd:1076}\textcolor{QuartoInternalColor1}{}\textcolor{QuartoInternalColor1}{}
\textcolor{QuartoInternalColor1}{ [4] }\textcolor{QuartoInternalColor1}{}\textcolor{QuartoInternalColor1}{h}\textcolor{QuartoInternalColor1}{}\textcolor{QuartoInternalColor1}{}\textcolor{QuartoInternalColor1}{(}\textcolor{QuartoInternalColor1}{}\textcolor{QuartoInternalColor1}{mat}\textcolor{QuartoInternalColor1}{::}\textcolor{QuartoInternalColor1}{Matrix}\textcolor{QuartoInternalColor1}{\{Float64\}}\textcolor{QuartoInternalColor1}{}\textcolor{QuartoInternalColor1}{}\textcolor{QuartoInternalColor1}{)}\textcolor{QuartoInternalColor1}{}
\textcolor{QuartoInternalColor1}{}\textcolor{QuartoInternalColor1}{   @}\textcolor{QuartoInternalColor1}{ }\textcolor{QuartoInternalColor2}{Main.Notebook}\textcolor{QuartoInternalColor1}{ }\textcolor{QuartoInternalColor1}{\textasciitilde{}/nextcloud/MATH2504\_2026/lecture/quarto/}\textcolor{QuartoInternalColor1}{}\textcolor{QuartoInternalColor1}{}\textcolor{QuartoInternalColor1}{lecture-unit-5.qmd:1077}\textcolor{QuartoInternalColor1}{}\textcolor{QuartoInternalColor1}{}
\textcolor{QuartoInternalColor1}{ [5] top-level scope}
\textcolor{QuartoInternalColor1}{}\textcolor{QuartoInternalColor1}{   @}\textcolor{QuartoInternalColor1}{ }\textcolor{QuartoInternalColor1}{\textasciitilde{}/nextcloud/MATH2504\_2026/lecture/quarto/}\textcolor{QuartoInternalColor1}{}\textcolor{QuartoInternalColor1}{}\textcolor{QuartoInternalColor1}{lecture-unit-5.qmd:1079}\textcolor{QuartoInternalColor1}{}\textcolor{QuartoInternalColor1}{}
\end{Highlighting}

\begin{verbatim}
ERROR: ArgumentError: matrix is singular or near singular
Stacktrace:
 [1] my_2_by_2_inv(A::Matrix{Float64})
   @ Main ~/git/mine/ProgrammingCourse-with-Julia-SimulationAnalysisAndLearningSystems/markdown/lecture-unit-4.jmd:5
 [2] f(mat::Matrix{Float64})
   @ Main ./REPL[33]:1
 [3] g(mat::Matrix{Float64})
   @ Main ./REPL[34]:1
 [4] h(mat::Matrix{Float64})
   @ Main ./REPL[35]:1
 [5] top-level scope
   @ REPL[36]:1
\end{verbatim}

\paragraph{Try, catch and finally}\label{try-catch-and-finally}

It's possible to ``catch'' and handle errors that happened using
\texttt{try}, \texttt{catch} and \texttt{finally}.

\begin{Shaded}
\begin{Highlighting}[]
\ControlFlowTok{try} 
   \FunctionTok{h}\NormalTok{(A)}
\ControlFlowTok{catch} \ConstantTok{e}
   \PreprocessorTok{@info} \StringTok{"Handling exception"} \FunctionTok{println}\NormalTok{(}\ConstantTok{e}\NormalTok{)}
\ControlFlowTok{end}
\end{Highlighting}
\end{Shaded}

\begin{verbatim}
ArgumentError("matrix is singular or near singular")
┌ Info: Handling exception
└   println(e) = nothing
\end{verbatim}

\begin{Shaded}
\begin{Highlighting}[]
\ControlFlowTok{try}
   \FloatTok{0} \OperatorTok{÷} \FloatTok{0}
\ControlFlowTok{catch} \ConstantTok{e}
   \PreprocessorTok{@show} \ConstantTok{e}
\ControlFlowTok{finally}
   \FunctionTok{println}\NormalTok{(}\StringTok{"done"}\NormalTok{)}
\ControlFlowTok{end}\NormalTok{;}
\end{Highlighting}
\end{Shaded}

\begin{verbatim}
e = DivideError()
done
\end{verbatim}

\paragraph{Resource cleanup}\label{resource-cleanup}

One place closures, and \texttt{do} syntax in particular, gets used is
when you want resources to be cleaned up automatically. Here's how one
method of \texttt{open} is defined in Julia:

\begin{Shaded}
\begin{Highlighting}[]
\KeywordTok{function} \FunctionTok{open}\NormalTok{(f, filename}\OperatorTok{::}\DataTypeTok{String}\NormalTok{)}
\NormalTok{   io }\OperatorTok{=} \FunctionTok{open}\NormalTok{(filename)}
   \ControlFlowTok{try}
   \FunctionTok{f}\NormalTok{(io)}
   \ControlFlowTok{finally}
   \FunctionTok{close}\NormalTok{(io)}
   \ControlFlowTok{end}
\KeywordTok{end}
\end{Highlighting}
\end{Shaded}

You can call the \texttt{open} function with an anonymous function like
so:

\begin{Shaded}
\begin{Highlighting}[]
\CommentTok{\# Read the first 128 bytes from the file}
\FunctionTok{open}\NormalTok{(io }\OperatorTok{{-}\textgreater{}} \FunctionTok{read}\NormalTok{(io, }\FloatTok{128}\NormalTok{), filename)}
\end{Highlighting}
\end{Shaded}

This way even when \texttt{read} might error out (e.g.~too few bytes),
the file is \emph{always} closed in the \texttt{finally} block above.
Using this pattern avoids mistakes and resource leakage.

When you want to use this pattern it is common to use \texttt{do} syntax
like so:

\begin{Shaded}
\begin{Highlighting}[]
\FunctionTok{open}\NormalTok{(filename) }\ControlFlowTok{do}\NormalTok{ (io)}
   \CommentTok{\# Read the first 128 bytes from the file}
   \ControlFlowTok{return} \FunctionTok{read}\NormalTok{(io, }\FloatTok{128}\NormalTok{) }\CommentTok{\# might throw an exception}
\ControlFlowTok{end}
\end{Highlighting}
\end{Shaded}

The \texttt{do} syntax is \emph{yet another} way of making a function.
It creates a new function and injects it into the first argument.

\subsubsection{Variable scope}\label{variable-scope}

See
\href{https://docs.julialang.org/en/v1/manual/variables-and-scoping/}{variables
and scoping in Julia docs}.

Julia actually has two separate scopes --- global or ``toplevel'' scope,
and local scope.

Local scope is the simplest, and we'll look at that first. By default
bindings inside a function body create \emph{new} variables that take
precedence (or ``shadows'') any variables of the same name outside the
function body.

\begin{Shaded}
\begin{Highlighting}[]
\NormalTok{x }\OperatorTok{=} \FloatTok{0}

\KeywordTok{function} \FunctionTok{f}\NormalTok{()}
\NormalTok{   x }\OperatorTok{=} \FloatTok{1} \CommentTok{\# create a new local variable}
   \ControlFlowTok{return} \FloatTok{2} \OperatorTok{*}\NormalTok{ x}
\KeywordTok{end}

\FunctionTok{f}\NormalTok{()}
\end{Highlighting}
\end{Shaded}

\begin{verbatim}
2
\end{verbatim}

You can use the \texttt{local} keyword to create new local bindings. The
above is 100\% equivalent to:

\begin{Shaded}
\begin{Highlighting}[]
\KeywordTok{function} \FunctionTok{f}\NormalTok{()}
   \KeywordTok{local}\NormalTok{ x }\OperatorTok{=} \FloatTok{1} \CommentTok{\# create a new local variable}
   \ControlFlowTok{return} \FloatTok{2} \OperatorTok{*}\NormalTok{ x}
\KeywordTok{end}
\end{Highlighting}
\end{Shaded}

If you want to modify the global variable, there is a \texttt{global}
keyword for that:

\begin{Shaded}
\begin{Highlighting}[]
\NormalTok{x }\OperatorTok{=} \FloatTok{0}

\KeywordTok{function} \FunctionTok{f}\NormalTok{()}
   \KeywordTok{global}\NormalTok{ x }\OperatorTok{=} \FloatTok{1} \CommentTok{\# modify the global variable}
   \ControlFlowTok{return} \FloatTok{2} \OperatorTok{*}\NormalTok{ x}
\KeywordTok{end}

\PreprocessorTok{@show}\NormalTok{ x}
\PreprocessorTok{@show} \FunctionTok{f}\NormalTok{()}
\PreprocessorTok{@show}\NormalTok{ x;}
\end{Highlighting}
\end{Shaded}

\begin{verbatim}
x = 0
f() = 2
x = 1
\end{verbatim}

Inside other code blocks such as \texttt{if}, \texttt{while} and
\texttt{for}, the default is to overwrite existing local variables if
they exist.

\emph{Q: What is the difference between these two functions?}

\begin{Shaded}
\begin{Highlighting}[]
\KeywordTok{function} \FunctionTok{f1}\NormalTok{()}
\NormalTok{   x }\OperatorTok{=} \FloatTok{0}
   \ControlFlowTok{for}\NormalTok{ i }\OperatorTok{=} \FloatTok{1}\OperatorTok{:}\FloatTok{10}
\NormalTok{   x }\OperatorTok{=}\NormalTok{ i}
   \ControlFlowTok{end}
   \ControlFlowTok{return}\NormalTok{ x}
\KeywordTok{end}

\KeywordTok{function} \FunctionTok{f2}\NormalTok{()}
\NormalTok{   x }\OperatorTok{=} \FloatTok{0}
   \ControlFlowTok{for}\NormalTok{ i }\OperatorTok{=} \FloatTok{1}\OperatorTok{:}\FloatTok{10}
   \KeywordTok{local}\NormalTok{ x }\OperatorTok{=}\NormalTok{ i}
   \ControlFlowTok{end}
   \ControlFlowTok{return}\NormalTok{ x}
\KeywordTok{end}
\end{Highlighting}
\end{Shaded}

\begin{verbatim}
f2 (generic function with 2 methods)
\end{verbatim}

By default variables created \emph{inside} a loop no longer exist after
the loop:

\begin{Shaded}
\begin{Highlighting}[]
\KeywordTok{function} \FunctionTok{f}\NormalTok{()}
   \ControlFlowTok{for}\NormalTok{ i }\OperatorTok{=} \FloatTok{1}\OperatorTok{:}\FloatTok{3}
\NormalTok{   x }\OperatorTok{=}\NormalTok{ i}
   \ControlFlowTok{end}
   \ControlFlowTok{return}\NormalTok{ x}
\KeywordTok{end}
\FunctionTok{f}\NormalTok{()}
\end{Highlighting}
\end{Shaded}

\begin{verbatim}
1
\end{verbatim}

To fix that you can define that variable in advance using
\texttt{local}:

\begin{Shaded}
\begin{Highlighting}[]
\KeywordTok{function} \FunctionTok{f}\NormalTok{()}
   \KeywordTok{local}\NormalTok{ x}
   \ControlFlowTok{for}\NormalTok{ i }\OperatorTok{=} \FloatTok{1}\OperatorTok{:}\FloatTok{3}
\NormalTok{   x }\OperatorTok{=}\NormalTok{ i}
   \ControlFlowTok{end}
   \ControlFlowTok{return}\NormalTok{ x}
\KeywordTok{end}
\FunctionTok{f}\NormalTok{()}
\end{Highlighting}
\end{Shaded}

\begin{verbatim}
3
\end{verbatim}

Occasionally you might want to use the value of what is being iterated,
where you can use \texttt{outer}:

\begin{Shaded}
\begin{Highlighting}[]
\KeywordTok{function} \FunctionTok{f}\NormalTok{()}
\NormalTok{   x }\OperatorTok{=} \FloatTok{0}
   \ControlFlowTok{for}\NormalTok{ outer x }\OperatorTok{=} \FloatTok{1}\OperatorTok{:}\FloatTok{3}
   \CommentTok{\# empty}
   \ControlFlowTok{end}
   \ControlFlowTok{return}\NormalTok{ x}
\KeywordTok{end}
\FunctionTok{f}\NormalTok{()}
\end{Highlighting}
\end{Shaded}

\begin{verbatim}
3
\end{verbatim}

That can be useful in combination to \texttt{break}, where you can see
where the loop terminated:

\begin{Shaded}
\begin{Highlighting}[]
\KeywordTok{function} \FunctionTok{f}\NormalTok{()}
\NormalTok{   x }\OperatorTok{=} \FloatTok{0}
   \ControlFlowTok{for}\NormalTok{ outer x }\OperatorTok{=} \FloatTok{1}\OperatorTok{:}\FloatTok{3}
   \ControlFlowTok{if}\NormalTok{ x }\OperatorTok{\textgreater{}} \FloatTok{1.5}
   \ControlFlowTok{break}
   \ControlFlowTok{end}
   \ControlFlowTok{end}
   \ControlFlowTok{return}\NormalTok{ x}
\KeywordTok{end}
\FunctionTok{f}\NormalTok{()}
\end{Highlighting}
\end{Shaded}

\begin{verbatim}
2
\end{verbatim}

Working with variables in different scopes can a little confusing and
error prone. My advice:

\begin{itemize}
\tightlist
\item
  Don't use non-const global variables unless there is no alternative.
  Prefer passing variables to functions, including using closures to
  capture state (see also ScopedValue in Julia 1.11).
\item
  Use distinct names for variables as much as possible
\item
  The larger the scope of a variable, the longer and more descriptive
  its name should be
\end{itemize}

Remember you want your programs to be \emph{simple} and \emph{easy to
read}. Variable names are cheap (though naming things well is hard!)

\subsection{More on mutation}\label{more-on-mutation}

As we saw earlier, a type can be \textbf{mutable} or not
(\textbf{immutable}). Values/objects of the immutable types are
typically stored on the stack. Values/objects of mutable types are
typically allocated and stored on the heap.

\begin{Shaded}
\begin{Highlighting}[]
\PreprocessorTok{@show} \FunctionTok{ismutable}\NormalTok{(}\FloatTok{7}\NormalTok{)}
\PreprocessorTok{@show} \FunctionTok{ismutable}\NormalTok{([}\FloatTok{7}\NormalTok{]);}
\end{Highlighting}
\end{Shaded}

\begin{verbatim}
ismutable(7) = false
ismutable([7]) = true
\end{verbatim}

When you pass a mutable object to a function, the function can change
the data contained within.

\begin{Shaded}
\begin{Highlighting}[]
\KeywordTok{function} \FunctionTok{mutating\_function}\NormalTok{(z}\OperatorTok{::}\DataTypeTok{Array\{Int\}}\NormalTok{) }
\NormalTok{   z[}\FloatTok{1}\NormalTok{] }\OperatorTok{=} \FloatTok{0}
\KeywordTok{end}

\NormalTok{x }\OperatorTok{=}\NormalTok{ [}\FloatTok{1}\NormalTok{]}
\FunctionTok{println}\NormalTok{(}\StringTok{"Before call: "}\NormalTok{, x)}
\FunctionTok{mutating\_function}\NormalTok{(x)}
\FunctionTok{println}\NormalTok{(}\StringTok{"After call: "}\NormalTok{, x)}
\end{Highlighting}
\end{Shaded}

\begin{verbatim}
Before call: [1]
After call: [0]
\end{verbatim}

Programmers tend to need to keep track of which objects are being
mutated and when. You'll see functions that mutate things often end in
\texttt{!}, e.g.~\texttt{push!}, \texttt{pop!}, \texttt{setindex!}, etc.
This is just a convention -- we could have called the above
\texttt{mutating\_function!}.

Note that \texttt{collection{[}index{]}} is shorthand for
\texttt{getindex(collection,\ index)} and
\texttt{collection{[}index{]}\ =\ value} is shorthand for
\texttt{setindex!(collection,\ index,\ value)}.

For immutable values, there is no way to modify their contents. They are
never modified by functions.

\paragraph{Variable bindings}\label{variable-bindings}

Variable bindings are not modified inside functions. The function below
introduces a \emph{new} binding \texttt{z} every time the function is
called, and \texttt{x} is unaffected.

\begin{Shaded}
\begin{Highlighting}[]
\KeywordTok{function} \FunctionTok{rebinding\_function}\NormalTok{(z}\OperatorTok{::}\DataTypeTok{Array\{Int\}}\NormalTok{) }
\NormalTok{   z }\OperatorTok{=}\NormalTok{ [}\FloatTok{0}\NormalTok{]}
\KeywordTok{end}

\NormalTok{x }\OperatorTok{=}\NormalTok{ [}\FloatTok{1}\NormalTok{]}
\FunctionTok{println}\NormalTok{(}\StringTok{"Before call: "}\NormalTok{, x)}
\FunctionTok{mutating\_function}\NormalTok{(x)}
\FunctionTok{println}\NormalTok{(}\StringTok{"After call: "}\NormalTok{, x)}
\end{Highlighting}
\end{Shaded}

\begin{verbatim}
Before call: [1]
After call: [0]
\end{verbatim}

So, the rule is that you can mutate the \emph{inside} of a mutable value
(\texttt{Array} or \texttt{mutable\ struct}) and \emph{other} variables
and/or data structures with references to that mutable value will be
able to see the effect of the mutation.

However variables are \emph{always} bound to \emph{new} values with
\texttt{=} (e.g.~\texttt{z\ =\ 0}) and the usual lexical scoping rules
apply.

To help with preserving data that might get modified, we often make a
copy with \texttt{copy} and \texttt{deepcopy}:

\begin{Shaded}
\begin{Highlighting}[]
\FunctionTok{println}\NormalTok{(}\StringTok{"Immutable:"}\NormalTok{)}
\NormalTok{a }\OperatorTok{=} \FloatTok{10}
\NormalTok{b }\OperatorTok{=}\NormalTok{ a}
\NormalTok{b }\OperatorTok{=} \FloatTok{20}
\PreprocessorTok{@show}\NormalTok{ a;}
\end{Highlighting}
\end{Shaded}

\begin{verbatim}
Immutable:
a = 10
\end{verbatim}

\begin{Shaded}
\begin{Highlighting}[]
\FunctionTok{println}\NormalTok{(}\StringTok{"}\SpecialCharTok{\textbackslash{}n}\StringTok{No copy:"}\NormalTok{)}
\NormalTok{a }\OperatorTok{=}\NormalTok{ [}\FloatTok{10}\NormalTok{]}
\NormalTok{b }\OperatorTok{=}\NormalTok{ a}
\NormalTok{b[}\FloatTok{1}\NormalTok{] }\OperatorTok{=} \FloatTok{20}
\PreprocessorTok{@show}\NormalTok{ a;}
\end{Highlighting}
\end{Shaded}

\begin{verbatim}

No copy:
a = [20]
\end{verbatim}

\begin{Shaded}
\begin{Highlighting}[]
\FunctionTok{println}\NormalTok{(}\StringTok{"}\SpecialCharTok{\textbackslash{}n}\StringTok{Copy:"}\NormalTok{)}
\NormalTok{a }\OperatorTok{=}\NormalTok{ [}\FloatTok{10}\NormalTok{]}
\NormalTok{b }\OperatorTok{=} \FunctionTok{copy}\NormalTok{(a)}
\NormalTok{b[}\FloatTok{1}\NormalTok{] }\OperatorTok{=} \FloatTok{20}
\PreprocessorTok{@show}\NormalTok{ a;}
\end{Highlighting}
\end{Shaded}

\begin{verbatim}

Copy:
a = [10]
\end{verbatim}

\begin{Shaded}
\begin{Highlighting}[]
\FunctionTok{println}\NormalTok{(}\StringTok{"}\SpecialCharTok{\textbackslash{}n}\StringTok{Shallow copy:"}\NormalTok{)}
\NormalTok{a }\OperatorTok{=}\NormalTok{ [[}\FloatTok{10}\NormalTok{]]}
\NormalTok{b }\OperatorTok{=} \FunctionTok{copy}\NormalTok{(a)}
\NormalTok{b[}\FloatTok{1}\NormalTok{][}\FloatTok{1}\NormalTok{] }\OperatorTok{=} \FloatTok{20}
\PreprocessorTok{@show}\NormalTok{ a;}
\end{Highlighting}
\end{Shaded}

\begin{verbatim}

Shallow copy:
a = [[20]]
\end{verbatim}

\begin{Shaded}
\begin{Highlighting}[]
\FunctionTok{println}\NormalTok{(}\StringTok{"}\SpecialCharTok{\textbackslash{}n}\StringTok{Deep copy:"}\NormalTok{)}
\NormalTok{a }\OperatorTok{=}\NormalTok{ [[}\FloatTok{10}\NormalTok{]]}
\NormalTok{b }\OperatorTok{=} \FunctionTok{deepcopy}\NormalTok{(a)}
\NormalTok{b[}\FloatTok{1}\NormalTok{][}\FloatTok{1}\NormalTok{] }\OperatorTok{=} \FloatTok{20}
\PreprocessorTok{@show}\NormalTok{ a;}
\end{Highlighting}
\end{Shaded}

\begin{verbatim}

Deep copy:
a = [[10]]
\end{verbatim}

\subsection{Abstract type tree}\label{abstract-type-tree}

We saw the \texttt{Union} abstract type earlier.

Julia also has an abstract type hierarchy (a tree). At the top of the
tree is the type \texttt{Any}, which encompasses every possible value in
Julia. All types have a supertype (the supertype of \texttt{Any} is
\texttt{Any}). Types that are not leaves of the tree have subtypes. Some
types are \textbf{abstract} while others are \textbf{concrete}. One
particularly distinctive feature of Julia's type system is that concrete
types may not subtype each other: all concrete types are final and may
only have abstract types as their supertypes.

\begin{Shaded}
\begin{Highlighting}[]
\NormalTok{x }\OperatorTok{=} \FloatTok{2.3}
\PreprocessorTok{@show} \FunctionTok{typeof}\NormalTok{(x)}
\PreprocessorTok{@show} \FunctionTok{supertype}\NormalTok{(}\DataTypeTok{Float64}\NormalTok{)}
\PreprocessorTok{@show} \FunctionTok{supertype}\NormalTok{(}\DataTypeTok{AbstractFloat}\NormalTok{)}
\PreprocessorTok{@show} \FunctionTok{supertype}\NormalTok{(}\DataTypeTok{Real}\NormalTok{)}
\PreprocessorTok{@show} \FunctionTok{supertype}\NormalTok{(}\DataTypeTok{Number}\NormalTok{)}
\PreprocessorTok{@show} \FunctionTok{supertype}\NormalTok{(}\DataTypeTok{Any}\NormalTok{);}
\end{Highlighting}
\end{Shaded}

\begin{verbatim}
typeof(x) = Float64
supertype(Float64) = AbstractFloat
supertype(AbstractFloat) = Real
supertype(Real) = Number
supertype(Number) = Any
supertype(Any) = Any
\end{verbatim}

There is an \textbf{is a} relationship:

\begin{Shaded}
\begin{Highlighting}[]
\FunctionTok{isa}\NormalTok{(}\FloatTok{2.3}\NormalTok{, }\DataTypeTok{Number}\NormalTok{)}
\end{Highlighting}
\end{Shaded}

\begin{verbatim}
true
\end{verbatim}

\begin{Shaded}
\begin{Highlighting}[]
\FunctionTok{isa}\NormalTok{(}\FloatTok{2.3}\NormalTok{, }\DataTypeTok{String}\NormalTok{)}
\end{Highlighting}
\end{Shaded}

\begin{verbatim}
false
\end{verbatim}

\begin{Shaded}
\begin{Highlighting}[]
\FloatTok{2.3}\NormalTok{ isa }\DataTypeTok{Float64}
\end{Highlighting}
\end{Shaded}

\begin{verbatim}
true
\end{verbatim}

Note that \texttt{x\ isa\ T} is the same as
\texttt{typeof(x)\ \textless{}:\ T}, where we say \texttt{\textless{}:}
as ``is a subtype of''.

\begin{Shaded}
\begin{Highlighting}[]
\PreprocessorTok{@show} \DataTypeTok{Float64} \OperatorTok{\textless{}:}\DataTypeTok{ Number}
\PreprocessorTok{@show} \DataTypeTok{String} \OperatorTok{\textless{}:}\DataTypeTok{ Number}\NormalTok{;}
\end{Highlighting}
\end{Shaded}

\begin{verbatim}
Float64 <: Number = true
String <: Number = false
\end{verbatim}

We can ask whether a given type is abstract or concrete.

\begin{Shaded}
\begin{Highlighting}[]
\PreprocessorTok{@show} \FunctionTok{isabstracttype}\NormalTok{(}\DataTypeTok{Float64}\NormalTok{)}
\PreprocessorTok{@show} \FunctionTok{isconcretetype}\NormalTok{(}\DataTypeTok{Float64}\NormalTok{);}
\end{Highlighting}
\end{Shaded}

\begin{verbatim}
isabstracttype(Float64) = false
isconcretetype(Float64) = true
\end{verbatim}

\begin{Shaded}
\begin{Highlighting}[]
\PreprocessorTok{@show} \FunctionTok{isabstracttype}\NormalTok{(}\DataTypeTok{Real}\NormalTok{)}
\PreprocessorTok{@show} \FunctionTok{isconcretetype}\NormalTok{(}\DataTypeTok{Real}\NormalTok{);}
\end{Highlighting}
\end{Shaded}

\begin{verbatim}
isabstracttype(Real) = true
isconcretetype(Real) = false
\end{verbatim}

Structs with undefined type parameters are not concrete:

\begin{Shaded}
\begin{Highlighting}[]
\PreprocessorTok{@show} \FunctionTok{isconcretetype}\NormalTok{(}\DataTypeTok{Complex}\NormalTok{);}
\end{Highlighting}
\end{Shaded}

\begin{verbatim}
isconcretetype(Complex) = false
\end{verbatim}

Once we provide the type parameters we do get a concrete type:

\begin{Shaded}
\begin{Highlighting}[]
\PreprocessorTok{@show} \FunctionTok{isconcretetype}\NormalTok{(}\DataTypeTok{Complex}\NormalTok{\{}\DataTypeTok{Float64}\NormalTok{\});}
\end{Highlighting}
\end{Shaded}

\begin{verbatim}
isconcretetype(Complex{Float64}) = true
\end{verbatim}

As mentioned, Julia has a type tree. Let's walk down from
\texttt{Number}:

\begin{Shaded}
\begin{Highlighting}[]
\ImportTok{using} \BuiltInTok{InteractiveUtils}\NormalTok{: subtypes}

\KeywordTok{function} \FunctionTok{type\_and\_children}\NormalTok{(}\KeywordTok{type}\NormalTok{, prefix }\OperatorTok{=} \StringTok{""}\NormalTok{, child\_prefix }\OperatorTok{=} \StringTok{""}\NormalTok{)}
   \ControlFlowTok{if} \FunctionTok{isconcretetype}\NormalTok{(}\KeywordTok{type}\NormalTok{)}
   \PreprocessorTok{@assert} \FunctionTok{isempty}\NormalTok{(}\FunctionTok{subtypes}\NormalTok{(}\KeywordTok{type}\NormalTok{))}

   \FunctionTok{println}\NormalTok{(prefix, }\KeywordTok{type}\NormalTok{, }\StringTok{": concrete"}\NormalTok{)}
   \ControlFlowTok{else}
   \FunctionTok{println}\NormalTok{(prefix, }\KeywordTok{type}\NormalTok{, }\FunctionTok{isabstracttype}\NormalTok{(}\KeywordTok{type}\NormalTok{) ? }\StringTok{": abstract"} \OperatorTok{:} \StringTok{": parameterised"}\NormalTok{)}

\NormalTok{   children }\OperatorTok{=} \FunctionTok{subtypes}\NormalTok{(}\KeywordTok{type}\NormalTok{)}
   \ControlFlowTok{for}\NormalTok{ (i, c) }\KeywordTok{in} \FunctionTok{enumerate}\NormalTok{(children)}
   \ControlFlowTok{if}\NormalTok{ i }\OperatorTok{==} \FunctionTok{length}\NormalTok{(children)}
   \FunctionTok{type\_and\_children}\NormalTok{(c, }\StringTok{"}\SpecialCharTok{$}\NormalTok{(child\_prefix)}\StringTok{ └─╴"}\NormalTok{, }\StringTok{"}\SpecialCharTok{$}\NormalTok{(child\_prefix)}\StringTok{    "}\NormalTok{)}
   \ControlFlowTok{else}
   \FunctionTok{type\_and\_children}\NormalTok{(c, }\StringTok{"}\SpecialCharTok{$}\NormalTok{(child\_prefix)}\StringTok{ ├─╴"}\NormalTok{, }\StringTok{"}\SpecialCharTok{$}\NormalTok{(child\_prefix)}\StringTok{ │  "}\NormalTok{)}
   \ControlFlowTok{end} 
   \ControlFlowTok{end}
   \ControlFlowTok{end}
\KeywordTok{end}

\FunctionTok{type\_and\_children}\NormalTok{(}\DataTypeTok{Number}\NormalTok{)}
\end{Highlighting}
\end{Shaded}

\begin{verbatim}
Number: abstract
 ├─╴Base.MultiplicativeInverses.MultiplicativeInverse: abstract
 │   ├─╴Base.MultiplicativeInverses.SignedMultiplicativeInverse: parameterised
 │   └─╴Base.MultiplicativeInverses.UnsignedMultiplicativeInverse: parameterised
 ├─╴Base.Complex: parameterised
 ├─╴Real: abstract
 │   ├─╴AbstractFloat: abstract
 │   │   ├─╴BigFloat: concrete
 │   │   ├─╴Core.BFloat16: concrete
 │   │   ├─╴Float16: concrete
 │   │   ├─╴Float32: concrete
 │   │   └─╴Float64: concrete
 │   ├─╴AbstractIrrational: abstract
 │   │   ├─╴Irrational: parameterised
 │   │   └─╴IrrationalConstants.IrrationalConstant: abstract
 │   │       ├─╴IrrationalConstants.Fourinvπ: concrete
 │   │       ├─╴IrrationalConstants.Fourπ: concrete
 │   │       ├─╴IrrationalConstants.Halfπ: concrete
 │   │       ├─╴IrrationalConstants.Inv2π: concrete
 │   │       ├─╴IrrationalConstants.Inv4π: concrete
 │   │       ├─╴IrrationalConstants.Invsqrt2: concrete
 │   │       ├─╴IrrationalConstants.Invsqrt2π: concrete
 │   │       ├─╴IrrationalConstants.Invsqrtπ: concrete
 │   │       ├─╴IrrationalConstants.Invπ: concrete
 │   │       ├─╴IrrationalConstants.Log2π: concrete
 │   │       ├─╴IrrationalConstants.Log4π: concrete
 │   │       ├─╴IrrationalConstants.Loghalf: concrete
 │   │       ├─╴IrrationalConstants.Logten: concrete
 │   │       ├─╴IrrationalConstants.Logtwo: concrete
 │   │       ├─╴IrrationalConstants.Logπ: concrete
 │   │       ├─╴IrrationalConstants.Quartπ: concrete
 │   │       ├─╴IrrationalConstants.Sqrt2: concrete
 │   │       ├─╴IrrationalConstants.Sqrt2π: concrete
 │   │       ├─╴IrrationalConstants.Sqrt3: concrete
 │   │       ├─╴IrrationalConstants.Sqrt4π: concrete
 │   │       ├─╴IrrationalConstants.Sqrthalfπ: concrete
 │   │       ├─╴IrrationalConstants.Sqrtπ: concrete
 │   │       ├─╴IrrationalConstants.Twoinvπ: concrete
 │   │       └─╴IrrationalConstants.Twoπ: concrete
 │   ├─╴Integer: abstract
 │   │   ├─╴Bool: concrete
 │   │   ├─╴Signed: abstract
 │   │   │   ├─╴BigInt: concrete
 │   │   │   ├─╴Int128: concrete
 │   │   │   ├─╴Int16: concrete
 │   │   │   ├─╴Int32: concrete
 │   │   │   ├─╴Int64: concrete
 │   │   │   └─╴Int8: concrete
 │   │   └─╴Unsigned: abstract
 │   │       ├─╴UInt128: concrete
 │   │       ├─╴UInt16: concrete
 │   │       ├─╴UInt32: concrete
 │   │       ├─╴UInt64: concrete
 │   │       └─╴UInt8: concrete
 │   ├─╴Rational: parameterised
 │   ├─╴StatsBase.PValue: concrete
 │   └─╴StatsBase.TestStat: concrete
 ├─╴Unitful.AbstractQuantity: abstract
 │   └─╴Unitful.Quantity: parameterised
 └─╴Unitful.LogScaled: abstract
     ├─╴Unitful.Gain{L} where L<:Unitful.LogInfo: parameterised
     └─╴Unitful.Level{L} where L<:Unitful.LogInfo: parameterised
\end{verbatim}

In Julia, you can define abstract types with the \texttt{abstract\ type}
keywords:

\begin{Shaded}
\begin{Highlighting}[]
\KeywordTok{abstract type} \DataTypeTok{Number}
\KeywordTok{end}

\KeywordTok{abstract type} \DataTypeTok{Real} \OperatorTok{\textless{}:}\DataTypeTok{ Number}
\KeywordTok{end}

\KeywordTok{abstract type} \DataTypeTok{AbstractFloat} \OperatorTok{\textless{}:}\DataTypeTok{ Real}
\KeywordTok{end}

\KeywordTok{primitive type} \DataTypeTok{Float64} \OperatorTok{\textless{}:}\DataTypeTok{ AbstractFloat}
   \FloatTok{64}
\KeywordTok{end}
\end{Highlighting}
\end{Shaded}

\subsubsection{Parameterized / generic
types}\label{parameterized-generic-types}

We've actually now seen all three types of abstract types in Julia --
the \texttt{abstract\ type}s that make up the type tree, the
\texttt{Union} type, and the parameterized types (abstract
\texttt{Complex} vs concrete \texttt{Complex\{Float64\}}).

Actually \texttt{Complex} is a shorthand. It's full type is written like
this:

\begin{Shaded}
\begin{Highlighting}[]
\DataTypeTok{Complex}\NormalTok{\{T\} }\KeywordTok{where}\NormalTok{ T }\OperatorTok{\textless{}:}\DataTypeTok{ Real}
\end{Highlighting}
\end{Shaded}

This object is of type \texttt{UnionAll}:

\begin{Shaded}
\begin{Highlighting}[]
\FunctionTok{typeof}\NormalTok{(}\DataTypeTok{Complex}\NormalTok{)}
\end{Highlighting}
\end{Shaded}

\begin{verbatim}
UnionAll
\end{verbatim}

You can read this like ``The abstract type which is the union of the
concrete types \texttt{Complex\{T\}} for all possible
\texttt{T\ \textless{}:\ Real}'' -- hence the shorthand
\texttt{UnionAll}. Parameterised types can have bounds, like the
components of complex numbers being real numbers.

\begin{Shaded}
\begin{Highlighting}[]
\PreprocessorTok{@show} \DataTypeTok{Complex}\NormalTok{\{}\DataTypeTok{Float64}\NormalTok{\} }\OperatorTok{\textless{}:}\DataTypeTok{ Complex}
\PreprocessorTok{@show} \FunctionTok{isconcretetype}\NormalTok{(}\DataTypeTok{Complex}\NormalTok{)}
\PreprocessorTok{@show} \FunctionTok{isconcretetype}\NormalTok{(}\DataTypeTok{Complex}\NormalTok{\{}\DataTypeTok{Float64}\NormalTok{\});}
\end{Highlighting}
\end{Shaded}

\begin{verbatim}
Complex{Float64} <: Complex = true
isconcretetype(Complex) = false
isconcretetype(Complex{Float64}) = true
\end{verbatim}

Sometimes it's useful to \textbf{use values which aren't types as type
parameters (``Value types'')}. One standard example is the
\textbf{\texttt{N} in \texttt{Array\{T,N\}}} which is just an integer
signifying the dimension of the array. Doing this helps the compiler
know some parts of the ``problem size'' at compile time based on the
type and can lead to much more efficient code or data storage (see
example \texttt{RegularPolygon} below).

Julia is pretty capable at figuring out the subtype (or subset)
relationships:

\begin{Shaded}
\begin{Highlighting}[]
\NormalTok{(}\DataTypeTok{Complex}\NormalTok{\{T\} }\KeywordTok{where}\NormalTok{ T }\OperatorTok{\textless{}:}\DataTypeTok{ AbstractFloat}\NormalTok{) }\OperatorTok{\textless{}:}\DataTypeTok{ Complex}
\end{Highlighting}
\end{Shaded}

\begin{verbatim}
true
\end{verbatim}

which follow from \texttt{AbstractFloat\ \textless{}:\ Real}.

You've seen other \texttt{UnionAll} types like \texttt{Vector},
\texttt{Matrix}, \texttt{Set}, \texttt{Dict}, etc.

Don't worry about this seeming complex -- consider this background
material!

\subsubsection{Union types}\label{union-types}

We saw earlier that Julia has a third abstract type called
\texttt{Union} which let's you reason about a finite set of (abstract or
concrete) types.

\begin{Shaded}
\begin{Highlighting}[]
\FloatTok{42}\OperatorTok{::}\DataTypeTok{Union\{Int, Float64\}}
\end{Highlighting}
\end{Shaded}

\begin{verbatim}
42
\end{verbatim}

\begin{Shaded}
\begin{Highlighting}[]
\FloatTok{3.14}\OperatorTok{::}\DataTypeTok{Union\{Int, Float64\}}
\end{Highlighting}
\end{Shaded}

\begin{verbatim}
3.14
\end{verbatim}

\begin{Shaded}
\begin{Highlighting}[]
\StringTok{"abc"}\OperatorTok{::}\DataTypeTok{Union\{Int, Float64\}}
\end{Highlighting}
\end{Shaded}

\begin{Highlighting}
\textcolor{black}{TypeError: TypeError(:typeassert, "", Union\{Float64, Int64\}, "abc")}
\textcolor{black}{TypeError: in typeassert, expected Union\{Float64, Int64\}, got a value of type String}
\textcolor{black}{Stacktrace:}
\textcolor{black}{ [1] top-level scope}
\textcolor{black}{}\textcolor{QuartoInternalColor1}{   @}\textcolor{QuartoInternalColor1}{ }\textcolor{QuartoInternalColor1}{\textasciitilde{}/nextcloud/MATH2504\_2026/lecture/quarto/}\textcolor{QuartoInternalColor1}{}\textcolor{QuartoInternalColor1}{}\textcolor{QuartoInternalColor1}{lecture-unit-5.qmd:1534}\textcolor{QuartoInternalColor1}{}\textcolor{QuartoInternalColor1}{}
\end{Highlighting}

\texttt{Union} can handle an arbitrary number of types,
\texttt{Union\{T1,\ T2,\ T3,\ ...\}}.

As a special case \texttt{Union\{T\}} is just the same as \texttt{T}. We
also have \texttt{Union\{T,\ T\}\ ==\ T}, etc.

The union of no types at all, \texttt{Union\{\}}, is a special builtin
type which is the opposite of \texttt{Any}. No value can exist with type
\texttt{Union\{\}}! Sometimes \texttt{Any} is called the ``top'' type
and \texttt{Union\{\}} is called the ``bottom'' type. It's used
internally by the compiler to rule out impossible situations, but it's
not something for you to worry about.

You've now seen \emph{every possible concrete type} and \emph{every
possible abstract type} in all of Julia. You've also looked a functions,
methods and multiple dispatch.

Don't worry if that was whirlwind. You don't need to memorise every
detail that I've covered!

Now we will practice using these tools to build larger programs, so you
are ready for Projects 1, 2 and 3.

\subsection{Crafting programs in
Julia}\label{crafting-programs-in-julia}

\subsubsection{Making types for your programs - dimensions and
units}\label{making-types-for-your-programs---dimensions-and-units}

You can define your own types. Any ``serious'' programming task would
almost always merit that you do that.

\paragraph{Dimensions and units}\label{dimensions-and-units}

One thing types let you do is ensure you do things ``safely'' without
making mistakes.

When solving physics problems its always useful to perform dimensional
analysis to ensure your answer makes sense, and to make sure you are
using a consistent set of units in your numeric computations.

\begin{itemize}
\tightlist
\item
  \textbf{Dimensions} - length, time, mass, area (length squared), etc
\item
  \textbf{Units} - metres, kilometres, centimetres, feet, inches, miles,
  etc
\end{itemize}

\paragraph{\texorpdfstring{Our first unit -
\texttt{Metre}}{Our first unit - Metre}}\label{our-first-unit---metre}

Julia lets us encode dimensions and units. There is a nice package
called
\href{https://painterqubits.github.io/Unitful.jl/stable/}{Unitful.jl}
that you can download and use. Lets see how a package like this works.

We can start simple with a single type:

\begin{Shaded}
\begin{Highlighting}[]
\KeywordTok{struct}\NormalTok{ Metre\{T }\OperatorTok{\textless{}:}\DataTypeTok{ Real}\NormalTok{\}}
\NormalTok{   value}\OperatorTok{::}\DataTypeTok{T}
\KeywordTok{end}
\end{Highlighting}
\end{Shaded}

That's all very good, but we can't actually \emph{do} anything with it.
When you introduce new data types you'll generally need some functions
or methods to interact with the data. Lets define some basic algebra:

\begin{Shaded}
\begin{Highlighting}[]
\BuiltInTok{Base}\NormalTok{.}\OperatorTok{:}\NormalTok{(}\OperatorTok{+}\NormalTok{)(x}\OperatorTok{::}\DataTypeTok{Metre}\NormalTok{, y}\OperatorTok{::}\DataTypeTok{Metre}\NormalTok{) }\OperatorTok{=} \FunctionTok{Metre}\NormalTok{(x.value }\OperatorTok{+}\NormalTok{ y.value)}
\BuiltInTok{Base}\NormalTok{.}\OperatorTok{:}\NormalTok{(}\OperatorTok{{-}}\NormalTok{)(x}\OperatorTok{::}\DataTypeTok{Metre}\NormalTok{, y}\OperatorTok{::}\DataTypeTok{Metre}\NormalTok{) }\OperatorTok{=} \FunctionTok{Metre}\NormalTok{(x.value }\OperatorTok{{-}}\NormalTok{ y.value)}
\BuiltInTok{Base}\NormalTok{.}\OperatorTok{:}\NormalTok{(}\OperatorTok{{-}}\NormalTok{)(x}\OperatorTok{::}\DataTypeTok{Metre}\NormalTok{) }\OperatorTok{=} \FunctionTok{Metre}\NormalTok{(}\OperatorTok{{-}}\NormalTok{x.value)}
\BuiltInTok{Base}\NormalTok{.}\OperatorTok{:}\NormalTok{(}\OperatorTok{*}\NormalTok{)(x}\OperatorTok{::}\DataTypeTok{Metre}\NormalTok{, y}\OperatorTok{::}\DataTypeTok{Real}\NormalTok{) }\OperatorTok{=} \FunctionTok{Metre}\NormalTok{(x.value }\OperatorTok{*}\NormalTok{ y)}
\BuiltInTok{Base}\NormalTok{.}\OperatorTok{:}\NormalTok{(}\OperatorTok{*}\NormalTok{)(x}\OperatorTok{::}\DataTypeTok{Real}\NormalTok{, y}\OperatorTok{::}\DataTypeTok{Metre}\NormalTok{) }\OperatorTok{=} \FunctionTok{Metre}\NormalTok{(x }\OperatorTok{*}\NormalTok{ y.value)}
\BuiltInTok{Base}\NormalTok{.}\OperatorTok{:}\NormalTok{(}\OperatorTok{/}\NormalTok{)(x}\OperatorTok{::}\DataTypeTok{Metre}\NormalTok{, y}\OperatorTok{::}\DataTypeTok{Metre}\NormalTok{) }\OperatorTok{=}\NormalTok{ x.value }\OperatorTok{/}\NormalTok{ y.value}
\BuiltInTok{Base}\NormalTok{.}\OperatorTok{:}\NormalTok{(}\OperatorTok{/}\NormalTok{)(x}\OperatorTok{::}\DataTypeTok{Metre}\NormalTok{, y}\OperatorTok{::}\DataTypeTok{Real}\NormalTok{) }\OperatorTok{=} \FunctionTok{Metre}\NormalTok{(x.value }\OperatorTok{/}\NormalTok{ y)}
\end{Highlighting}
\end{Shaded}

Here we have added new methods to existing functions (from
\texttt{Base}). What we are doing here is opting-in to an existing
interface. If it makes sense to use existing interfaces - do so! Later
we'll create our own interface, but baby-steps first.

(In an object-oriented language, you would probably add a class with
some data and some function methods. It is very similar except for where
the methods live!)

We can make this a little ``nicer'' to use with one more definition.

\begin{Shaded}
\begin{Highlighting}[]
\KeywordTok{const}\NormalTok{ m }\OperatorTok{=} \FunctionTok{Metre}\NormalTok{(}\FloatTok{1}\NormalTok{)}
\end{Highlighting}
\end{Shaded}

This way we can write the code rather naturally:

\begin{Shaded}
\begin{Highlighting}[]
\FloatTok{3.5}\NormalTok{m}
\end{Highlighting}
\end{Shaded}

(Remember that's just short for \texttt{3.5\ *\ m}).

We can also make it print something similar to the terminal:

\begin{Shaded}
\begin{Highlighting}[]
\BuiltInTok{Base}\NormalTok{.}\FunctionTok{show}\NormalTok{(io}\OperatorTok{::}\DataTypeTok{IO}\NormalTok{, x}\OperatorTok{::}\DataTypeTok{Metre}\NormalTok{) }\OperatorTok{=} \FunctionTok{print}\NormalTok{(io, x.value, }\StringTok{"m"}\NormalTok{)}
\end{Highlighting}
\end{Shaded}

\subparagraph{How does this make things
safe?}\label{how-does-this-make-things-safe}

If I try and do something invalid like \texttt{1\ +\ 2m}, Julia will
throw an error. This means algebra with incorrect dimensions will not
compute in Julia. You'll get an error message and have the opportunity
to fix your code.

\subparagraph{Isn't all of this slow?}\label{isnt-all-of-this-slow}

Surprisingly, in Julia this doesn't slow anything down. The
\texttt{Metre\{Float64\}} object uses the same 64 bits to represent as a
bare \texttt{Float64}. The compiler is smart enough to know it can use
the same representation for either type interchangeably, and your
program won't spend any effort ``constructing structs'' or ``getting
fields'' at run time.

\paragraph{Other length units}\label{other-length-units}

How would you deal with other kinds of units of length?

\begin{itemize}
\tightlist
\item
  kilometres
\item
  centimetres
\item
  inches
\item
  miles
\end{itemize}

\emph{Q: Does anyone have any other favourite unit of length?}

As you can see, there might be a lot of different units of length. And
units of time, units of mass, units of temperature, etc. And the variety
grows further when you consider compound dimensions (like area, or
velocity). And then we need to consider \emph{binary} operators like
\texttt{+}, where the number of possible combinations is quadratic in
the number of units being considered.

How should we deal with all these combinations? How will the system know
how to add centimetres and yards? What is \texttt{3yd\ +\ 68cm}?

\begin{enumerate}
\def\labelenumi{\arabic{enumi}.}
\tightlist
\item
  Write every single method by hand?
\item
  Have a ``blessed'' unit (like SI units - kilograms, metres, seconds,
  Kelvin, etc)
\item
  Have some set of clever rules (e.g.~if both inputs are imperial, stay
  imperial)
\end{enumerate}

Option \#1 is not practical. Here we are going to look at option \#2.
Julia actually uses option \#3 for numeric conversion (what is
\texttt{0x20\ +\ 3.14}?), but we'll keep it a bit simpler.

This is a good place to introduce some abstract types. Let's focus on
just length units.

\begin{Shaded}
\begin{Highlighting}[]
\KeywordTok{abstract type}\NormalTok{ AbstractLength}
\KeywordTok{end}\NormalTok{;}
\end{Highlighting}
\end{Shaded}

This will be the ``supertype'' for our length types. (Remember ``super''
is Latin for ``above'' -- your supervisor is your overseer, etc).

\begin{Shaded}
\begin{Highlighting}[]
\KeywordTok{struct}\NormalTok{ Metre\{T\} }\OperatorTok{\textless{}:}\DataTypeTok{ AbstractLength}
\NormalTok{   value}\OperatorTok{::}\DataTypeTok{T}
\KeywordTok{end}

\KeywordTok{struct}\NormalTok{ Kilometre\{T\} }\OperatorTok{\textless{}:}\DataTypeTok{ AbstractLength}
\NormalTok{   value}\OperatorTok{::}\DataTypeTok{T}
\KeywordTok{end}

\KeywordTok{struct}\NormalTok{ Centimetre\{T\} }\OperatorTok{\textless{}:}\DataTypeTok{ AbstractLength}
\NormalTok{   value}\OperatorTok{::}\DataTypeTok{T}
\KeywordTok{end}

\KeywordTok{struct}\NormalTok{ Inch\{T\} }\OperatorTok{\textless{}:}\DataTypeTok{ AbstractLength}
\NormalTok{   value}\OperatorTok{::}\DataTypeTok{T}
\KeywordTok{end}

\KeywordTok{struct}\NormalTok{ Mile\{T\} }\OperatorTok{\textless{}:}\DataTypeTok{ AbstractLength}
\NormalTok{   value}\OperatorTok{::}\DataTypeTok{T}
\KeywordTok{end}\NormalTok{;}
\end{Highlighting}
\end{Shaded}

\subparagraph{Conversion to and from
metres}\label{conversion-to-and-from-metres}

For each of these, we might want a way to convert these to metres. We
can use the \texttt{Metre} type, and attach methods to that. Such
methods are known as ``constructors''.

\begin{Shaded}
\begin{Highlighting}[]
\FunctionTok{Metre}\NormalTok{(x}\OperatorTok{::}\DataTypeTok{Kilometre}\NormalTok{) }\OperatorTok{=} \FunctionTok{Metre}\NormalTok{(}\FloatTok{1000} \OperatorTok{*}\NormalTok{ x.value)}
\FunctionTok{Metre}\NormalTok{(x}\OperatorTok{::}\DataTypeTok{Centimetre}\NormalTok{) }\OperatorTok{=} \FunctionTok{Metre}\NormalTok{(}\FloatTok{0.01} \OperatorTok{*}\NormalTok{ x.value)}
\FunctionTok{Metre}\NormalTok{(x}\OperatorTok{::}\DataTypeTok{Inch}\NormalTok{) }\OperatorTok{=} \FunctionTok{Metre}\NormalTok{(}\FloatTok{0.0254} \OperatorTok{*}\NormalTok{ x.value)}
\FunctionTok{Metre}\NormalTok{(x}\OperatorTok{::}\DataTypeTok{Mile}\NormalTok{) }\OperatorTok{=} \FunctionTok{Metre}\NormalTok{(}\FloatTok{1760} \OperatorTok{*} \FloatTok{3} \OperatorTok{*} \FloatTok{12} \OperatorTok{*} \FloatTok{0.0254} \OperatorTok{*}\NormalTok{ x.value);}
\end{Highlighting}
\end{Shaded}

In the above, I know an inch is 2.54 centimetres, so I used that. I
don't recall how many metres per mile, but I can look up how many yards
per mile, and I recall there are 3 feet per yard and 12 inches per foot.

There is one more we should add, just to make sure things are never
ambiguous for Julia:

\begin{Shaded}
\begin{Highlighting}[]
\FunctionTok{Metre}\NormalTok{(x}\OperatorTok{::}\DataTypeTok{Metre}\NormalTok{) }\OperatorTok{=}\NormalTok{ x;}
\end{Highlighting}
\end{Shaded}

We'll also find it useful to go the other direction:

\begin{Shaded}
\begin{Highlighting}[]
\FunctionTok{Kilometre}\NormalTok{(x}\OperatorTok{::}\DataTypeTok{Metre}\NormalTok{) }\OperatorTok{=} \FunctionTok{Kilometre}\NormalTok{(}\FloatTok{0.001} \OperatorTok{*}\NormalTok{ x.value)}
\FunctionTok{Centimetre}\NormalTok{(x}\OperatorTok{::}\DataTypeTok{Metre}\NormalTok{) }\OperatorTok{=} \FunctionTok{Centimetre}\NormalTok{(}\FloatTok{100} \OperatorTok{*}\NormalTok{ x.value)}
\FunctionTok{Inch}\NormalTok{(x}\OperatorTok{::}\DataTypeTok{Metre}\NormalTok{) }\OperatorTok{=} \FunctionTok{Inch}\NormalTok{(x.value }\OperatorTok{/} \FloatTok{0.0254}\NormalTok{)}
\FunctionTok{Mile}\NormalTok{(x}\OperatorTok{::}\DataTypeTok{Metre}\NormalTok{) }\OperatorTok{=} \FunctionTok{Mile}\NormalTok{(x.value }\OperatorTok{/}\NormalTok{ (}\FloatTok{1760} \OperatorTok{*} \FloatTok{3} \OperatorTok{*} \FloatTok{12} \OperatorTok{*} \FloatTok{0.0254}\NormalTok{));}
\end{Highlighting}
\end{Shaded}

\subparagraph{Conversion between any two
types}\label{conversion-between-any-two-types}

If we are converting between two length units neither of which is
\texttt{Metre}, we might want to simply convert the input to metres
first. Doing this we end up with:

\begin{Shaded}
\begin{Highlighting}[]
\FunctionTok{Kilometre}\NormalTok{(x}\OperatorTok{::}\DataTypeTok{AbstractLength}\NormalTok{) }\OperatorTok{=} \FunctionTok{Kilometre}\NormalTok{(}\FunctionTok{Metre}\NormalTok{(x))}
\FunctionTok{Kilometre}\NormalTok{(x}\OperatorTok{::}\DataTypeTok{Kilometre}\NormalTok{) }\OperatorTok{=}\NormalTok{ x}

\FunctionTok{Centimetre}\NormalTok{(x}\OperatorTok{::}\DataTypeTok{AbstractLength}\NormalTok{) }\OperatorTok{=} \FunctionTok{Centimetre}\NormalTok{(}\FunctionTok{Metre}\NormalTok{(x))}
\FunctionTok{Centimetre}\NormalTok{(x}\OperatorTok{::}\DataTypeTok{Centimetre}\NormalTok{) }\OperatorTok{=}\NormalTok{ x}

\FunctionTok{Inch}\NormalTok{(x}\OperatorTok{::}\DataTypeTok{AbstractLength}\NormalTok{) }\OperatorTok{=} \FunctionTok{Inch}\NormalTok{(}\FunctionTok{Metre}\NormalTok{(x))}
\FunctionTok{Inch}\NormalTok{(x}\OperatorTok{::}\DataTypeTok{Inch}\NormalTok{) }\OperatorTok{=}\NormalTok{ x}

\FunctionTok{Mile}\NormalTok{(x}\OperatorTok{::}\DataTypeTok{AbstractLength}\NormalTok{) }\OperatorTok{=} \FunctionTok{Mile}\NormalTok{(}\FunctionTok{Metre}\NormalTok{(x))}
\FunctionTok{Mile}\NormalTok{(x}\OperatorTok{::}\DataTypeTok{Mile}\NormalTok{) }\OperatorTok{=}\NormalTok{ x;}
\end{Highlighting}
\end{Shaded}

Note that if we have \texttt{N} different length types, we'd only need
to define \texttt{2N} different conversions (to and from metres) instead
of \texttt{N\^{}2} conversions (between every two possibilities). If
\texttt{N} is large, this is hugely beneficial!

Another fun fact is two different programmers can define different
\texttt{AbstractLength} units without talking to each other. A third
programmer can use the code from the first two, and those units will
seamlessly work together with zero effort. This is great for sharing
working in teams, between teams, or between entirely separate
organizations (e.g.~open-source code).

Note that in general, we can get Julia to automatically do this step for
every \texttt{AbstractLength} type. The syntax for this one is a bit
crazy:

\begin{Shaded}
\begin{Highlighting}[]
\NormalTok{(}\OperatorTok{::}\DataTypeTok{Type\{T\}}\NormalTok{)(x}\OperatorTok{::}\DataTypeTok{AbstractLength}\NormalTok{) }\KeywordTok{where}\NormalTok{ \{T }\OperatorTok{\textless{}:}\DataTypeTok{ AbstractLength}\NormalTok{\} }\OperatorTok{=} \FunctionTok{T}\NormalTok{(}\FunctionTok{Metre}\NormalTok{(x));}
\NormalTok{(}\OperatorTok{::}\DataTypeTok{Type\{T\}}\NormalTok{)(x}\OperatorTok{::}\DataTypeTok{T}\NormalTok{) }\KeywordTok{where}\NormalTok{ \{T }\OperatorTok{\textless{}:}\DataTypeTok{ AbstractLength}\NormalTok{\} }\OperatorTok{=}\NormalTok{ x;}
\end{Highlighting}
\end{Shaded}

(Don't worry -- we won't expect you to create something like that in
this course!)

\subparagraph{Algebra with metres}\label{algebra-with-metres}

We have seen these already:

\begin{Shaded}
\begin{Highlighting}[]
\BuiltInTok{Base}\NormalTok{.}\OperatorTok{:}\NormalTok{(}\OperatorTok{+}\NormalTok{)(x}\OperatorTok{::}\DataTypeTok{Metre}\NormalTok{, y}\OperatorTok{::}\DataTypeTok{Metre}\NormalTok{) }\OperatorTok{=} \FunctionTok{Metre}\NormalTok{(x.value }\OperatorTok{+}\NormalTok{ y.value)}
\BuiltInTok{Base}\NormalTok{.}\OperatorTok{:}\NormalTok{(}\OperatorTok{{-}}\NormalTok{)(x}\OperatorTok{::}\DataTypeTok{Metre}\NormalTok{, y}\OperatorTok{::}\DataTypeTok{Metre}\NormalTok{) }\OperatorTok{=} \FunctionTok{Metre}\NormalTok{(x.value }\OperatorTok{{-}}\NormalTok{ y.value)}
\BuiltInTok{Base}\NormalTok{.}\OperatorTok{:}\NormalTok{(}\OperatorTok{{-}}\NormalTok{)(x}\OperatorTok{::}\DataTypeTok{Metre}\NormalTok{) }\OperatorTok{=} \FunctionTok{Metre}\NormalTok{(}\OperatorTok{{-}}\NormalTok{x.value)}
\BuiltInTok{Base}\NormalTok{.}\OperatorTok{:}\NormalTok{(}\OperatorTok{*}\NormalTok{)(x}\OperatorTok{::}\DataTypeTok{Metre}\NormalTok{, y}\OperatorTok{::}\DataTypeTok{Real}\NormalTok{) }\OperatorTok{=} \FunctionTok{Metre}\NormalTok{(x.value }\OperatorTok{*}\NormalTok{ y)}
\BuiltInTok{Base}\NormalTok{.}\OperatorTok{:}\NormalTok{(}\OperatorTok{*}\NormalTok{)(x}\OperatorTok{::}\DataTypeTok{Real}\NormalTok{, y}\OperatorTok{::}\DataTypeTok{Metre}\NormalTok{) }\OperatorTok{=} \FunctionTok{Metre}\NormalTok{(x }\OperatorTok{*}\NormalTok{ y.value)}
\BuiltInTok{Base}\NormalTok{.}\OperatorTok{:}\NormalTok{(}\OperatorTok{/}\NormalTok{)(x}\OperatorTok{::}\DataTypeTok{Metre}\NormalTok{, y}\OperatorTok{::}\DataTypeTok{Metre}\NormalTok{) }\OperatorTok{=}\NormalTok{ x.value }\OperatorTok{/}\NormalTok{ y.value}
\BuiltInTok{Base}\NormalTok{.}\OperatorTok{:}\NormalTok{(}\OperatorTok{/}\NormalTok{)(x}\OperatorTok{::}\DataTypeTok{Metre}\NormalTok{, y}\OperatorTok{::}\DataTypeTok{Real}\NormalTok{) }\OperatorTok{=} \FunctionTok{Metre}\NormalTok{(x.value }\OperatorTok{/}\NormalTok{ y)}
\end{Highlighting}
\end{Shaded}

\subparagraph{Algebra with abstract
types}\label{algebra-with-abstract-types}

Now we have enough to perform algebra between all these types. Let's
look at addition, subtraction and division:

\begin{Shaded}
\begin{Highlighting}[]
\BuiltInTok{Base}\NormalTok{.}\OperatorTok{:}\NormalTok{(}\OperatorTok{+}\NormalTok{)(x}\OperatorTok{::}\DataTypeTok{AbstractLength}\NormalTok{, y}\OperatorTok{::}\DataTypeTok{AbstractLength}\NormalTok{) }\OperatorTok{=} \FunctionTok{Metre}\NormalTok{(x) }\OperatorTok{+} \FunctionTok{Metre}\NormalTok{(y)}
\BuiltInTok{Base}\NormalTok{.}\OperatorTok{:}\NormalTok{(}\OperatorTok{{-}}\NormalTok{)(x}\OperatorTok{::}\DataTypeTok{AbstractLength}\NormalTok{, y}\OperatorTok{::}\DataTypeTok{AbstractLength}\NormalTok{) }\OperatorTok{=} \FunctionTok{Metre}\NormalTok{(x) }\OperatorTok{{-}} \FunctionTok{Metre}\NormalTok{(y)}
\BuiltInTok{Base}\NormalTok{.}\OperatorTok{:}\NormalTok{(}\OperatorTok{/}\NormalTok{)(x}\OperatorTok{::}\DataTypeTok{AbstractLength}\NormalTok{, y}\OperatorTok{::}\DataTypeTok{AbstractLength}\NormalTok{) }\OperatorTok{=} \FunctionTok{Metre}\NormalTok{(x) }\OperatorTok{/} \FunctionTok{Metre}\NormalTok{(y);}
\end{Highlighting}
\end{Shaded}

So to add, we just convert to metres first and then keep going.

For the other operations we can ``keep'' the existing type, but we need
to actually know the ``unit'' separately to its type parameter
(e.g.~\texttt{Centimetre} vs \texttt{Centimetre\{Int64\}} or
\texttt{Centimetre\{Float64\}}). We define a brand-new function to
describe this:

\begin{Shaded}
\begin{Highlighting}[]
\FunctionTok{unit}\NormalTok{(}\OperatorTok{::}\DataTypeTok{Metre}\NormalTok{) }\OperatorTok{=}\NormalTok{ Metre}
\FunctionTok{unit}\NormalTok{(}\OperatorTok{::}\DataTypeTok{Kilometre}\NormalTok{) }\OperatorTok{=}\NormalTok{ Kilometre}
\FunctionTok{unit}\NormalTok{(}\OperatorTok{::}\DataTypeTok{Centimetre}\NormalTok{) }\OperatorTok{=}\NormalTok{ Centimetre}
\FunctionTok{unit}\NormalTok{(}\OperatorTok{::}\DataTypeTok{Inch}\NormalTok{) }\OperatorTok{=}\NormalTok{ Inch}
\FunctionTok{unit}\NormalTok{(}\OperatorTok{::}\DataTypeTok{Mile}\NormalTok{) }\OperatorTok{=}\NormalTok{ Mile}
\end{Highlighting}
\end{Shaded}

\begin{verbatim}
unit (generic function with 5 methods)
\end{verbatim}

Using this we can define ``generic'' methods for these functions:

\begin{Shaded}
\begin{Highlighting}[]
\BuiltInTok{Base}\NormalTok{.}\OperatorTok{:}\NormalTok{(}\OperatorTok{{-}}\NormalTok{)(x}\OperatorTok{::}\DataTypeTok{AbstractLength}\NormalTok{) }\OperatorTok{=} \FunctionTok{unit}\NormalTok{(x)(}\OperatorTok{{-}}\NormalTok{x.value)}
\BuiltInTok{Base}\NormalTok{.}\OperatorTok{:}\NormalTok{(}\OperatorTok{*}\NormalTok{)(x}\OperatorTok{::}\DataTypeTok{AbstractLength}\NormalTok{, y}\OperatorTok{::}\DataTypeTok{Real}\NormalTok{) }\OperatorTok{=} \FunctionTok{unit}\NormalTok{(x)(x.value }\OperatorTok{*}\NormalTok{ y)}
\BuiltInTok{Base}\NormalTok{.}\OperatorTok{:}\NormalTok{(}\OperatorTok{*}\NormalTok{)(x}\OperatorTok{::}\DataTypeTok{Real}\NormalTok{, y}\OperatorTok{::}\DataTypeTok{AbstractLength}\NormalTok{) }\OperatorTok{=} \FunctionTok{unit}\NormalTok{(y)(x }\OperatorTok{*}\NormalTok{ y.value)}
\BuiltInTok{Base}\NormalTok{.}\OperatorTok{:}\NormalTok{(}\OperatorTok{/}\NormalTok{)(x}\OperatorTok{::}\DataTypeTok{AbstractLength}\NormalTok{, y}\OperatorTok{::}\DataTypeTok{Real}\NormalTok{) }\OperatorTok{=} \FunctionTok{unit}\NormalTok{(x)(x.value }\OperatorTok{/}\NormalTok{ y);}
\end{Highlighting}
\end{Shaded}

If we want to be a bit more fancy, we can get Julia to skip the
conversion for addition, subtraction and division when it is
unnecessary. There are different ways to do this like fancy dispatch
rules, but you can just work with types directly \emph{inside} the
functions:

\begin{Shaded}
\begin{Highlighting}[]
\KeywordTok{function} \BuiltInTok{Base}\NormalTok{.}\OperatorTok{:}\NormalTok{(}\OperatorTok{+}\NormalTok{)(x}\OperatorTok{::}\DataTypeTok{AbstractLength}\NormalTok{, y}\OperatorTok{::}\DataTypeTok{AbstractLength}\NormalTok{)}
   \ControlFlowTok{if} \FunctionTok{unit}\NormalTok{(x) }\OperatorTok{==} \FunctionTok{unit}\NormalTok{(y)}
   \ControlFlowTok{return} \FunctionTok{unit}\NormalTok{(x)(x.value }\OperatorTok{+}\NormalTok{ y.value)}
   \ControlFlowTok{else}
   \ControlFlowTok{return} \FunctionTok{Metre}\NormalTok{(x) }\OperatorTok{+} \FunctionTok{Metre}\NormalTok{(y)}
   \ControlFlowTok{end}
\KeywordTok{end}

\KeywordTok{function} \BuiltInTok{Base}\NormalTok{.}\OperatorTok{:}\NormalTok{(}\OperatorTok{{-}}\NormalTok{)(x}\OperatorTok{::}\DataTypeTok{AbstractLength}\NormalTok{, y}\OperatorTok{::}\DataTypeTok{AbstractLength}\NormalTok{)}
   \ControlFlowTok{if} \FunctionTok{unit}\NormalTok{(x) }\OperatorTok{==} \FunctionTok{unit}\NormalTok{(y)}
   \ControlFlowTok{return} \FunctionTok{unit}\NormalTok{(x)(x.value }\OperatorTok{{-}}\NormalTok{ y.value)}
   \ControlFlowTok{else}
   \ControlFlowTok{return} \FunctionTok{Metre}\NormalTok{(x) }\OperatorTok{{-}} \FunctionTok{Metre}\NormalTok{(y)}
   \ControlFlowTok{end}
\KeywordTok{end}

\KeywordTok{function} \BuiltInTok{Base}\NormalTok{.}\OperatorTok{:}\NormalTok{(}\OperatorTok{/}\NormalTok{)(x}\OperatorTok{::}\DataTypeTok{AbstractLength}\NormalTok{, y}\OperatorTok{::}\DataTypeTok{AbstractLength}\NormalTok{)}
   \ControlFlowTok{if} \FunctionTok{unit}\NormalTok{(x) }\OperatorTok{==} \FunctionTok{unit}\NormalTok{(y)}
   \ControlFlowTok{return}\NormalTok{ x.value }\OperatorTok{/}\NormalTok{ y.value}
   \ControlFlowTok{else}
   \ControlFlowTok{return} \FunctionTok{Metre}\NormalTok{(x) }\OperatorTok{/} \FunctionTok{Metre}\NormalTok{(y)}
   \ControlFlowTok{end}
\KeywordTok{end}\NormalTok{;}
\end{Highlighting}
\end{Shaded}

(Note that the \texttt{if} part will be dealt with at compile time.)

\subparagraph{Niceties}\label{niceties}

\begin{Shaded}
\begin{Highlighting}[]
\KeywordTok{const}\NormalTok{ m }\OperatorTok{=} \FunctionTok{Metre}\NormalTok{(}\FloatTok{1}\NormalTok{) }\CommentTok{\# done above}
\KeywordTok{const}\NormalTok{ km }\OperatorTok{=} \FunctionTok{Kilometre}\NormalTok{(}\FloatTok{1}\NormalTok{)}
\KeywordTok{const}\NormalTok{ cm }\OperatorTok{=} \FunctionTok{Centimetre}\NormalTok{(}\FloatTok{1}\NormalTok{)}
\KeywordTok{const}\NormalTok{ inch }\OperatorTok{=} \FunctionTok{Inch}\NormalTok{(}\FloatTok{1}\NormalTok{)}
\KeywordTok{const}\NormalTok{ mile }\OperatorTok{=} \FunctionTok{Mile}\NormalTok{(}\FloatTok{1}\NormalTok{);}
\end{Highlighting}
\end{Shaded}

\begin{Shaded}
\begin{Highlighting}[]
\BuiltInTok{Base}\NormalTok{.}\FunctionTok{show}\NormalTok{(io}\OperatorTok{::}\DataTypeTok{IO}\NormalTok{, x}\OperatorTok{::}\DataTypeTok{Metre}\NormalTok{) }\OperatorTok{=} \FunctionTok{print}\NormalTok{(io, x.value, }\StringTok{"m"}\NormalTok{)}
\BuiltInTok{Base}\NormalTok{.}\FunctionTok{show}\NormalTok{(io}\OperatorTok{::}\DataTypeTok{IO}\NormalTok{, x}\OperatorTok{::}\DataTypeTok{Kilometre}\NormalTok{) }\OperatorTok{=} \FunctionTok{print}\NormalTok{(io, x.value, }\StringTok{"km"}\NormalTok{)}
\BuiltInTok{Base}\NormalTok{.}\FunctionTok{show}\NormalTok{(io}\OperatorTok{::}\DataTypeTok{IO}\NormalTok{, x}\OperatorTok{::}\DataTypeTok{Centimetre}\NormalTok{) }\OperatorTok{=} \FunctionTok{print}\NormalTok{(io, x.value, }\StringTok{"cm"}\NormalTok{)}
\BuiltInTok{Base}\NormalTok{.}\FunctionTok{show}\NormalTok{(io}\OperatorTok{::}\DataTypeTok{IO}\NormalTok{, x}\OperatorTok{::}\DataTypeTok{Inch}\NormalTok{) }\OperatorTok{=} \FunctionTok{print}\NormalTok{(io, x.value, }\StringTok{"inch"}\NormalTok{)}
\BuiltInTok{Base}\NormalTok{.}\FunctionTok{show}\NormalTok{(io}\OperatorTok{::}\DataTypeTok{IO}\NormalTok{, x}\OperatorTok{::}\DataTypeTok{Mile}\NormalTok{) }\OperatorTok{=} \FunctionTok{print}\NormalTok{(io, x.value, }\StringTok{"mile"}\NormalTok{);}
\end{Highlighting}
\end{Shaded}

\subparagraph{Some results}\label{some-results}

\emph{Q: How many kilometres in a mile?}

\begin{Shaded}
\begin{Highlighting}[]
\NormalTok{mile }\OperatorTok{/}\NormalTok{ km}
\end{Highlighting}
\end{Shaded}

\begin{verbatim}
1.6093439999999999
\end{verbatim}

\emph{Q: Jenny is 5'9'' (5 feet, 9 inches) tall, and has a ladder 210cm
long. Can she reach a ceiling 4m high to change a light bulb?}

\begin{Shaded}
\begin{Highlighting}[]
\NormalTok{(}\FloatTok{5}\OperatorTok{*}\FloatTok{12} \OperatorTok{+} \FloatTok{9}\NormalTok{)inch }\OperatorTok{+} \FloatTok{210}\NormalTok{cm}
\end{Highlighting}
\end{Shaded}

\begin{verbatim}
3.8526m
\end{verbatim}

\emph{A: Her head won't quite hit the ceiling, but Jenny will be able to
reach the light with her arms.}

\subparagraph{New interfaces}\label{new-interfaces}

So far we have mostly overloaded existing functions built into
\texttt{Base} Julia. There is nothing stopping you defining your own
functions -- in fact you'll likely define more new functions that
overload existing ones. Here we also added the \texttt{unit} function to
help out.

But how did we know to overload the \texttt{Base} functions? Well, these
functions are \emph{interfaces} used throughout Julia. We can define new
interfaces for our new \emph{abstract} types. (Typically, a function we
write that accepts some concrete type is not a \emph{generic interface}
-- it is just a one-off function with a single method).

Now we have our \texttt{AbstractLength} interface, we can document how
people \emph{implement} it so people could add new types, like
\texttt{Nanometres} and \texttt{Yard}. To \emph{implement the
\texttt{AbstractLength} interface} you'd need to implement:

\begin{itemize}
\tightlist
\item
  Conversion to and from \texttt{Metre}
\item
  Define a method for \texttt{unit}
\item
  Optionally, define a \texttt{const} unit value and overload
  \texttt{show}
\end{itemize}

After this, everything will ``just work''™. This is the power of
interfaces!

\paragraph{Other dimensions}\label{other-dimensions}

We could do the same for time, mass, temperature, etc. But then we need
to deal with compound dimensions like area, velocity, energy, etc.
Again, we need a clever interface to deal with all this complexity. This
(and more) is implemented in the
\href{https://painterqubits.github.io/Unitful.jl/stable/}{Unitful.jl}
package.

\subsubsection{Making types and functions together --
shapes}\label{making-types-and-functions-together-shapes}

This time we'll practice making new types together with functions.

For this example we'll look at geometric shapes, starting with
two-dimensional shapes.

\paragraph{Abstract type and
interface}\label{abstract-type-and-interface}

We can start with an abstract type for two-dimensional shapes.

\begin{Shaded}
\begin{Highlighting}[]
\KeywordTok{abstract type}\NormalTok{ Shape2D}
\KeywordTok{end}\NormalTok{;}
\end{Highlighting}
\end{Shaded}

Even \emph{before} we talk about what kinds of 2D shapes there are, we
want to think about what you can \emph{do} with 2D shapes. For our
purposes let's look at finding a couple of properties like perimeter and
area.

\begin{Shaded}
\begin{Highlighting}[]
\StringTok{"""}
\StringTok{   perimeter(shape::2DShape)}

\StringTok{Return the perimeter of \textasciigrave{}shape\textasciigrave{}.}
\StringTok{"""}
\KeywordTok{function}\NormalTok{ perimeter}
\KeywordTok{end}

\StringTok{"""}
\StringTok{   area(shape::2DShape)}

\StringTok{Return the area of \textasciigrave{}shape\textasciigrave{}.}
\StringTok{"""}
\KeywordTok{function}\NormalTok{ area}
\KeywordTok{end}\NormalTok{;}
\end{Highlighting}
\end{Shaded}

Here we've defined two functions, each with zero methods but some
documentation explaining their purpose (interfaces in Julia are
informal, and are specified through documentation). You can find the
documentation with \texttt{?area} at the REPL.

\paragraph{Types of 2D shape}\label{types-of-2d-shape}

There are various kinds of 2D shapes. Let's introduce a few here, and
organize them by the number of sides. Here is some triangles:

\begin{Shaded}
\begin{Highlighting}[]
\KeywordTok{abstract type}\NormalTok{ Triangle }\OperatorTok{\textless{}:}\DataTypeTok{ Shape2D}
\KeywordTok{end}

\KeywordTok{struct}\NormalTok{ EquilateralTriangle\{T\} }\OperatorTok{\textless{}:}\DataTypeTok{ Triangle}
\NormalTok{   side}\OperatorTok{::}\DataTypeTok{T}
\KeywordTok{end}

\KeywordTok{struct}\NormalTok{ RightAngledTriangle\{T\} }\OperatorTok{\textless{}:}\DataTypeTok{ Triangle}
\NormalTok{   height}\OperatorTok{::}\DataTypeTok{T}
\NormalTok{   width}\OperatorTok{::}\DataTypeTok{T}
\KeywordTok{end}\NormalTok{;}
\end{Highlighting}
\end{Shaded}

And here are some quadrilaterals:

\begin{Shaded}
\begin{Highlighting}[]
\KeywordTok{abstract type}\NormalTok{ Quadrilateral }\OperatorTok{\textless{}:}\DataTypeTok{ Shape2D}
\KeywordTok{end}

\KeywordTok{struct}\NormalTok{ Square\{T\} }\OperatorTok{\textless{}:}\DataTypeTok{ Quadrilateral}
\NormalTok{   side}\OperatorTok{::}\DataTypeTok{T}
\KeywordTok{end}

\KeywordTok{struct}\NormalTok{ Rectangle\{T\} }\OperatorTok{\textless{}:}\DataTypeTok{ Quadrilateral}
\NormalTok{   height}\OperatorTok{::}\DataTypeTok{T}
\NormalTok{   width}\OperatorTok{::}\DataTypeTok{T}
\KeywordTok{end}\NormalTok{;}
\end{Highlighting}
\end{Shaded}

And here is a circle, a direct subtype of \texttt{Shape2D}:

\begin{Shaded}
\begin{Highlighting}[]
\KeywordTok{struct}\NormalTok{ Circle\{T\} }\OperatorTok{\textless{}:}\DataTypeTok{ Shape2D}
\NormalTok{   radius}\OperatorTok{::}\DataTypeTok{T}
\KeywordTok{end}\NormalTok{;}
\end{Highlighting}
\end{Shaded}

Sometimes it's useful to \textbf{use values which aren't types as type
parameters (``Value types'')}. One standard example is the \texttt{N} in
\texttt{Array\{T,N\}} which is just an integer signifying the dimension
of the array. Doing this helps the compiler know some parts of the
``problem size'' at compile time based on the type and can lead to much
more efficient code or data storage. It can also let us cover several
different sub-cases with one type definition. In our \texttt{Shape2D}
example we can use it to define a regular polygon with \texttt{N} sides
as follows:

\begin{Shaded}
\begin{Highlighting}[]
\KeywordTok{struct}\NormalTok{ RegularPolygon\{N, T\} }\OperatorTok{\textless{}:}\DataTypeTok{ Shape2D}
\NormalTok{    side\_length}\OperatorTok{::}\DataTypeTok{T}
\KeywordTok{end} 

\CommentTok{\# and a constructor for convenience}
\KeywordTok{function} \FunctionTok{RegularPolygon}\DataTypeTok{\{N\}}\NormalTok{(side) }\KeywordTok{where}\NormalTok{ N}
   \FunctionTok{RegularPolygon}\DataTypeTok{\{N,typeof(side)\}}\NormalTok{(side)}
\KeywordTok{end}

\CommentTok{\# const Square = RegularPolygon\{4\}  \# We already defined squares above}
\KeywordTok{const}\NormalTok{ Pentagon }\OperatorTok{=}\NormalTok{ RegularPolygon\{}\FloatTok{5}\NormalTok{\}}
\KeywordTok{const}\NormalTok{ Hexagon }\OperatorTok{=}\NormalTok{ RegularPolygon\{}\FloatTok{6}\NormalTok{\}}
\end{Highlighting}
\end{Shaded}

\begin{verbatim}
RegularPolygon{6}
\end{verbatim}

(One limitation of this is that we we can't have
\texttt{RegularPolgyon\{4\}\ \textless{}:\ Quadrilateral} in this
generic definition so if \texttt{Quadrilateral} was important to our
design we may have some thinking to do here.)

\paragraph{\texorpdfstring{Implementing the \texttt{Shape2D}
interface}{Implementing the Shape2D interface}}\label{implementing-the-shape2d-interface}

For each of these shapes we can determine the perimeter:

\begin{Shaded}
\begin{Highlighting}[]
\KeywordTok{function} \FunctionTok{perimeter}\NormalTok{(shape}\OperatorTok{::}\DataTypeTok{EquilateralTriangle}\NormalTok{)}
   \ControlFlowTok{return} \FloatTok{3} \OperatorTok{*}\NormalTok{ shape.side}
\KeywordTok{end}

\KeywordTok{function} \FunctionTok{perimeter}\NormalTok{(shape}\OperatorTok{::}\DataTypeTok{RightAngledTriangle}\NormalTok{)}
\NormalTok{   angle }\OperatorTok{=} \FunctionTok{atan}\NormalTok{(shape.height }\OperatorTok{/}\NormalTok{ shape.width)}
   \ControlFlowTok{return}\NormalTok{ shape.width }\OperatorTok{+}\NormalTok{ shape.height }\OperatorTok{+}\NormalTok{ shape.width }\OperatorTok{/} \FunctionTok{cos}\NormalTok{(angle)}
\KeywordTok{end}

\KeywordTok{function} \FunctionTok{perimeter}\NormalTok{(shape}\OperatorTok{::}\DataTypeTok{Square}\NormalTok{)}
   \ControlFlowTok{return} \FloatTok{4} \OperatorTok{*}\NormalTok{ shape.side}
\KeywordTok{end}

\KeywordTok{function} \FunctionTok{perimeter}\NormalTok{(shape}\OperatorTok{::}\DataTypeTok{Rectangle}\NormalTok{)}
   \ControlFlowTok{return} \FloatTok{2} \OperatorTok{*}\NormalTok{ shape.width }\OperatorTok{+} \FloatTok{2} \OperatorTok{*}\NormalTok{ shape.height}
\KeywordTok{end}

\KeywordTok{function} \FunctionTok{perimeter}\NormalTok{(shape}\OperatorTok{::}\DataTypeTok{Circle}\NormalTok{)}
   \ControlFlowTok{return} \FloatTok{2}\NormalTok{π }\OperatorTok{*}\NormalTok{ shape.radius}
\KeywordTok{end}

\KeywordTok{function} \FunctionTok{perimeter}\NormalTok{(shape}\OperatorTok{::}\DataTypeTok{RegularPolygon\{N, T\}}\NormalTok{) }\KeywordTok{where}\NormalTok{ \{N, T\}}
\NormalTok{    return N }\OperatorTok{*}\NormalTok{ shape.side\_length}
\KeywordTok{end}\NormalTok{;}
\end{Highlighting}
\end{Shaded}

And the area:

\begin{Shaded}
\begin{Highlighting}[]
\KeywordTok{function} \FunctionTok{area}\NormalTok{(shape}\OperatorTok{::}\DataTypeTok{EquilateralTriangle}\NormalTok{)}
   \ControlFlowTok{return} \FloatTok{0.5} \OperatorTok{*} \FunctionTok{sind}\NormalTok{(}\FloatTok{60}\NormalTok{) }\OperatorTok{*}\NormalTok{ shape.side }\OperatorTok{*}\NormalTok{ shape.side}
\KeywordTok{end}

\KeywordTok{function} \FunctionTok{area}\NormalTok{(shape}\OperatorTok{::}\DataTypeTok{RightAngledTriangle}\NormalTok{)}
   \ControlFlowTok{return} \FloatTok{0.5} \OperatorTok{*}\NormalTok{ shape.width }\OperatorTok{+}\NormalTok{ shape.height}
\KeywordTok{end}

\KeywordTok{function} \FunctionTok{area}\NormalTok{(shape}\OperatorTok{::}\DataTypeTok{Square}\NormalTok{)}
   \ControlFlowTok{return}\NormalTok{ shape.side }\OperatorTok{*}\NormalTok{ shape.side}
\KeywordTok{end}

\KeywordTok{function} \FunctionTok{area}\NormalTok{(shape}\OperatorTok{::}\DataTypeTok{Rectangle}\NormalTok{)}
   \ControlFlowTok{return}\NormalTok{ shape.width }\OperatorTok{*}\NormalTok{ shape.height}
\KeywordTok{end}

\KeywordTok{function} \FunctionTok{area}\NormalTok{(shape}\OperatorTok{::}\DataTypeTok{Circle}\NormalTok{)}
   \ControlFlowTok{return} \ConstantTok{π} \OperatorTok{*}\NormalTok{ shape.radius }\OperatorTok{*}\NormalTok{ shape.radius}
\KeywordTok{end}

\KeywordTok{function} \FunctionTok{area}\NormalTok{(shape}\OperatorTok{::}\DataTypeTok{RegularPolygon\{N, T\}}\NormalTok{) }\KeywordTok{where}\NormalTok{ \{N, T\}}
\NormalTok{    return  (N }\OperatorTok{*}\NormalTok{ shape.side\_length}\OperatorTok{\^{}}\FloatTok{2}\NormalTok{) }\OperatorTok{/}\NormalTok{ (}\FloatTok{4} \OperatorTok{*} \FunctionTok{tand}\NormalTok{(}\FloatTok{180}\OperatorTok{/}\NormalTok{N))}
\KeywordTok{end}\NormalTok{;}
\end{Highlighting}
\end{Shaded}

We can use these pretty straightforwardly.

\begin{Shaded}
\begin{Highlighting}[]
\PreprocessorTok{@show} \FunctionTok{perimeter}\NormalTok{(}\FunctionTok{Rectangle}\NormalTok{(}\FloatTok{3.5}\NormalTok{, }\FloatTok{2.0}\NormalTok{))}
\PreprocessorTok{@show} \FunctionTok{perimeter}\NormalTok{(}\FunctionTok{Hexagon}\NormalTok{(}\FloatTok{1.5}\NormalTok{))}
\PreprocessorTok{@show} \FunctionTok{area}\NormalTok{(}\FunctionTok{RightAngledTriangle}\NormalTok{(}\FloatTok{10}\NormalTok{, }\FloatTok{10}\NormalTok{));}
\end{Highlighting}
\end{Shaded}

\begin{verbatim}
perimeter(Rectangle(3.5, 2.0)) = 11.0
perimeter(Hexagon(1.5)) = 9.0
area(RightAngledTriangle(10, 10)) = 15.0
\end{verbatim}

At any point in time we can introduce new subtypes of \texttt{Shape2D},
\texttt{Triangle} or \texttt{Rectangle} so long as we define their
perimeter and area. The person who wrote \texttt{Shape2d} does not even
need to be aware.

\subsubsection{Combining interfaces}\label{combining-interfaces}

Today we have introduced two new abstract types with their own
interfaces. One for units of length, and one for shapes.

We have said Julia is good for quickly writing mathematical code, and
good for writing code that executes fast. Here we can see how Julia
``scales'' well when combining code together -- everything will ``just
work''.

Take this:

\begin{Shaded}
\begin{Highlighting}[]
\FunctionTok{perimeter}\NormalTok{(}\FunctionTok{Square}\NormalTok{(}\FloatTok{4}\NormalTok{cm))}
\end{Highlighting}
\end{Shaded}

\begin{verbatim}
16cm
\end{verbatim}

Here, the code that handles \texttt{perimeter} doesn't need to know
anything about \texttt{AbstractLength}. The only thing it relies on is
that the lengths in question work with \texttt{+}, \texttt{*},
\texttt{/}, etc.

\begin{Shaded}
\begin{Highlighting}[]
\FunctionTok{perimeter}\NormalTok{(}\FunctionTok{Circle}\NormalTok{(}\FloatTok{10.0}\NormalTok{cm))}
\end{Highlighting}
\end{Shaded}

\begin{verbatim}
62.83185307179586cm
\end{verbatim}

However, in our units implementation we can't (yet) multiply lengths to
get units of area (dimensions of length squared).

\begin{Shaded}
\begin{Highlighting}[]
\FunctionTok{area}\NormalTok{(}\FunctionTok{Circle}\NormalTok{(}\FloatTok{10.0}\NormalTok{cm))}
\end{Highlighting}
\end{Shaded}

\begin{Highlighting}
\textcolor{black}{MethodError: MethodError(*, (31.41592653589793cm, 10.0cm), 0x00000000000099b9)}
\textcolor{black}{MethodError: no method matching *(::Centimetre\{Float64\}, ::Centimetre\{Float64\})}
\textcolor{black}{The function `*` exists, but no method is defined for this combination of argument types.}
\textcolor{black}{}\textcolor{QuartoInternalColor1}{Closest candidates are:}
\textcolor{QuartoInternalColor1}{}\textcolor{QuartoInternalColor1}{  *(::Any, ::Any, }\textcolor{QuartoInternalColor1}{::Any}\textcolor{QuartoInternalColor1}{, }\textcolor{QuartoInternalColor1}{::Any...}\textcolor{QuartoInternalColor1}{)}
\textcolor{QuartoInternalColor1}{}\textcolor{QuartoInternalColor1}{}\textcolor{QuartoInternalColor1}{   @}\textcolor{QuartoInternalColor1}{ }\textcolor{QuartoInternalColor1}{Base}\textcolor{QuartoInternalColor1}{ }\textcolor{QuartoInternalColor1}{}\textcolor{QuartoInternalColor1}{operators.jl:642}\textcolor{QuartoInternalColor1}{}\textcolor{QuartoInternalColor1}{}
\textcolor{QuartoInternalColor1}{}\textcolor{QuartoInternalColor1}{  *(}\textcolor{QuartoInternalColor1}{::Real}\textcolor{QuartoInternalColor1}{, ::AbstractLength)}
\textcolor{QuartoInternalColor1}{}\textcolor{QuartoInternalColor1}{}\textcolor{QuartoInternalColor1}{   @}\textcolor{QuartoInternalColor1}{ }\textcolor{QuartoInternalColor2}{Main.Notebook}\textcolor{QuartoInternalColor1}{ }\textcolor{QuartoInternalColor1}{\textasciitilde{}/nextcloud/MATH2504\_2026/lecture/quarto/}\textcolor{QuartoInternalColor1}{}\textcolor{QuartoInternalColor1}{}\textcolor{QuartoInternalColor1}{lecture-unit-5.qmd:1778}\textcolor{QuartoInternalColor1}{}\textcolor{QuartoInternalColor1}{}
\textcolor{QuartoInternalColor1}{}\textcolor{QuartoInternalColor1}{  *(}\textcolor{QuartoInternalColor1}{::SpecialFunctions.SimplePoly}\textcolor{QuartoInternalColor1}{, ::Any)}
\textcolor{QuartoInternalColor1}{}\textcolor{QuartoInternalColor1}{}\textcolor{QuartoInternalColor1}{   @}\textcolor{QuartoInternalColor1}{ }\textcolor{QuartoInternalColor3}{SpecialFunctions}\textcolor{QuartoInternalColor1}{ }\textcolor{QuartoInternalColor1}{\textasciitilde{}/.julia/packages/SpecialFunctions/mf0qH/src/}\textcolor{QuartoInternalColor1}{}\textcolor{QuartoInternalColor1}{}\textcolor{QuartoInternalColor1}{expint.jl:8}\textcolor{QuartoInternalColor1}{}\textcolor{QuartoInternalColor1}{}
\textcolor{QuartoInternalColor1}{}\textcolor{QuartoInternalColor1}{  ...}
\textcolor{QuartoInternalColor1}{Stacktrace:}
\textcolor{QuartoInternalColor1}{ [1] }\textcolor{QuartoInternalColor1}{}\textcolor{QuartoInternalColor1}{*}\textcolor{QuartoInternalColor1}{}
\textcolor{QuartoInternalColor1}{}\textcolor{QuartoInternalColor1}{   @}\textcolor{QuartoInternalColor1}{ }\textcolor{QuartoInternalColor1}{./}\textcolor{QuartoInternalColor1}{}\textcolor{QuartoInternalColor1}{}\textcolor{QuartoInternalColor1}{operators.jl:642}\textcolor{QuartoInternalColor1}{}\textcolor{QuartoInternalColor1}{}\textcolor{QuartoInternalColor1}{ [inlined]}\textcolor{QuartoInternalColor1}{}
\textcolor{QuartoInternalColor1}{ [2] }\textcolor{QuartoInternalColor1}{}\textcolor{QuartoInternalColor1}{area}\textcolor{QuartoInternalColor1}{}\textcolor{QuartoInternalColor1}{}\textcolor{QuartoInternalColor1}{(}\textcolor{QuartoInternalColor1}{}\textcolor{QuartoInternalColor1}{shape}\textcolor{QuartoInternalColor1}{::}\textcolor{QuartoInternalColor1}{Circle}\textcolor{QuartoInternalColor1}{\{Centimetre\{Float64\}\}}\textcolor{QuartoInternalColor1}{}\textcolor{QuartoInternalColor1}{}\textcolor{QuartoInternalColor1}{)}\textcolor{QuartoInternalColor1}{}
\textcolor{QuartoInternalColor1}{}\textcolor{QuartoInternalColor1}{   @}\textcolor{QuartoInternalColor1}{ }\textcolor{QuartoInternalColor2}{Main.Notebook}\textcolor{QuartoInternalColor1}{ }\textcolor{QuartoInternalColor1}{\textasciitilde{}/nextcloud/MATH2504\_2026/lecture/quarto/}\textcolor{QuartoInternalColor1}{}\textcolor{QuartoInternalColor1}{}\textcolor{QuartoInternalColor1}{lecture-unit-5.qmd:2023}\textcolor{QuartoInternalColor1}{}\textcolor{QuartoInternalColor1}{}
\textcolor{QuartoInternalColor1}{ [3] top-level scope}
\textcolor{QuartoInternalColor1}{}\textcolor{QuartoInternalColor1}{   @}\textcolor{QuartoInternalColor1}{ }\textcolor{QuartoInternalColor1}{\textasciitilde{}/nextcloud/MATH2504\_2026/lecture/quarto/}\textcolor{QuartoInternalColor1}{}\textcolor{QuartoInternalColor1}{}\textcolor{QuartoInternalColor1}{lecture-unit-5.qmd:2062}\textcolor{QuartoInternalColor1}{}\textcolor{QuartoInternalColor1}{}
\end{Highlighting}

If instead we used a more thorough implementation of units, like
\href{https://painterqubits.github.io/Unitful.jl/stable/}{Unitful.jl},
this will work fine.

\begin{Shaded}
\begin{Highlighting}[]
\ImportTok{using} \BuiltInTok{Unitful}

\FunctionTok{area}\NormalTok{(}\FunctionTok{Circle}\NormalTok{(}\FloatTok{10.0}\NormalTok{u}\StringTok{"cm"}\NormalTok{))}
\end{Highlighting}
\end{Shaded}

\begin{verbatim}
314.1592653589793 cm^2
\end{verbatim}

\subsection{Going forward}\label{going-forward}

We are now done with the material for Unit 5. We have learned:

\begin{enumerate}
\def\labelenumi{\arabic{enumi}.}
\tightlist
\item
  What are the ``concrete'' types represents sets of values in Julia?

  \begin{itemize}
  \tightlist
  \item
    \texttt{primitive\ type}
  \item
    \texttt{struct}, \texttt{Tuple}, \texttt{NamedTuple}
  \item
    \texttt{Array}, \texttt{String}
  \end{itemize}
\item
  What are the ``abstract'' types that represent sets of types in Julia?

  \begin{itemize}
  \tightlist
  \item
    \texttt{abstract\ type}
  \item
    Generic type parameters, \texttt{UnionAll}
  \item
    \texttt{Union} of types
  \end{itemize}
\item
  How do functions work in Julia?

  \begin{itemize}
  \tightlist
  \item
    functions and closures
  \item
    methods
  \item
    multiple dispatch
  \end{itemize}
\item
  How do we create types and interfaces in Julia?

  \begin{itemize}
  \tightlist
  \item
    length units
  \item
    2D shapes
  \item
    how interfaces compose naturally
  \end{itemize}
\end{enumerate}

This will place you in good stead for Units 6-8, and in particular
Projects 1-3 where you will need to craft non-trivial programs to create
a solution.

\subsection{Bonus material}\label{bonus-material}

Consider the following material optional. Have a look through if you are
curious, have extra time, or would like extra exposure and practice to
defining types and interfaces.

\subsubsection{Online statistics}\label{online-statistics}

Lets consider an example where we want to collect some quick running
statistics from some data coming in. E.g.:

\begin{Shaded}
\begin{Highlighting}[]
\ImportTok{using} \BuiltInTok{Statistics}

\CommentTok{\# A function that returns a data point {-} could be streamed from disk or the internet}
\FunctionTok{fetch\_new\_data}\NormalTok{() }\OperatorTok{=} \FloatTok{100}\FunctionTok{*rand}\NormalTok{()}

\KeywordTok{function} \FunctionTok{print\_running\_stats}\NormalTok{(n}\OperatorTok{::}\DataTypeTok{Integer}\NormalTok{, running\_stats\_storage)}
   \ControlFlowTok{for}\NormalTok{ i }\KeywordTok{in} \FloatTok{1}\OperatorTok{:}\NormalTok{n}
   \CommentTok{\# Get some data}
\NormalTok{   data\_point }\OperatorTok{=} \FunctionTok{fetch\_new\_data}\NormalTok{()}

   \CommentTok{\# Collect data for statistics}
   \FunctionTok{push!}\NormalTok{(running\_stats\_storage, data\_point)}

   \CommentTok{\# Periodically we look at summary statistics}
   \ControlFlowTok{if}\NormalTok{ i }\OperatorTok{\%} \FloatTok{20} \OperatorTok{==} \FloatTok{0}
   \FunctionTok{println}\NormalTok{(}\StringTok{"{-}{-}{-}{-}{-}{-}{-}"}\NormalTok{)}
   \FunctionTok{println}\NormalTok{(}\StringTok{"Count: "}\NormalTok{, }\FunctionTok{length}\NormalTok{(running\_stats\_storage))}
   \FunctionTok{println}\NormalTok{(}\StringTok{"Mean: "}\NormalTok{, }\FunctionTok{mean}\NormalTok{(running\_stats\_storage))}
   \FunctionTok{println}\NormalTok{(}\StringTok{"Max: "}\NormalTok{, }\FunctionTok{maximum}\NormalTok{(running\_stats\_storage))}
   \ControlFlowTok{end}
   \ControlFlowTok{end}
\KeywordTok{end}

\BuiltInTok{Random}\NormalTok{.}\FunctionTok{seed!}\NormalTok{(}\FloatTok{0}\NormalTok{)}
\NormalTok{running\_stats\_storage }\OperatorTok{=} \DataTypeTok{Float64}\NormalTok{[] }\CommentTok{\# Some place to hold the running data}
\FunctionTok{print\_running\_stats}\NormalTok{(}\FloatTok{100}\NormalTok{, running\_stats\_storage)}
\end{Highlighting}
\end{Shaded}

\begin{verbatim}
-------
Count: 20
Mean: 43.111155613892045
Max: 90.47275767596541
-------
Count: 40
Mean: 44.657640026434834
Max: 90.47275767596541
-------
Count: 60
Mean: 47.483269933729034
Max: 98.13388392678519
-------
Count: 80
Mean: 46.553989816155244
Max: 98.13388392678519
-------
Count: 100
Mean: 45.13778124972069
Max: 98.13388392678519
\end{verbatim}

Consider the scenario where you have a lot of data to stream and
\texttt{n} gets very big\ldots{} so you don't want to recompute the mean
and max every time. In fact --- you might not even want to store the a
copy of the data in RAM at all!

Here's an approach to do ``online'' statistics:

\begin{Shaded}
\begin{Highlighting}[]
\KeywordTok{mutable struct}\NormalTok{ RunningStats}
\NormalTok{   count}\OperatorTok{::}\DataTypeTok{Int}
\NormalTok{   sum}\OperatorTok{::}\DataTypeTok{Float64}
\NormalTok{   max}\OperatorTok{::}\DataTypeTok{Float64}

   \FunctionTok{RunningStats}\NormalTok{() }\OperatorTok{=} \FunctionTok{new}\NormalTok{(}\FloatTok{0}\NormalTok{, }\FloatTok{0.0}\NormalTok{, }\OperatorTok{{-}}\ConstantTok{Inf}\NormalTok{)}
\KeywordTok{end}

\NormalTok{running\_stats\_storage }\OperatorTok{=} \FunctionTok{RunningStats}\NormalTok{()}
\end{Highlighting}
\end{Shaded}

\begin{verbatim}
RunningStats(0, 0.0, -Inf)
\end{verbatim}

We can now make specific methods for \texttt{push!}, \texttt{length},
\texttt{sum}, \texttt{mean} and \texttt{maximum} for this new type:

\begin{Shaded}
\begin{Highlighting}[]
\ImportTok{import} \BuiltInTok{Base}\NormalTok{: push!, length, sum, maximum}
\ImportTok{import} \BuiltInTok{Statistics}\NormalTok{: mean}

\FunctionTok{length}\NormalTok{(rsd}\OperatorTok{::}\DataTypeTok{RunningStats}\NormalTok{) }\OperatorTok{=}\NormalTok{ rsd.count}
\FunctionTok{sum}\NormalTok{(rsd}\OperatorTok{::}\DataTypeTok{RunningStats}\NormalTok{) }\OperatorTok{=}\NormalTok{ rsd.sum}
\FunctionTok{mean}\NormalTok{(rsd}\OperatorTok{::}\DataTypeTok{RunningStats}\NormalTok{) }\OperatorTok{=}\NormalTok{ rsd.sum }\OperatorTok{/}\NormalTok{ rsd.count}
\FunctionTok{maximum}\NormalTok{(rsd}\OperatorTok{::}\DataTypeTok{RunningStats}\NormalTok{) }\OperatorTok{=}\NormalTok{ rsd.max}

\KeywordTok{function} \FunctionTok{push!}\NormalTok{(rsd}\OperatorTok{::}\DataTypeTok{RunningStats}\NormalTok{, data\_point)}
   \CommentTok{\# Update the count}
\NormalTok{   rsd.count }\OperatorTok{+=} \FloatTok{1}

   \CommentTok{\# Update the sum}
\NormalTok{   rsd.sum }\OperatorTok{+=}\NormalTok{ data\_point}

   \CommentTok{\# Update the maximum}
   \ControlFlowTok{if}\NormalTok{ rsd.max }\OperatorTok{\textless{}}\NormalTok{ data\_point}
\NormalTok{   rsd.max }\OperatorTok{=}\NormalTok{ data\_point}
   \ControlFlowTok{end}
\KeywordTok{end}

\BuiltInTok{Random}\NormalTok{.}\FunctionTok{seed!}\NormalTok{(}\FloatTok{0}\NormalTok{)}
\NormalTok{running\_stats\_storage }\OperatorTok{=} \FunctionTok{RunningStats}\NormalTok{()}
\FunctionTok{print\_running\_stats}\NormalTok{(}\FloatTok{100}\NormalTok{, running\_stats\_storage)}
\end{Highlighting}
\end{Shaded}

\begin{verbatim}
-------
Count: 20
Mean: 43.111155613892045
Max: 90.47275767596541
-------
Count: 40
Mean: 44.65764002643482
Max: 90.47275767596541
-------
Count: 60
Mean: 47.483269933729034
Max: 98.13388392678519
-------
Count: 80
Mean: 46.553989816155244
Max: 98.13388392678519
-------
Count: 100
Mean: 45.137781249720675
Max: 98.13388392678519
\end{verbatim}

\paragraph{Running statistics of other
types}\label{running-statistics-of-other-types}

Going a bit more generic we could have also had,

\begin{Shaded}
\begin{Highlighting}[]
\KeywordTok{mutable struct}\NormalTok{ FlexRunningStats\{T }\OperatorTok{\textless{}:}\DataTypeTok{ Number}\NormalTok{\}}
\NormalTok{   data}\OperatorTok{::}\DataTypeTok{Vector\{T\}}
\NormalTok{   count}\OperatorTok{::}\DataTypeTok{Int}
\NormalTok{   sum}\OperatorTok{::}\DataTypeTok{T}
\NormalTok{   max}\OperatorTok{::}\DataTypeTok{T}

   \FunctionTok{FlexRunningStats}\DataTypeTok{\{T\}}\NormalTok{() }\KeywordTok{where}\NormalTok{ \{T\} }\OperatorTok{=} \FunctionTok{new}\DataTypeTok{\{T\}}\NormalTok{(T[], }\FloatTok{0}\NormalTok{, }\FunctionTok{zero}\NormalTok{(T), }\FunctionTok{typemin}\NormalTok{(T))}
\KeywordTok{end}

\FunctionTok{length}\NormalTok{(rsd}\OperatorTok{::}\DataTypeTok{FlexRunningStats}\NormalTok{) }\OperatorTok{=}\NormalTok{ rsd.count}
\FunctionTok{sum}\NormalTok{(rsd}\OperatorTok{::}\DataTypeTok{FlexRunningStats}\NormalTok{) }\OperatorTok{=}\NormalTok{ rsd.sum}
\FunctionTok{mean}\NormalTok{(rsd}\OperatorTok{::}\DataTypeTok{FlexRunningStats}\NormalTok{) }\OperatorTok{=}\NormalTok{ rsd.sum }\OperatorTok{/}\NormalTok{ rsd.count}
\FunctionTok{maximum}\NormalTok{(rsd}\OperatorTok{::}\DataTypeTok{FlexRunningStats}\NormalTok{) }\OperatorTok{=}\NormalTok{ rsd.max}

\KeywordTok{function} \FunctionTok{push!}\NormalTok{(rsd}\OperatorTok{::}\DataTypeTok{FlexRunningStats}\NormalTok{, data\_point)}
   \CommentTok{\# Insert the new datapoint}
   \FunctionTok{push!}\NormalTok{(rsd.data, data\_point)}

   \CommentTok{\# Update the count}
\NormalTok{   rsd.count }\OperatorTok{+=} \FloatTok{1}

   \CommentTok{\# Update the sum}
\NormalTok{   rsd.sum }\OperatorTok{+=}\NormalTok{ data\_point}

   \CommentTok{\# Update the maximum}
   \ControlFlowTok{if}\NormalTok{ rsd.max }\OperatorTok{\textless{}}\NormalTok{ data\_point}
\NormalTok{   rsd.max }\OperatorTok{=}\NormalTok{ data\_point}
   \ControlFlowTok{end}
\KeywordTok{end}

\FunctionTok{fetch\_new\_data}\NormalTok{() }\OperatorTok{=} \FunctionTok{rand}\NormalTok{(}\FloatTok{0}\OperatorTok{:}\FloatTok{10}\OperatorTok{\^{}}\FloatTok{4}\NormalTok{) }\CommentTok{\# override this to return integers}

\BuiltInTok{Random}\NormalTok{.}\FunctionTok{seed!}\NormalTok{(}\FloatTok{0}\NormalTok{)}
\NormalTok{running\_stats\_storage }\OperatorTok{=} \FunctionTok{FlexRunningStats}\DataTypeTok{\{Int\}}\NormalTok{()}
\FunctionTok{print\_running\_stats}\NormalTok{(}\FloatTok{100}\NormalTok{, running\_stats\_storage)}
\end{Highlighting}
\end{Shaded}

\begin{verbatim}
-------
Count: 20
Mean: 4310.95
Max: 9048
-------
Count: 40
Mean: 4465.675
Max: 9048
-------
Count: 60
Mean: 4748.3
Max: 9814
-------
Count: 80
Mean: 4655.375
Max: 9814
-------
Count: 100
Mean: 4513.76
Max: 9814
\end{verbatim}

\begin{Shaded}
\begin{Highlighting}[]
\PreprocessorTok{@which} \FunctionTok{mean}\NormalTok{(running\_stats\_storage)}
\end{Highlighting}
\end{Shaded}

\begin{verbatim}
mean(rsd::Main.Notebook.FlexRunningStats)
     @ Main.Notebook ~/nextcloud/MATH2504_2026/lecture/quarto/lecture-unit-5.qmd:2194
\end{verbatim}

But there is a problem with the above. What if \texttt{T} was
\texttt{Complex}? Should we have done \texttt{\textless{}:\ Number} or
\texttt{\textless{}:\ Real}?

In generic code its good to consider what \emph{interfaces} we want to
use, so we can figure out which type constraints are appropriate. Julia
will go ahead and assemble the pieces of your program together into a
coherent whole.

\subsubsection{Interfaces}\label{interfaces}

So above we used a few interface methods to implement our
online-statistics accumulator.

\begin{itemize}
\tightlist
\item
  \texttt{push!} - add an element to a collection (mutates the input)
\item
  \texttt{length} - the number of element in a collection
\item
  \texttt{sum} - add up all the element in a collection
\item
  \texttt{maximum} - the maximum number of elements in a collection
\item
  \texttt{mean} - the mean of all the elements in a collection
\end{itemize}

Let's have a look at some common interfaces

\paragraph{Numbers}\label{numbers}

\begin{itemize}
\tightlist
\item
  Common arithmetic operations \texttt{+}, \texttt{-}, \texttt{*},
  \texttt{/}, etc.
\item
  \texttt{promote(x,\ y)} - take two numbers and return two numbers of
  the same type, e.g \texttt{promote\_type(3.14,\ 1)\ =\ (3.14,\ 1.0)}
\item
  \texttt{promote\_rule(T1,\ T2)} -
  e.g.~\texttt{promote\_rule(Float64,\ Int64)\ =\ Float64}.
\end{itemize}

Some subtypes of \texttt{Number} like \texttt{AbstractFloat} have more
interface methods.

\paragraph{Iterables}\label{iterables}

\begin{itemize}
\tightlist
\item
  \texttt{iterate(iter)} - used in \texttt{for} loops
\item
  \texttt{length(iter)} - the number of elements to iterate
\item
  \texttt{eltype(iter)} - the type of the elements to iterate
\end{itemize}

Some things iterate but we might not have known length or element type,
so there you can define some \emph{traits} to describe these facts:

\begin{itemize}
\tightlist
\item
  \texttt{IteratorSize(iter)} - return whether \texttt{length} is known
\item
  \texttt{IteratorEltype(iter)} - return whether \texttt{eltype} is
  known
\end{itemize}

\paragraph{Arrays}\label{arrays-1}

Arrays are iterable and satisfy the the above. They are also indexable

\begin{itemize}
\tightlist
\item
  \texttt{getindex(a,\ i)} - the function behind \texttt{a{[}i{]}}
\item
  \texttt{setindex!(a,\ v,\ i)} - the function behind
  \texttt{a{[}i{]}\ =\ v}
\item
  \texttt{size(a)} - gives a tuple of sizes,
  e.g.~\texttt{size({[}1,2,3{]})\ =\ (3,)}
\end{itemize}

In fact you can create a custom array in Julia with just those methods
above! Things like ranges \texttt{1:10} are subtypes of
\texttt{AbstractArray}.

Some types like \texttt{Vector} are resizable and for these types there
are functions like \texttt{push!}, \texttt{pop!}, \texttt{insert!},
\texttt{deleteat!}, \texttt{append!}, \texttt{empty!}, etc to manipulate
them as arbitrary-sized lists.

Arrays are also defined as arithmetic objects according to linear
algebra, and therefore have \texttt{+}, \texttt{-}, \texttt{*},
\texttt{/}, etc defined. The \texttt{LinearAlgebra} standard library
package has lots of functions to invert or decompose matrices, find
eigenvalues, etc.

\paragraph{Sets}\label{sets}

Julia has a \texttt{Set\{T\}} type and an \texttt{AbstractSet\{T\}}
supertype. Sets are iterable and support things like:

\begin{itemize}
\tightlist
\item
  \texttt{in}
\item
  \texttt{union}
\item
  \texttt{intersect}
\item
  \texttt{setdiff}
\item
  \texttt{symdiff}
\end{itemize}

Here's where interfaces get interesting. Things like \texttt{Array} also
support \texttt{in}, but it is slow (must check every element). With
\texttt{Set} lookup is fast (e.g.~O(1) instead of O(N)) - so the
operations above are fast.

\paragraph{Dictionaries}\label{dictionaries}

Julia has a \texttt{Dict\{K,\ T\}} type and an
\texttt{AbstractDict\{K,\ T\}} supertype. Like arrays, dictionaries are
iterable and indexable. You can imagine they are like a \texttt{Set} but
have associate values (in fact their implementation is related).

\begin{itemize}
\tightlist
\item
  \texttt{keys(dict)} - return the set of keys
\item
  \texttt{haskey(dict,\ key)} - check if a key exists in the dictionary
\item
  \texttt{getindex(dict,\ key)} - get a value associated with a key
\item
  \texttt{setindex!(dict,\ value,\ key)} - insert or update a key
\item
  \texttt{delete!(dict,\ key)} - remove a key
\item
  \texttt{keytype(dict)} - the type of the dictionary keys
\item
  \texttt{valtype(dict)} - the type of the dictionary values
\end{itemize}

You generally iterate a dictionary by specifying if you want to iterate
the keys, values, or both:

\begin{itemize}
\tightlist
\item
  \texttt{pairs(dict)} - an iterable of
  \texttt{key\ =\textgreater{}\ value} pairs (the default)
\item
  \texttt{keys(dict)} - an iterable of the dictionary keys only
\item
  \texttt{values(dict)} - an iterable of the dictionary values only
\end{itemize}

\subsubsection{Using structs for complex, nested data
structures}\label{using-structs-for-complex-nested-data-structures}

You can make interesting data structures by referencing other instances
of yourself, like this tree structure below:

\begin{Shaded}
\begin{Highlighting}[]
\BuiltInTok{Random}\NormalTok{.}\FunctionTok{seed!}\NormalTok{(}\FloatTok{0}\NormalTok{)}

\KeywordTok{struct}\NormalTok{ Node}
\NormalTok{   id}\OperatorTok{::}\DataTypeTok{UInt16}
\NormalTok{   friends}\OperatorTok{::}\DataTypeTok{Vector\{Node\}}
   \FunctionTok{Node}\NormalTok{() }\OperatorTok{=} \FunctionTok{new}\NormalTok{(}\FunctionTok{rand}\NormalTok{(}\DataTypeTok{UInt16}\NormalTok{), [])}
   \FunctionTok{Node}\NormalTok{(friend}\OperatorTok{::}\DataTypeTok{Node}\NormalTok{) }\OperatorTok{=} \FunctionTok{new}\NormalTok{(}\FunctionTok{rand}\NormalTok{(}\DataTypeTok{UInt16}\NormalTok{),[friend])}
\KeywordTok{end}

\StringTok{"""}
\StringTok{Makes \textquotesingle{}n\textasciigrave{} children to node, each with a single friend}
\StringTok{"""}
\KeywordTok{function} \FunctionTok{make\_children}\NormalTok{(node}\OperatorTok{::}\DataTypeTok{Node}\NormalTok{, n}\OperatorTok{::}\DataTypeTok{Int}\NormalTok{, friend}\OperatorTok{::}\DataTypeTok{Node}\NormalTok{)}
   \ControlFlowTok{for}\NormalTok{ \_ }\KeywordTok{in} \FloatTok{1}\OperatorTok{:}\NormalTok{n}
\NormalTok{   new\_node }\OperatorTok{=} \FunctionTok{Node}\NormalTok{(friend)}
   \FunctionTok{push!}\NormalTok{(node.friends, new\_node)}
   \ControlFlowTok{end}
\KeywordTok{end}

\NormalTok{root }\OperatorTok{=} \FunctionTok{Node}\NormalTok{()}
\FunctionTok{make\_children}\NormalTok{(root, }\FloatTok{3}\NormalTok{, root)}
\ControlFlowTok{for}\NormalTok{ node }\KeywordTok{in}\NormalTok{ root.friends}
   \FunctionTok{make\_children}\NormalTok{(node, }\FloatTok{2}\NormalTok{,root)}
\ControlFlowTok{end}
\NormalTok{root}
\end{Highlighting}
\end{Shaded}

\begin{verbatim}
Node(0x67db, Node[Node(0x118c, Node[Node(#= circular reference @-4 =#), Node(0xa95f, Node[Node(#= circular reference @-6 =#)]), Node(0x1dc7, Node[Node(#= circular reference @-6 =#)])]), Node(0xdcb5, Node[Node(#= circular reference @-4 =#), Node(0x1c00, Node[Node(#= circular reference @-6 =#)]), Node(0xb3b6, Node[Node(#= circular reference @-6 =#)])]), Node(0x1602, Node[Node(#= circular reference @-4 =#), Node(0x4a1d, Node[Node(#= circular reference @-6 =#)]), Node(0x074f, Node[Node(#= circular reference @-6 =#)])])])
\end{verbatim}

\subsubsection{More to be covered as part of Unit
6:}\label{more-to-be-covered-as-part-of-unit-6}

\href{https://docs.julialang.org/en/v1/manual/constructors/}{See
constructors in Julia docs}. More on this in Unit 6.

\href{https://docs.julialang.org/en/v1/manual/conversion-and-promotion/}{See
conversion and promotion in Julia docs}

\href{https://docs.julialang.org/en/v1/manual/interfaces/}{See
interfaces in Julia docs}




\end{document}
